% !TEX program = xelatex
\documentclass[14pt,a4paper]{extarticle}

\usepackage{fontspec}
\setmainfont{Times New Roman}
\setsansfont{Arial}
\setmonofont{Consolas}
\usepackage[russian]{babel}
\usepackage{geometry}
\geometry{left=30mm,right=15mm,top=20mm,bottom=20mm}
\usepackage{setspace}
\onehalfspacing
\usepackage{amsmath,amssymb,mathtools}
\usepackage{booktabs,longtable,array,tabularx}
\usepackage{enumitem}
\usepackage{caption}
\usepackage{titlesec}
\usepackage{float}
\usepackage{xcolor}
\usepackage{fancyhdr}
\usepackage{lastpage}
\usepackage{microtype}
\usepackage{graphicx}
\usepackage[hidelinks]{hyperref}
\graphicspath{{figures/article_sources/}}

\setlength{\parindent}{12.5mm}
\setlength{\parskip}{0pt}
\setlength{\emergencystretch}{3em}
\setlist{nosep,leftmargin=12.5mm}
\renewcommand{\arraystretch}{1.18}
\renewcommand{\tablename}{Таблица}
\renewcommand{\figurename}{Рисунок}
\addto\captionsrussian{%
  \renewcommand{\refname}{СПИСОК ИСПОЛЬЗОВАННЫХ ИСТОЧНИКОВ}%
}
\titleformat{\section}[hang]
  {\normalfont\Large\bfseries\raggedright\hyphenpenalty=10000}
  {\thesection}{1em}{}
\titleformat{\subsection}[hang]
  {\normalfont\large\bfseries}{\thesubsection}{1em}{}
\titleformat{\subsubsection}[hang]
  {\normalfont\normalsize\bfseries}{\thesubsubsection}{1em}{}
\newcommand{\sectionbreak}{\clearpage}
\DeclareCaptionLabelSeparator{gostdash}{~---~}
\captionsetup{font=small,labelsep=gostdash}
\captionsetup[table]{position=top,skip=6pt,justification=raggedleft,
singlelinecheck=false}
\captionsetup[figure]{position=bottom,skip=6pt,justification=centering,
singlelinecheck=false}

\pagestyle{fancy}
\fancyhf{}
\fancyfoot[C]{\thepage}
\renewcommand{\headrulewidth}{0pt}

\newcommand{\source}[1]{\par\noindent\textit{Источник:} #1\par}
\newcommand{\phys}{\par\noindent\textbf{Физическая сущность.} }
\newcommand{\method}{\par\noindent\textbf{Методика.} }
\newcommand{\recommend}{\par\noindent\textbf{Рекомендация по выбору.} }
\newcommand{\subst}{\par\noindent\textbf{Подстановка.} }
\newcommand{\result}{\par\noindent\textbf{Результат.} }
\newcommand{\criterion}{\par\noindent\textbf{Критерий решения.} }
\newcommand{\stepgoal}[1]{\par\noindent\textbf{Цель шага.} #1\par}
\newcommand{\dd}{\mathop{}\!\mathrm{d}}
\newenvironment{whereblock}
 {\par\noindent где\par
  \begin{list}{}{\setlength{\leftmargin}{12.5mm}
                  \setlength{\itemindent}{-12.5mm}
                  \setlength{\itemsep}{0pt}
                  \setlength{\parsep}{0pt}
                  \setlength{\topsep}{0pt}}}
 {\end{list}}
\newcommand{\whereitem}[2]{\item[] #1 --- #2}

\title{Методика расчёта высокочастотных силовых дросселей\\
с дискретным или распределённым зазором\\
для частот до 1 МГц}
\author{Универсальная методика\\без привязки к конкретному изделию}
\date{2026}

\begin{document}
\maketitle
\thispagestyle{empty}
\newpage

\section*{РЕФЕРАТ}
\addcontentsline{toc}{section}{РЕФЕРАТ}

Методика объёмом \pageref{LastPage}~с. устанавливает
воспроизводимую последовательность проектирования высокочастотного силового
дросселя: от восстановления электрического режима и предварительного выбора
магнитопровода до разработки обмотки, расчёта потерь, частотной характеристики,
электромагнитных помех, численной и экспериментальной верификации. Основной
критерий энергетической эффективности --- минимум потерь дросселя во всём
заданном профиле режимов; электромагнитная совместимость проверяется по
дифференциальным и синфазным путям распространения помех. Для конструкций с
жёстким ограничением габарита отдельно рассмотрены гибридные магнитопроводы с
ферритовым основным телом, порошковыми вставками центрального стержня и
несколькими дискретными воздушными зазорами.

\noindent\textbf{Ключевые слова:} высокочастотный дроссель, магнитопровод,
воздушный зазор, гибридный магнитопровод, Sendust, множественный зазор,
литцендрат, потери, коэффициент полезного действия, собственный резонанс,
амплитудно-частотная характеристика, синфазные помехи, метод конечных элементов.

\newpage
\tableofcontents
\newpage

\section*{ТЕРМИНЫ, ОПРЕДЕЛЕНИЯ И СОКРАЩЕНИЯ}
\addcontentsline{toc}{section}{ТЕРМИНЫ, ОПРЕДЕЛЕНИЯ И СОКРАЩЕНИЯ}

\textbf{Высокочастотный силовой дроссель} --- индуктивный компонент силового
преобразователя, который накапливает энергию магнитного поля и ограничивает
изменение тока на частоте коммутации и её гармониках.

\textbf{Дискретный воздушный зазор} --- один или несколько локальных
немагнитных участков магнитной цепи. \textbf{Распределённый зазор} ---
эквивалентный немагнитный промежуток, распределённый в объёме порошкового
магнитного материала.

\textbf{Гибридный магнитопровод} --- составная магнитная цепь, участки которой
выполнены из материалов с различными магнитной проницаемостью и удельными
потерями. В рассматриваемой конструкции основное тело выполнено из феррита, а
одна или две вставки центрального стержня --- из феррита либо Sendust
(порошковый сплав Fe--Si--Al); между вставками расположены дискретные воздушные
зазоры.

\textbf{Секущая индуктивность} $L_{\text{сек}}$ --- отношение потокосцепления
к току в заданной рабочей точке. \textbf{Дифференциальная индуктивность}
$L_{\text{диф}}$ --- производная потокосцепления по току; именно она определяет
малое приращение и пульсацию тока около рабочей точки.

\textbf{Собственный резонанс} --- частота, на которой индуктивные и ёмкостные
реактивности дросселя взаимно компенсируются. Выше первого собственного
резонанса дроссель обычно перестаёт быть индуктивным элементом.

\textbf{Дифференциальная помеха} (differential-mode interference, DM) ---
помеховый ток, текущий по силовой линии и возвращающийся по другой силовой
линии. \textbf{Синфазная помеха} (common-mode interference, CM) --- помеховый
ток, текущий в одном направлении по силовым проводникам и возвращающийся через
паразитные ёмкости, корпус, защитное заземление или окружающие конструкции.

\textbf{Амплитудно-частотная характеристика} (АЧХ) --- зависимость модуля
комплексного сопротивления или коэффициента передачи от частоты. Для
однозначной интерпретации АЧХ дросселя всегда рассматривают совместно с
фазочастотной характеристикой.

\textbf{Метод конечных элементов} (finite element method, FEM) --- численный
метод решения полевых задач после разбиения расчётной области на конечные
элементы.

\textbf{Эквивалентная сеть питания} (line impedance stabilization network,
LISN) --- измерительная сеть с нормированным импедансом, применяемая при
измерении кондуктивных электромагнитных помех.

\newpage

\section{НАЗНАЧЕНИЕ И ГРАНИЦЫ ПРИМЕНИМОСТИ}
\label{sec:scope}

Настоящая методика устанавливает последовательность проектирования силовых
дросселей с проволочной, литцендратной, фольговой или планарной обмоткой.
Рассматриваются ферритовые магнитопроводы с одним или несколькими дискретными
зазорами, а также порошковые магнитопроводы с распределённым зазором.
Предельная частота коммутации составляет $1\,\text{МГц}$.

Указанная частота является границей алгоритма, а не гарантией применимости
конкретного материала. Материал допускается использовать только в диапазоне
частоты, температуры, размаха магнитной индукции и постоянного
подмагничивания, для которого имеются паспортные или измеренные данные.
Экстраполяция кривых потерь и магнитной проницаемости за подтверждённый
диапазон не допускается.

Методика предназначена для:

\begin{enumerate}
\item предварительного выбора магнитопровода, материала, числа витков и
воздушного зазора;
\item расчёта проводников и предварительной оценки потерь;
\item подготовки исходных данных для метода конечных элементов
(finite element method, FEM);
\item совместной электромагнитно-тепловой оптимизации;
\item проверки электромагнитных помех и паразитных параметров;
\item формирования программы испытаний опытного образца.
\end{enumerate}

Аналитический расчёт не заменяет измерение $L(I,T)$, комплексного импеданса
$Z(f)$, температуры и электромагнитных помех. Для дросселя с дискретным
зазором обязательна проверка поля у кромок зазора.

Базовые соотношения магнитной цепи и последовательность предварительного
выбора магнитопровода приняты по учебным пособиям
\cite{obrusnik,legostaev}; высокочастотные уточнения приведены в профильных
источниках соответствующих разделов.

\section{ИСХОДНЫЕ ДАННЫЕ И РАСЧЁТНАЯ ПОСЛЕДОВАТЕЛЬНОСТЬ}
\label{sec:inputs}

\subsection{Обязательные исходные данные}

\begin{table}[H]
\begin{singlespace}
\caption{Обязательные исходные данные}
\label{tab:required-inputs}
\centering
\scriptsize
\begin{tabularx}{\textwidth}{>{\raggedright\arraybackslash}p{38mm}
>{\raggedright\arraybackslash}X
>{\raggedright\arraybackslash}p{48mm}}
\toprule
Группа & Требуемые величины & Предпочтительный источник \\
\midrule
Силовая схема &
$u_L(t)$, $i_L(t)$, диапазон входного напряжения, нагрузка, номинальная
частота либо допустимый диапазон частоты, признак возможности её выбора,
переходные и аварийные режимы &
Схемотехническое моделирование и осциллограммы \\
Профиль эксплуатации &
Длительность или вероятность каждого режима, температура окружающей среды,
режим охлаждения &
Техническое задание и профиль нагрузки изделия \\
Индуктивность &
$L_{\text{ном}}$, допуск, минимальная дифференциальная индуктивность
$L_{\text{диф,min}}(I,T)$ &
Техническое задание и динамическая модель преобразователя \\
Ограничения конструкции &
предельные габариты, масса, способ крепления, доступная площадь печатной
платы, высота, требования к изоляции и допустимые технологии намотки &
Техническое задание и технологические нормы \\
Каталог кандидатов &
доступные марки материала, формы и типоразмеры магнитопроводов, каркасы,
провода и стандартные прокладки зазора &
Формируется в ходе шагов выбора, не задаётся пользователем заранее \\
Ограничения обмотки &
Допустимые материалы и классы нагревостойкости, минимальные изоляционные
расстояния, доступные технологии намотки или печатной платы &
Техническое задание и технологические нормы изготовителя \\
Охлаждение &
Температура окружающей среды, конвекция, контактные поверхности,
тепловые сопротивления &
Тепловая модель и условия установки \\
Помехи &
Нормируемая полоса, схема эквивалентной сети питания, допустимые
дифференциальные и синфазные помехи &
Техническое задание и применимый стандарт изделия \\
Критерий выбора &
приоритет коэффициента полезного действия, помех, габарита, массы и стоимости;
числовые ограничения и требуемый запас &
Техническое задание; при отсутствии --- правила подраздела об оптимизации \\
\bottomrule
\end{tabularx}
\end{singlespace}
\end{table}

\subsection{Расчётная последовательность}

\method Разделы методики применяют последовательно. Возврат к предыдущему
шагу является нормальной проектной итерацией. Перед первым расчётом разработчик
заполняет таблицу~\ref{tab:required-inputs}; неизвестное значение помечает как
«требует модели» или «требует измерения», но не заменяет произвольной
константой.

\begin{enumerate}
\item по разделу~\ref{sec:electrical-mode} определяют формы $u_L(t)$ и
$i_L(t)$, максимальные ток, потокосцепление и энергию;
\item по подразделу~\ref{subsec:material-quantities} рассчитывают величины, которые предстоит
сопоставлять с графиками материала: $f_{\text{с}}$, $T_{\max}$,
$B_{\text{DC}}$, $\widehat B_{\text{AC}}$ и $B_{\max}$;
\item по подразделу~\ref{subsec:material-choice} оставляют несколько
допустимых марок материала и назначают предварительное
$B_{\text{доп,пред}}$;
\item по подразделу~\ref{subsec:size-conductor-preselection} выбирают
предварительные $J_{\text{доп}}$, коэффициент заполнения и сечение меди;
\item по подразделу~\ref{subsec:core-size-choice} выбирают из каталога
несколько конкретных пар
«типоразмер~-- материал», используя $A_e$, $A_{\min}$, $V_e$, окно и
габариты; коэффициент $A_{L0}$ на этом шаге не является критерием выбора;
\item по подразделу~\ref{subsec:turn-gap-loop} для каждого кандидата
перебирают допустимые целые
числа витков, а для каждого числа витков рассчитывают $A_L$, зазор,
индукцию, заполнение окна и предварительные потери;
\item если габарит фиксирован и обычный ферритовый кандидат не обеспечивает
требуемую характеристику $L(I,T)$, по разделу~\ref{sec:hybrid-core}
рассчитывают варианты с ферритовыми или Sendust-вставками и несколькими
зазорами;
\item по разделу~\ref{sec:winding} до начала FEM выбирают диаметр и число
жил, строят послойную укладку и получают фактические размеры обмотки;
\item по разделу~\ref{sec:losses} уточняют потери в сердечнике и обмотке;
\item по разделу~\ref{sec:parasitics} рассчитывают частотную характеристику,
дифференциальные и синфазные пути помех;
\item по разделу~\ref{sec:fem} проверяют уже определённую геометрию
магнитопровода, зазора, обмотки и окружающих деталей;
\item по разделу~\ref{sec:thermal} замыкают температурную итерацию;
\item по разделу~\ref{sec:optimization} сопоставляют варианты и проверяют
все ограничения;
\item выпускают чертёж, изготавливают образец и замыкают расчёт
измерениями.
\end{enumerate}

Результатом предварительного проектирования должен быть не абстрактный
«сердечник», а комплект:
точное обозначение и код заказа магнитопровода, марка материала, число
витков, схема и допуск зазора, обозначение проводника, схема укладки,
каркас или печатная плата, изоляция и ожидаемые электрические и тепловые
параметры.


\section{ОПРЕДЕЛЕНИЕ ЭЛЕКТРИЧЕСКОГО РЕЖИМА}
\label{sec:electrical-mode}

\subsection{Период, вольт-секундная площадь и ток}

\textbf{Цель шага} --- получить формы напряжения и тока, необходимые для
расчёта числа витков, запаса по насыщению и потерь.

Период коммутации определяют по формуле
\begin{equation}
 T_{\text{с}}=\frac{1}{f_{\text{с}}},
 \label{eq:period}
\end{equation}
где $T_{\text{с}}$ --- период коммутации, с;\par\noindent
$f_{\text{с}}$ --- частота коммутации, Гц.

В расчёте магнитного режима используют напряжение на индуктивной
составляющей дросселя
\begin{equation}
 u_L(t)=u_{\text{выв}}(t)-u_R(t),
 \qquad
 u_R(t)\approx R_{\text{DC}}(\vartheta)i(t),
 \label{eq:inductive-voltage}
\end{equation}
\begin{whereblock}
\whereitem{$u_L(t)$}{мгновенное напряжение на индуктивной составляющей,
В;}
\whereitem{$u_{\text{выв}}(t)$}{мгновенное напряжение между выводами
реального дросселя, В;}
\whereitem{$u_R(t)$}{составляющая напряжения, обусловленная потерями в
обмотке, В;}
\whereitem{$R_{\text{DC}}(\vartheta)$}{сопротивление обмотки постоянному
току при температуре $\vartheta$, Ом;}
\whereitem{$\vartheta$}{температура обмотки, $^{\circ}\text{С}$;}
\whereitem{$i(t)$}{мгновенный ток обмотки, А.}
\end{whereblock}

Приближение во второй части формулы~\eqref{eq:inductive-voltage} применяют
ниже первого собственного резонанса, если добавочная составляющая
сопротивления обмотки мала. В ином случае $u_R(t)$ получают из частотной
модели или измерения. Смешивать напряжение между выводами $u_{\text{выв}}$
и напряжение $u_L$ в вольт-секундном расчёте нельзя.

В установившемся периодическом режиме среднее значение $u_L(t)$ за период
равно нулю:
\begin{equation}
 \int_{0}^{T_{\text{с}}}u_L(t)\,dt=0,
 \label{eq:volt-second-balance}
\end{equation}
где $t$ --- время, с.

Положительную вольт-секундную площадь определяют по формуле
\begin{equation}
 \Lambda_+=\int_{t_1}^{t_2}u_L(t)\,dt,
 \label{eq:volt-second-area}
\end{equation}
где $\Lambda_+$ --- положительная вольт-секундная площадь, В$\cdot$с;\par\noindent
$t_1$, $t_2$ --- границы интервала положительного напряжения, с.

Для линейной индуктивности размах пульсации тока равен
\begin{equation}
 \Delta I_{\text{пп}}=\frac{\Lambda_+}{L},
 \label{eq:ripple-linear}
\end{equation}
где $\Delta I_{\text{пп}}$ --- размах пульсации тока от минимума до
максимума, А;\par\noindent
$L$ --- индуктивность дросселя, Гн.

При выраженной нелинейности сердечника ток находят интегрированием
дифференциального уравнения
\begin{equation}
 \frac{di}{dt}=\frac{u_L(t)}
 {L_{\text{диф}}\!\left(i,T\right)},
 \label{eq:current-nonlinear}
\end{equation}
где $i$ --- мгновенный ток, А;\par\noindent
$L_{\text{диф}}=d\Psi/di$ --- дифференциальная индуктивность, Гн;\par\noindent
$\Psi$ --- потокосцепление, Вб$\cdot$виток;\par\noindent
$T$ --- температура, К.

Для треугольной пульсации с постоянной составляющей применяют
\begin{align}
 I_{\max}&=I_{\text{ср}}+\frac{\Delta I_{\text{пп}}}{2},
 \label{eq:ipeak}\\
 I_{\text{действ}}&=
 \sqrt{I_{\text{ср}}^2+\frac{\Delta I_{\text{пп}}^2}{12}},
 \label{eq:irms-tri}
\end{align}
где $I_{\max}$ --- максимальный ток, А;\par\noindent
$I_{\text{ср}}$ --- средний ток, А;\par\noindent
$I_{\text{действ}}$ --- действующее значение тока, А.

\recommend Значения $I_{\text{ср}}$, $I_{\max}$ и $I_{\text{действ}}$
определяют для всех установившихся и переходных режимов. Для
широтно-импульсной модуляции с переменными частотой или заполнением
расчёт выполняют во всём диапазоне управления, а не только в номинальной
точке.

\subsection{Потокосцепление и запасаемая энергия}
\label{subsec:flux-linkage-energy}

Напряжение определяет производную потокосцепления:
\begin{equation}
 u_L(t)=\frac{d\Psi(t)}{dt},
 \label{eq:flux-linkage-voltage}
\end{equation}
где $u_L(t)$ --- напряжение на индуктивной составляющей дросселя, В;
\par\noindent
$\Psi(t)$ --- мгновенное потокосцепление обмотки, Вб$\cdot$виток.

При обработке измеренного напряжения $u_L(t)$ восстанавливают по
формуле~\eqref{eq:inductive-voltage}. Изменение потокосцепления на выбранном
интервале определяют по формуле
\begin{equation}
 \Delta\Psi(t_1,t_2)=\int_{t_1}^{t_2}u_L(t)\,dt,
 \label{eq:flux-linkage}
\end{equation}
где $\Delta\Psi(t_1,t_2)$ --- изменение потокосцепления от момента $t_1$
до момента $t_2$, Вб$\cdot$виток;
\par\noindent
$t_1$, $t_2$ --- границы интервала интегрирования, с.

\phys Формула~\eqref{eq:flux-linkage} определяет только изменение
потокосцепления. По одной вольт-секундной площади нельзя найти
$\left|\Psi\right|_{\max}$, поскольку неизвестна постоянная составляющая,
создаваемая средним током дросселя.

Для каждого расчётного режима $r$ временную зависимость потокосцепления
восстанавливают относительно опорного момента $t_0$:
\begin{equation}
 \Psi_r(t)=\Psi_r(t_0)+
 \int_{t_0}^{t}u_{L,r}(\tau)\,d\tau,
 \label{eq:flux-linkage-trajectory}
\end{equation}
\begin{whereblock}
\whereitem{$\Psi_r(t)$}{потокосцепление в $r$-м расчётном режиме,
Вб$\cdot$виток;}
\whereitem{$\Psi_r(t_0)$}{опорное потокосцепление в момент $t_0$,
Вб$\cdot$виток;}
\whereitem{$u_{L,r}(\tau)$}{напряжение на индуктивной составляющей
дросселя в $r$-м режиме, В;}
\whereitem{$\tau$}{переменная интегрирования по времени, с.}
\end{whereblock}

Опорный момент удобно выбирать в точке минимального тока периода. Если
известна нелинейная характеристика $\Psi(I,T)$, опорное значение берут
непосредственно из неё. Если известна только дифференциальная индуктивность,
то при пренебрежении остаточным потокосцеплением применяют
\begin{equation}
 \Psi_r(t_0)=
 \int_{0}^{I_r(t_0)}
 L_{\text{диф}}(I,T_r)\,dI,
 \label{eq:flux-linkage-reference}
\end{equation}
\begin{whereblock}
\whereitem{$I_r(t_0)$}{ток дросселя в опорный момент $t_0$ для $r$-го
режима, А;}
\whereitem{$L_{\text{диф}}(I,T_r)$}{дифференциальная индуктивность при токе
$I$ и температуре $T_r$, Гн;}
\whereitem{$T_r$}{температура магнитопровода в $r$-м режиме, К.}
\end{whereblock}

Для линейной модели с постоянной индуктивностью расчёт упрощается:
\begin{equation}
 \Psi_r(t)=L_r i_{L,r}(t),
 \qquad
 \left|\Psi\right|_{\max,r}=
 L_r\max_{0\leq t\leq T_{\text{с},r}}\left|i_{L,r}(t)\right|,
 \label{eq:flux-linkage-linear-maximum}
\end{equation}
\begin{whereblock}
\whereitem{$L_r$}{индуктивность, принятая постоянной в $r$-м режиме, Гн;}
\whereitem{$i_{L,r}(t)$}{мгновенный ток дросселя в $r$-м режиме, А;}
\whereitem{$T_{\text{с},r}$}{период изменения электрического режима в
$r$-м режиме, с;}
\whereitem{$\left|\Psi\right|_{\max,r}$}{максимальный модуль
потокосцепления в $r$-м режиме, Вб$\cdot$виток.}
\end{whereblock}

Для однополярного тока с треугольной пульсацией предварительная оценка до
выбора конструкции имеет вид
\begin{equation}
 \left|\Psi\right|_{\max,r}^{(0)}=
 L_{\text{расч},r}
 \left(
 I_{\text{ср},r}+\frac{\Delta I_{\text{пп},r}}{2}
 \right)
 =L_{\text{расч},r}I_{\max,r},
 \label{eq:flux-linkage-preliminary}
\end{equation}
\begin{whereblock}
\whereitem{$\left|\Psi\right|_{\max,r}^{(0)}$}{предварительная оценка
максимального модуля потокосцепления в $r$-м режиме, Вб$\cdot$виток;}
\whereitem{$L_{\text{расч},r}$}{индуктивность, принятая при расчёте тока в
$r$-м режиме, Гн;}
\whereitem{$I_{\text{ср},r}$}{средний ток дросселя в $r$-м режиме, А;}
\whereitem{$\Delta I_{\text{пп},r}$}{размах пульсации тока в $r$-м режиме,
А;}
\whereitem{$I_{\max,r}$}{максимальный ток дросселя в $r$-м режиме, А.}
\end{whereblock}

После выбора магнитопровода, числа витков и зазора предварительную оценку
заменяют результатом нелинейной модели или измеренной характеристикой
$\Psi(I,T)$. Окончательное значение для расчёта числа витков и проверки
насыщения определяют по формуле
\begin{equation}
 \left|\Psi\right|_{\max}=
 \max_{r\in\mathcal R}
 \left[
 \max_{0\leq t\leq T_{\text{с},r}}\left|\Psi_r(t)\right|
 \right],
 \label{eq:flux-linkage-global-maximum}
\end{equation}
\begin{whereblock}
\whereitem{$\left|\Psi\right|_{\max}$}{наибольший модуль потокосцепления
среди всех рассчитанных моментов времени и режимов, Вб$\cdot$виток;}
\whereitem{$\mathcal R$}{множество установившихся, переходных, аварийных и
предельных по допускам расчётных режимов;}
\whereitem{$T_{\text{с},r}$}{период изменения электрического режима в
$r$-м режиме, с;}
\whereitem{$r$}{номер расчётного режима.}
\end{whereblock}

\recommend При интегрировании измеренной периодической формы напряжения
сначала проверяют условие~\eqref{eq:volt-second-balance}. Небольшую постоянную
составляющую, вызванную смещением измерительного канала, удаляют до
интегрирования. Физическую вольт-секундную несимметрию переходного или
аварийного режима не удаляют.

\criterion До выбора конструкции используют только предварительную оценку
по формуле~\eqref{eq:flux-linkage-preliminary}. Если не заданы ни расчётная
индуктивность, ни опорное значение $\Psi(I,T)$, по напряжению можно определить
только $\Delta\Psi$; величину $\left|\Psi\right|_{\max}$ в этом случае
помечают как неопределённую. Подставлять вместо неё положительную
вольт-секундную площадь $\Lambda_+$ при однополярном токе недопустимо.

Для линейной индуктивности энергия магнитного поля равна
\begin{equation}
 W=\frac{L I^2}{2},
 \label{eq:energy-linear}
\end{equation}
где $W$ --- энергия магнитного поля, Дж;\par\noindent
$I$ --- ток, А.

Для однозначной квазистатической характеристики без учёта гистерезиса
энергию вычисляют по формуле
\begin{equation}
 W(I,T)=
 \int_{0}^{\Psi(I,T)}i(\widetilde\Psi,T)\,d\widetilde\Psi,
 \label{eq:energy-nonlinear}
\end{equation}
где $\widetilde\Psi$ --- переменная интегрирования по потокосцеплению,
Вб$\cdot$виток.

Величина $\int_0^I\Psi(i,T)\,di$ является магнитной коэнергией; для
нелинейной характеристики она в общем случае не равна энергии. При наличии
гистерезиса траектория $\Psi(i)$ зависит от предыстории, поэтому потери за
цикл рассчитывают отдельно по площади петли либо по подтверждённой модели
потерь.

\recommend Максимальную вольт-секундную площадь и максимальную энергию
определяют независимо. Эти максимумы могут соответствовать разным
сочетаниям входного напряжения, нагрузки, температуры и допуска
индуктивности.

\section{ВЫБОР МАРКИ МАТЕРИАЛА, ТИПОРАЗМЕРА, ЧИСЛА ВИТКОВ И ЗАЗОРА}
\label{sec:magnetic-design}

\subsection{Расчёт величин для сопоставления с паспортом материала}
\label{subsec:material-quantities}

\textbf{Цель шага} --- преобразовать электрические исходные данные в
величины, нанесённые на графики изготовителя. На этом шаге конкретный
магнитопровод ещё не выбран, поэтому расчёт индукции повторяют для
каждого кандидата и каждого проверяемого числа витков.

Для кандидата с минимальным сечением $A_{\min}$ предварительную оценку
максимальной магнитной индукции, средней по этому сечению, определяют по
формуле
\begin{equation}
 B_{\max}(N)=
 \frac{\left|\Psi\right|_{\max}}{N A_{\min}},
 \label{eq:bmax-candidate}
\end{equation}
где $B_{\max}(N)$ --- предварительная максимальная индукция, средняя по
минимальному сечению, при числе витков $N$, Тл;\par\noindent
$\left|\Psi\right|_{\max}$ --- максимальный модуль потокосцепления,
на первой итерации оценённый по формуле~\eqref{eq:flux-linkage-preliminary},
а после выбора конструкции определённый по
формуле~\eqref{eq:flux-linkage-global-maximum},
Вб$\cdot$виток;\par\noindent
$N$ --- проверяемое целое число витков;\par\noindent
$A_{\min}$ --- минимальная площадь поперечного сечения
магнитопровода, м$^2$.

Формула~\eqref{eq:bmax-candidate} предполагает, что полезный магнитный
поток одинаков во всех витках, а поток рассеяния и неоднородность поля в
сечении несущественны. Локальный максимум магнитной индукции, особенно у
кромок зазора, по этой формуле не определяют: его получают в сходящейся
полевой модели по разделу~\ref{sec:fem}.

Размах переменной магнитной индукции определяют по формуле
\begin{equation}
 \Delta B_{\text{пп}}(N)=
 \frac{\Delta\Psi_{\text{пп}}}{N A_{\min}},
 \label{eq:delta-b}
\end{equation}
где $\Delta B_{\text{пп}}$ --- размах магнитной индукции от минимума до
максимума, Тл;\par\noindent
$\Delta\Psi_{\text{пп}}$ --- размах потокосцепления за период,
Вб$\cdot$виток.

Для симметричной переменной составляющей амплитуда, используемая на
графиках потерь TDK, равна
\begin{equation}
 \widehat B_{\text{AC}}=
 \frac{\Delta B_{\text{пп}}}{2},
 \label{eq:bhat-ac}
\end{equation}
где $\widehat B_{\text{AC}}$ --- амплитуда переменной составляющей
магнитной индукции, Тл.

Постоянную составляющую определяют по формуле
\begin{equation}
 B_{\text{DC}}=
 \frac{B_{\max}+B_{\min}}{2},
 \label{eq:bdc}
\end{equation}
где $B_{\text{DC}}$ --- постоянная составляющая магнитной индукции, Тл;\par\noindent
$B_{\min}$ --- минимальная индукция за период, Тл.

Для несимметричной формы сначала восстанавливают $B(t)$ интегрированием
напряжения по формуле~\eqref{eq:flux-linkage}, затем непосредственно
определяют $B_{\max}$, $B_{\min}$ и $\Delta B_{\text{пп}}$. Обозначения
$B$, $\widehat B$ и $\Delta B$ у разных изготовителей могут различаться;
перед чтением графика проверяют раздел \textit{Symbols and terms}
конкретного паспорта.

\begin{table}[H]
\caption{Сопоставление расчётных величин с паспортом материала}
\label{tab:material-datasheet-map}
\centering
\scriptsize
\begin{tabularx}{\textwidth}{>{\raggedright\arraybackslash}p{29mm}
>{\raggedright\arraybackslash}p{38mm}
>{\raggedright\arraybackslash}p{47mm} X}
\toprule
Что рассчитано & Что искать в паспорте & Условие предварительного принятия &
Ограничение применения \\
\midrule
$f_{\text{с}}$ и существенные гармоники &
\textit{Optimum frequency range}, график комплексной проницаемости &
$f_{\text{с}}$ входит в подтверждённый диапазон &
Диапазон является фильтром кандидатов, но не гарантирует допустимые потери \\
$T_{\max}$ &
$B_S(T)$, $p_v(T)$, температура Кюри $T_C$ &
имеются данные не ниже расчётной температуры &
$T_C$ не является допустимой рабочей температурой \\
$B_{\max}$ &
\textit{Flux density} $B_S(T)$, динамическая петля $B(H,T)$ &
до выбора конструкции принять
$B_{\text{доп,пред}}=0{,}60B_S(T_{\max})$; окончательный критерий
проверить после расчёта конструкции &
$B_S$ берут при той температуре и условиях измерения, которые указаны в
паспорте \\
$\widehat B_{\text{AC}}$, $f_{\text{с}}$, $T$ &
$p_v(\widehat B)$, $p_v(f)$, $p_v(T)$ &
$P_{\text{Fe}}=p_vV_e$ укладывается в тепловой бюджет &
Для несинусоидальной формы применяют iGSE в границах её применимости либо
измерение, а не одну точку графика \\
$B_{\text{DC}}$ или локальное $H_{\text{DC}}$ &
\textit{DC magnetic bias}, амплитудная проницаемость &
снижение дифференциальной проницаемости учтено в $L_{\text{диф,min}}$ &
Кривая малого сигнала на тороиде не заменяет нелинейную модель
магнитопровода с зазором \\
Требуемые форма и размер &
\textit{Core shapes}, каталог типоразмеров &
материал выпускается в выбранной форме и доступен для заказа &
Название материала без кода конкретного сердечника недостаточно \\
\bottomrule
\end{tabularx}
\end{table}

\subsection{Предварительный выбор марки материала}
\label{subsec:material-choice}

\stepgoal{Задать допустимую индукцию не произвольным процентом, а по
паспортной индукции насыщения при наибольшей расчётной температуре и с
явным запасом на неопределённость.}

Допустимую магнитную индукцию определяют по формуле
\begin{equation}
 B_{\text{доп}}(T)=k_B B_S(T)-u_B,
 \label{eq:b-allow}
\end{equation}
\begin{whereblock}
\whereitem{$B_{\text{доп}}$}{допустимая магнитная индукция, Тл;}
\whereitem{$k_B$}{коэффициент использования паспортной индукции насыщения;}
\whereitem{$B_S(T)$}{паспортная индукция насыщения при температуре $T$, Тл;}
\whereitem{$u_B$}{абсолютный запас на неучтённые допуски и численную
неопределённость, Тл.}
\end{whereblock}

\subsubsection{Предварительное применение до выбора конструкции}

\criterion Если ещё не выбраны типоразмер магнитопровода, число витков и
зазор, величины $B_{\max,\text{ном}}$, $B_{\max,\text{нх}}$ и
$u_{B,\text{числ}}$ не рассчитывают. Они не являются исходными данными.
На этом этапе неизвестны даже $A_{\min}$ и $N$, необходимые для
определения магнитной индукции.

Для предварительного отбора материала принимают
\begin{equation}
 B_{\text{доп,пред}}(T_{\max})=
 0{,}60 B_S(T_{\max}),
 \label{eq:b-allow-preliminary}
\end{equation}
\begin{whereblock}
\whereitem{$B_{\text{доп,пред}}$}{предварительно допустимая магнитная
индукция до выбора конструкции, Тл;}
\whereitem{$T_{\max}$}{наибольшая расчётная температура
магнитопровода, К.}
\end{whereblock}

Коэффициенты $k_B=0{,}60$, $0{,}65$ и $0{,}70\ldots0{,}75$ являются
внутренними правилами запаса настоящей методики, а не рекомендациями
изготовителя материала. Во всех случаях приёмочный критерий по минимальной
дифференциальной индуктивности имеет приоритет перед процентным запасом по
$B_S$.

Формула~\eqref{eq:b-allow-preliminary} соответствует подстановке
$k_B=0{,}60$ и $u_B=0$ в формулу~\eqref{eq:b-allow}. Весь общий запас
на неизвестную конструкцию включён в консервативный коэффициент $k_B$.
До выбора типоразмера материал проверяют только по
$B_S(T_{\max})$, рабочему диапазону частоты, паспортным потерям и
доступности требуемой формы.

\begin{table}[H]
\caption{Появление величин для проверки магнитной индукции}
\label{tab:b-uncertainty-stages}
\centering
\small
\begin{tabularx}{\textwidth}{
>{\raggedright\arraybackslash}p{30mm}
>{\raggedright\arraybackslash}p{52mm}
>{\raggedright\arraybackslash}X}
\toprule
Этап & Уже выбранные или полученные данные & Применяемое правило \\
\midrule
Отбор материала &
Электрический режим, $T_{\max}$, паспортное $B_S(T)$ &
$k_B=0{,}60$, $u_B=0$;
$B_{\text{доп,пред}}=0{,}60B_S(T_{\max})$ \\
Расчёт кандидата &
Материал, типоразмер, $A_{\min}$, целое $N$, номинальные $L$ и $g$ &
рассчитать $B_{\max,\text{ном}}$ \\
Анализ допусков &
Допуски тока, площади, индуктивности, зазора, температуры и свойства
материала &
рассчитать $B_{\max,\text{нх}}$ перебором предельных сочетаний \\
Окончательная проверка &
Сходящаяся полевая модель и, при наличии, результаты измерения образца &
определить $u_{B,\text{числ}}$ и проверить окончательный критерий \\
\bottomrule
\end{tabularx}
\end{table}

\subsubsection{Расчёт после выбора типоразмера и числа витков}

После выбора кандидатного магнитопровода и числа витков номинальную
максимальную магнитную индукцию определяют по формуле
\begin{equation}
 B_{\max,\text{ном}}=
 \frac{
 \left|\Psi_{\max,\text{ном}}\right|
 }{
 N A_{\min,\text{ном}}
 },
 \label{eq:bmax-nominal}
\end{equation}
\begin{whereblock}
\whereitem{$B_{\max,\text{ном}}$}{аналитическая оценка максимальной
магнитной индукции, средней по минимальному сечению, при номинальных
параметрах выбранной конструкции, Тл;}
\whereitem{$\Psi_{\max,\text{ном}}$}{максимальный модуль потокосцепления
при номинальных параметрах, Вб$\cdot$виток;}
\whereitem{$N$}{целое число витков выбранной обмотки;}
\whereitem{$A_{\min,\text{ном}}$}{номинальная минимальная площадь
сечения выбранного магнитопровода, м$^2$.}
\end{whereblock}

Потокосцепление $\Psi_{\max,\text{ном}}$ получают по
формулам~\eqref{eq:flux-linkage-trajectory}
и~\eqref{eq:flux-linkage-global-maximum}: опорное значение задают по
характеристике $\Psi(I,T)$, после чего интегрируют фактическое $u_L(t)$.
Одна формула~\eqref{eq:flux-linkage} без опорного значения даёт только
размах потокосцепления и для определения $\Psi_{\max,\text{ном}}$
недостаточна. Только для предварительно линейного дросселя допускается
приближение
\begin{equation}
 \Psi_{\max,\text{ном}}\approx
 L_{\text{ном}} I_{\max,\text{ном}},
 \label{eq:psi-linear-preliminary}
\end{equation}
\begin{whereblock}
\whereitem{$L_{\text{ном}}$}{номинальная индуктивность дросселя, Гн;}
\whereitem{$I_{\max,\text{ном}}$}{максимальный ток в номинальном
расчётном режиме с учётом пульсации, А.}
\end{whereblock}

При выраженной зависимости $L(I,T)$ формулу
\eqref{eq:psi-linear-preliminary} не используют. Потокосцепление получают
из нелинейной характеристики $\Psi(I,T)$ магнитной цепи.

\subsubsection{Расчёт наихудшего сочетания}

Индекс «нх» означает наихудшее сочетание. Сначала формируют конечное
множество $\mathcal J$ предельных сочетаний допусков и рабочих режимов.
Для каждого сочетания заново рассчитывают ток, потокосцепление и среднюю
по минимальному сечению магнитную индукцию. Искомое значение равно
\begin{equation}
 B_{\max,\text{нх}}=
 \max_{j\in\mathcal J}
 \left[
 \frac{
 \left|\Psi_{\max,j}\right|
 }{
 N A_{\min,j}
 }
 \right],
 \label{eq:bmax-worst-case}
\end{equation}
\begin{whereblock}
\whereitem{$B_{\max,\text{нх}}$}{наибольшая аналитическая оценка магнитной
индукции, средней по минимальному сечению, среди всех рассчитанных
предельных сочетаний, Тл;}
\whereitem{$\mathcal J$}{множество проверяемых сочетаний допусков и
рабочих режимов;}
\whereitem{$j$}{номер сочетания;}
\whereitem{$\Psi_{\max,j}$}{максимальный модуль потокосцепления в
$j$-м сочетании, Вб$\cdot$виток;}
\whereitem{$A_{\min,j}$}{минимальная фактическая площадь сечения в
$j$-м сочетании, м$^2$.}
\end{whereblock}

В множество $\mathcal J$ включают:
\begin{enumerate}
\item максимальный рабочий и аварийный токи с допусками измерения и
регулирования;
\item минимальную допустимую площадь сечения магнитопровода;
\item предельные значения индуктивности и воздушного зазора;
\item наибольшую рабочую температуру;
\item наиболее неблагоприятную допустимую характеристику материала;
\item все установившиеся, переходные и аварийные режимы, заданные
техническим заданием.
\end{enumerate}

Отдельные процентные поправки не складывают с результатом
\eqref{eq:bmax-worst-case}, если соответствующие допуски уже вошли в
множество $\mathcal J$. Если перечень допусков ещё не определён,
$B_{\max,\text{нх}}$ получить нельзя и вариант остаётся предварительным.

\subsubsection{Определение численной и модельной неопределённости}

Численную составляющую определяют только после построения полевой модели
методом конечных элементов (Finite Element Method, FEM) и проверки
сходимости сетки. Изменение результата при последнем сгущении сетки равно
\begin{equation}
 \Delta B_{\text{сетка}}=
 \left|
 B_{\max}^{(m)}-
 B_{\max}^{(m-1)}
 \right|,
 \label{eq:b-mesh-uncertainty}
\end{equation}
\begin{whereblock}
\whereitem{$\Delta B_{\text{сетка}}$}{изменение локального максимума
магнитной индукции при последнем сгущении сетки, Тл;}
\whereitem{$B_{\max}^{(m)}$}{локальный максимум на последней, более
мелкой сетке, Тл;}
\whereitem{$B_{\max}^{(m-1)}$}{локальный максимум на предыдущей
сетке, Тл.}
\end{whereblock}

Если имеется измерение образца или подтверждённая контрольная модель,
расхождение определяют по формуле
\begin{equation}
 \Delta B_{\text{сопост}}=
 \left|
 B_{\max,\text{расч}}-
 B_{\max,\text{контр}}
 \right|,
 \label{eq:b-validation-uncertainty}
\end{equation}
\begin{whereblock}
\whereitem{$B_{\max,\text{расч}}$}{максимальная магнитная индукция,
полученная проверяемой моделью, Тл;}
\whereitem{$B_{\max,\text{контр}}$}{максимальная магнитная индукция,
полученная измерением образца или подтверждённой контрольной моделью,
Тл.}
\end{whereblock}

Верхнюю инженерную оценку, не являющуюся статистической стандартной
неопределённостью, принимают по правилу
\begin{equation}
 u_{B,\text{числ}}=
 \begin{cases}
  \Delta B_{\text{сетка}},
  &\text{без контрольного результата};\\
  \max\!\left(\Delta B_{\text{сетка}},
  \Delta B_{\text{сопост}}\right),
  &\text{при наличии контрольного результата},
 \end{cases}
 \label{eq:b-numerical-uncertainty}
\end{equation}
\begin{whereblock}
\whereitem{$u_{B,\text{числ}}$}{верхняя оценка остаточной численной и
модельной неопределённости, Тл;}
\whereitem{$\Delta B_{\text{сетка}}$}{изменение локального максимума при
последнем сгущении сетки, Тл;}
\whereitem{$\Delta B_{\text{сопост}}$}{расхождение с измерением образца
или подтверждённой контрольной моделью, Тл.}
\end{whereblock}

При отсутствии измерения формула~\eqref{eq:b-numerical-uncertainty}
учитывает только подтверждённую численную составляющую, а непроверенную
модельную составляющую компенсируют сохранением $k_B=0{,}60$. Если сетка
не проверена на сходимость, значение $u_{B,\text{числ}}$ не назначают
произвольным числом.

\subsubsection{Аналитическая проверка запаса по магнитной индукции}

Если окончательный критерий записывают относительно номинального
результата, абсолютный запас равен
\begin{equation}
 u_B=
 \max\left(
 0,\,
 B_{\max,\text{нх}}-B_{\max,\text{ном}}
 \right)
 +u_{B,\text{числ}},
 \label{eq:b-uncertainty}
\end{equation}
где $u_B$ --- абсолютный запас для применения совместно с номинальным
$B_{\max,\text{ном}}$, Тл.

Более прямой способ аналитической проверки имеет вид
\begin{equation}
 B_{\max,\text{нх}}+u_{B,\text{числ}}
 \leq k_B B_S(T_{\max}).
 \label{eq:b-worst-case-direct}
\end{equation}

Формулы~\eqref{eq:b-allow} и~\eqref{eq:b-uncertainty} с номинальным
$B_{\max,\text{ном}}$ и прямая проверка
\eqref{eq:b-worst-case-direct} являются двумя эквивалентными способами.
При использовании прямой проверки разность
$B_{\max,\text{нх}}-B_{\max,\text{ном}}$ повторно в $u_B$ не включают.
До проверки локального максимума методом конечных элементов результат
является аналитическим критерием отбраковки, а не окончательным
подтверждением конструкции.

\textbf{Пример аналитической проверки.}
При $B_S(T_{\max})=0{,}430\,\text{Тл}$, $k_B=0{,}65$,
$B_{\max,\text{ном}}=0{,}240\,\text{Тл}$,
$B_{\max,\text{нх}}=0{,}252\,\text{Тл}$ и
$u_{B,\text{числ}}=0{,}003\,\text{Тл}$ получают
\begin{align}
 u_B&=(0{,}252-0{,}240)+0{,}003
 =0{,}015\,\text{Тл},
 \label{eq:b-uncertainty-example}\\
 B_{\text{доп}}&=
 0{,}65\cdot0{,}430-0{,}015
 =0{,}2645\,\text{Тл}.
 \label{eq:b-allow-example}
\end{align}
Номинальный результат удовлетворяет условию
$0{,}240\leq0{,}2645\,\text{Тл}$. Эквивалентная прямая проверка даёт
$0{,}252+0{,}003=0{,}255\,\text{Тл}
\leq0{,}65\cdot0{,}430=0{,}2795\,\text{Тл}$.

\recommend Значение $k_B=0{,}65$ применяют после расчёта всех сочетаний
$\mathcal J$. Значение $k_B=0{,}70\ldots0{,}75$ допускается только после
измерения характеристики $\Psi(I,T)$ опытного образца и проверки
локального $B_{\max}$ на сходящейся сетке.

\criterion До выбора конструкции применяют только
формулу~\eqref{eq:b-allow-preliminary}. После определения конструкции
для аналитического отсева используют только один из двух вариантов: сравнивают
$B_{\max,\text{ном}}$ с $B_{\text{доп}}$ по формулам
\eqref{eq:b-allow} и~\eqref{eq:b-uncertainty} либо непосредственно
проверяют формулу~\eqref{eq:b-worst-case-direct}. Окончательную локальную
проверку выполняют после полевого расчёта по
формуле~\eqref{eq:acceptance-b-fem}. Температуру Кюри не
используют вместо $B_S(T)$ и не принимают как допустимую рабочую
температуру.

Марку материала выбирают в два прохода.

\textbf{Первый проход} формирует список из двух--четырёх кандидатов:
\begin{enumerate}
 \item исключают материалы, диапазон частоты которых не охватывает
 $f_{\text{с}}$;
 \item исключают материалы без кривых потерь для требуемых
 $\widehat B_{\text{AC}}$ и температуры;
 \item исключают материалы, для которых
 $B_{\text{доп,пред}}$ по формуле~\eqref{eq:b-allow-preliminary}
 заведомо недостаточно;
 \item проверяют наличие материала в требуемых формах и типоразмерах;
 \item оставляют несколько материалов с наименьшими ожидаемыми
 удельными потерями в рабочей области.
\end{enumerate}

\textbf{Второй проход} выполняют внутри перебора
подраздела~\ref{subsec:turn-gap-loop}: для каждой пары
«типоразмер~-- материал» рассчитывают фактические $B_{\max}$,
$B_{\text{DC}}$, $\widehat B_{\text{AC}}$, $P_{\text{Fe}}$ и
температуру. Только после этого фиксируют марку материала.

\begin{table}[H]
\caption{Пример чтения паспорта материала N88}
\label{tab:n88-datasheet-example}
\centering
\scriptsize
\begin{tabularx}{\textwidth}{>{\raggedright\arraybackslash}p{39mm}
>{\raggedright\arraybackslash}p{49mm} X}
\toprule
Паспортный параметр & Значение N88 & Как использовать \\
\midrule
Начальная проницаемость $\mu_i$ &
$1600\pm25\,\%$ при $25\,^{\circ}\text{С}$ &
для оценки замкнутого магнитопровода и малосигнального режима; не для
выбора индукции насыщения \\
Индукция насыщения $B_S$ &
$510\,\text{мТл}$ при $25\,^{\circ}\text{С}$;
$430\,\text{мТл}$ при $100\,^{\circ}\text{С}$;
$370\,\text{мТл}$ при $140\,^{\circ}\text{С}$;
условия измерения $H=1200\,\text{А/м}$, $f=10\,\text{кГц}$ &
до выбора конструкции вычислить $B_{\text{доп,пред}}$ по
формуле~\eqref{eq:b-allow-preliminary}; после расчёта конструкции
применить формулу~\eqref{eq:b-allow} \\
Оптимальный диапазон частоты &
$25\ldots500\,\text{кГц}$ &
N88 является кандидатом для $200\,\text{кГц}$, но не подтверждён данным
паспортом для $1\,\text{МГц}$ \\
Удельные потери $p_v$ &
$430\,\text{кВт/м}^3$ при $100\,\text{кГц}$,
$200\,\text{мТл}$, $100\,^{\circ}\text{С}$;
$360\,\text{кВт/м}^3$ при $300\,\text{кГц}$,
$100\,\text{мТл}$, $100\,^{\circ}\text{С}$;
$230\,\text{кВт/м}^3$ при $500\,\text{кГц}$,
$50\,\text{мТл}$, $100\,^{\circ}\text{С}$ &
использовать только точку или интерполяцию внутри представленных кривых;
указанная индукция на графиках TDK является $\widehat B$ \\
Постоянное подмагничивание &
графики при $25$ и $100\,^{\circ}\text{С}$, измеренные на ETD29 &
применять как проверку тенденции; точную $L_{\text{диф}}$ выбранного
зазорного сердечника получать из нелинейной модели \\
Температура Кюри $T_C$ & $260\,^{\circ}\text{С}$ &
не использовать как допустимую рабочую температуру \\
\bottomrule
\end{tabularx}
\end{table}

\source \mbox{TDK Electronics}, материал N88~\cite{tdkn88}. Условия измерения,
единицы и вид возбуждения являются частью паспортного значения.

\subsection{Предварительный габарит и проводник}
\label{subsec:size-conductor-preselection}

\stepgoal{Отсеять заведомо малые типоразмеры и получить первое, ещё не
окончательное, сечение меди.}

Коэффициент пика тока определяют по формуле
\begin{equation}
 k_{\text{п}}=
 \frac{I_{\max}}{I_{\text{действ}}},
 \label{eq:current-peak-factor}
\end{equation}
\begin{whereblock}
\whereitem{$k_{\text{п}}$}{коэффициент пика тока;}
\whereitem{$I_{\max}$}{максимальный ток, А;}
\whereitem{$I_{\text{действ}}$}{действующее значение полного тока, А.}
\end{whereblock}

Для первого выбора типоразмера используют размерно однородное базовое
соотношение произведения площадей
\begin{equation}
 A_{p,\text{треб}}=A_e A_w\ge
 \frac{2W_{\max}}
 {k_{\text{п}}k_w B_m J_{\text{доп}}},
 \label{eq:area-product-legostaev}
\end{equation}
\begin{whereblock}
\whereitem{$A_{p,\text{треб}}$}{минимально требуемое произведение площадей,
м$^4$;}
\whereitem{$A_e$}{эффективная площадь сечения магнитопровода, м$^2$;}
\whereitem{$A_w$}{полезная площадь окна магнитопровода, м$^2$;}
\whereitem{$W_{\max}$}{максимальная запасаемая энергия, Дж;}
\whereitem{$k_{\text{п}}$}{коэффициент пика тока;}
\whereitem{$k_w$}{предварительный коэффициент заполнения окна активной
медью;}
\whereitem{$B_m$}{предварительно допустимая максимальная индукция, Тл;}
\whereitem{$J_{\text{доп}}$}{предварительно допустимая плотность тока,
А/м$^2$.}
\end{whereblock}

Формула~\eqref{eq:area-product-legostaev} следует из энергетического
соотношения и ограничения площади меди; она соответствует базовому
выражению Легостаева~\cite[формула~(2.17)]{legostaev}. Эмпирическое
продолжение по формулам~(2.18), (2.19) и таблице~2.1 того же источника
зависит от принятой системы единиц и нормировки площади, поэтому в расчёте
по СИ его не применяют. Значения $J_{\text{доп}}$ и $k_w$ задают явно по
таблице~\ref{tab:kw-j}.

Для первого прохода принимают
$B_m=B_{\text{доп,пред}}(T_{\max})$. Формулу рассчитывают отдельно для
каждого режима с его энергией и коэффициентом пика, после чего в качестве
$A_{p,\text{треб}}$ выбирают наибольшую правую часть. Смешивать энергию
одного режима с коэффициентом пика другого режима нельзя.

\begin{table}[H]
\caption{Начальные значения заполнения окна и плотности тока}
\label{tab:kw-j}
\centering
\small
\begin{tabularx}{\textwidth}{>{\raggedright\arraybackslash}p{43mm}
>{\raggedleft\arraybackslash}p{37mm} X}
\toprule
Параметр & Начальное значение & Правило выбора \\
\midrule
$k_w$ для литцендрата & $0{,}20$ &
$0{,}25$ допускается только после выбора реального кабеля и построения
укладки; значение свыше $0{,}30$ не используют для первого расчёта \\
$k_w$ для круглого одножильного провода & $0{,}25$ &
$0{,}30\ldots0{,}35$ допускается для простой проверенной укладки \\
$J_{\text{доп}}$ при естественной конвекции &
$2{,}5\,\text{А/мм}^2$ &
$3{,}0\,\text{А/мм}^2$ допускается при температуре окружающей среды не выше
заданной и после тепловой проверки \\
$J_{\text{доп}}$ при принудительном охлаждении &
$3{,}0\,\text{А/мм}^2$ &
значение до $4{,}0\,\text{А/мм}^2$ назначают по рассчитанным потерям и
температуре, а не только по плотности тока \\
\bottomrule
\end{tabularx}
\end{table}

\source Значения $k_w$ согласованы с примерами
Легостаева~\cite{legostaev}, грубой верхней оценкой Sullivan и
Zhang~\cite{sullivan} и фактическим заполнением оптимизированных конструкций
Papamanolis с соавторами~\cite{papamanolis}. Они являются проектными
начальными значениями, а не свойствами материалов.

Предварительное требуемое произведение площадей определяют по
формуле~\eqref{eq:area-product-legostaev}. Для каждого каталожного
кандидата вычисляют
\begin{equation}
 A_{p,\text{кат}}=A_eA_w,
 \label{eq:catalog-area-product}
\end{equation}
где $A_{p,\text{кат}}$ --- каталожное произведение площадей, м$^4$;\par\noindent
$A_e$ --- эффективная площадь сечения магнитопровода, м$^2$;\par\noindent
$A_w$ --- доступная площадь окна, м$^2$.

Кандидаты с $A_{p,\text{кат}}$ меньше требуемого значения исключают.
Произведение площадей является только фильтром габарита; оно не
определяет конкретное число витков, зазор или марку материала.

До выбора окна определяют предварительное активное сечение меди:
\begin{equation}
 S_{\text{Cu,треб}}=
 \frac{I_{\text{действ}}}{J_{\text{доп}}},
 \label{eq:copper-area}
\end{equation}
где $S_{\text{Cu,треб}}$ --- требуемая суммарная площадь меди, м$^2$;\par\noindent
$I_{\text{действ}}$ --- действующее значение тока, А;\par\noindent
$J_{\text{доп}}$ --- предварительно допустимая плотность тока, А/м$^2$.

По $S_{\text{Cu,треб}}$ выбирают один или несколько реальных проводников
из каталога. Для каждого проводника фиксируют активную площадь
$S_{\text{Cu}}$ и наружную площадь с изоляцией $S_{\text{нар}}$.
Для планарной обмотки вместо $S_{\text{нар}}$ сразу задают допустимые
ширину, толщину, число слоёв и переходов по технологии печатной платы.

\subsection{Выбор конкретного типоразмера магнитопровода}
\label{subsec:core-size-choice}

На этом шаге выбирают не геометрическое название «E-сердечник», а
конкретный каталожный комплект в конкретном материале.

\begin{table}[H]
\caption{Использование параметров паспорта конкретного магнитопровода}
\label{tab:core-datasheet-map}
\centering
\scriptsize
\begin{tabularx}{\textwidth}{>{\raggedright\arraybackslash}p{30mm}
>{\raggedright\arraybackslash}p{45mm} X}
\toprule
Параметр паспорта & Для чего используется & Критерий \\
\midrule
$A_{\min}$ &
расчёт $B_{\max}$ и $\Delta B_{\text{пп}}$ &
обеспечивается запас по насыщению \\
$A_e$ &
расчёт зазора, среднего потока и произведения площадей &
используется значение для комплекта, а не одной половины \\
$V_e$ &
перевод удельных потерь в потери комплекта &
$P_{\text{Fe}}=p_vV_e$ входит в тепловой бюджет \\
$l_e$ и $\sum l/A$ &
магнитная модель замкнутого сердечника &
требуются при отсутствии готового $A_{L0}$ \\
$A_w$, размеры окна, каркас &
проверка размещения проводника и изоляции &
проверяются полезное окно и реальный чертёж намотки \\
$A_{L0}$ замкнутого комплекта &
вычитание магнитного сопротивления феррита при расчёте зазора &
используется только после выбора типоразмера; не выбирает габарит \\
$\mu_e$ &
проверка согласованности $A_{L0}$ и геометрии &
не является рабочей проницаемостью сердечника с большим зазором \\
$P_V$ или $P_{\text{Fe}}$ в строке заказа &
контрольная точка потерь готового комплекта &
действительна только при указанных $f$, $\widehat B$ и $T$ \\
Код заказа &
однозначное задание материала, формы и исполнения &
обязательно переносится в спецификацию \\
\bottomrule
\end{tabularx}
\end{table}

Если $A_{L0}$ не указан, но приведены эффективная проницаемость
$\mu_e$ и геометрический коэффициент $\sum l/A$, его оценивают по
формуле
\begin{equation}
 A_{L0}\approx
 \frac{\mu_0\mu_e}{\displaystyle\sum l/A},
 \label{eq:al0-from-geometry}
\end{equation}
\begin{whereblock}
\whereitem{$A_{L0}$}{коэффициент индуктивности замкнутого комплекта,
Гн/виток$^2$;}
\whereitem{$\mu_0=1{,}256\,637\,061\,27\cdot10^{-6}\,\text{Гн/м}$}{магнитная
постоянная по CODATA 2022~\cite{codata};}
\whereitem{$\mu_e$}{эффективная относительная магнитная проницаемость
комплекта;}
\whereitem{$\sum l/A$}{сумма отношений длины участков магнитной цепи к их
сечению, м$^{-1}$.}
\end{whereblock}

Для инженерного расчёта с относительной точностью хуже $10^{-8}$ допускается
округление $\mu_0\approx4\pi\cdot10^{-7}\,\text{Гн/м}$.

Для E120/71/40 в N88 паспорт содержит
$\sum l/A=0{,}166\,\text{мм}^{-1}$ и $\mu_e=1650$:
\begin{equation}
 A_{L0}\approx
 \frac{4\pi\cdot10^{-7}\cdot1650}
 {0{,}166\cdot10^3}
 =12{,}49\,\text{мкГн/виток}^2,
 \label{eq:al0-e120-check}
\end{equation}
что согласуется с паспортными
$12500\,\text{нГн/виток}^2$. Для расчёта зазора используют
не оценку~\eqref{eq:al0-from-geometry}, а непосредственно указанное
изготовителем $A_{L0}$ с его допуском.

Например, паспорт E120/71/40~\cite{tdkcore} для комплекта N88 указывает:
$A_e=1782\,\text{мм}^2$, $A_{\min}=1720\,\text{мм}^2$,
$l_e=295\,\text{мм}$, $V_e=525690\,\text{мм}^3$,
$A_{L0}=12500\,\text{нГн/виток}^2$ с допуском
$+30/-20\,\%$ и код заказа \texttt{B66614G0000X188}.
Площадь окна в этом паспорте не приведена; её следует получить из
размерного чертежа и чертежа каркаса. Потери менее $9{,}5\,\text{Вт}$,
указанные в строке заказа, относятся только к
$\widehat B=100\,\text{мТл}$, $25\,\text{кГц}$ и
$120\,^{\circ}\text{С}$ и не могут быть перенесены на другую рабочую
точку.

Кандидат оставляют для перебора, если одновременно:
\begin{enumerate}
 \item выполняются габаритные ограничения;
 \item $A_{p,\text{кат}}$ не меньше предварительно требуемого;
 \item предварительная оценка показывает, что в окне размещается требуемое
 по ограничению магнитной индукции число витков выбранного проводника;
 \item существует способ изготовления и контроля требуемого зазора;
 \item материал имеет данные потерь в рабочем диапазоне.
\end{enumerate}

\subsection{Совместный перебор числа витков и зазора}
\label{subsec:turn-gap-loop}

\textbf{Цель шага} --- для каждого конкретного комплекта получить
конечные целые $N$ и физический $g$, не назначая их независимо.

\begin{figure}[htbp]
\centering
\includegraphics[width=0.92\textwidth]{gap_arrangements.png}
\caption{Варианты размещения одного, двух и трёх дискретных зазоров
в магнитопроводе (по Liao и Chen~\cite[рисунок~2]{liao})}
\label{fig:gap-arrangements}
\end{figure}

На рисунке~\ref{fig:gap-arrangements} показано, что одинаковый суммарный
зазор можно реализовать различным числом участков. Разделение одного
большого зазора на несколько меньших может уменьшить область повышенного
поля, но эффект не является монотонным и зависит от взаимного положения
зазоров и обмотки. В расчёте учитывают каждый физический зазор,
технологический разброс его длины и расположение меди относительно всех
кромок.

Для одного эквивалентного зазора предварительный коэффициент выпучивания
поля определяют по формуле
\begin{equation}
 F_f(g)=1+q\frac{g}{\sqrt{A_e}}
 \ln\left(\frac{2G}{g}\right),
 \label{eq:fringing-factor}
\end{equation}
\begin{whereblock}
\whereitem{$F_f$}{коэффициент выпучивания поля у зазора;}
\whereitem{$q$}{коэффициент формы поперечного сечения магнитопровода;}
\whereitem{$g$}{суммарная длина рассматриваемого зазора, м;}
\whereitem{$A_e$}{эффективная площадь сечения магнитопровода, м$^2$;}
\whereitem{$G$}{размер магнитопровода или зоны намотки, принятый в модели
источника, м.}
\end{whereblock}

\begin{table}[H]
\caption{Выбор коэффициента формы $q$}
\label{tab:q-factor}
\centering
\small
\begin{tabularx}{\textwidth}{>{\raggedright\arraybackslash}p{47mm}
>{\raggedleft\arraybackslash}p{30mm} X}
\toprule
Форма центрального стержня & Диапазон $q$ & Начальное значение \\
\midrule
Круглый стержень PQ или ETD & $0{,}85\ldots0{,}95$ &
$0{,}90$; обе границы проверяют, если запас по индуктивности менее
$10\,\%$ \\
Прямоугольный стержень E, EE или EI & $1{,}00\ldots1{,}10$ &
$1{,}05$; без FEM рассчитывают обе границы \\
\bottomrule
\end{tabularx}
\end{table}

\source Формула и диапазоны $q$ приведены Liao и
Chen~\cite[формула~(21)]{liao}. Она является предварительной: для нескольких
зазоров, большого отношения $g/\sqrt{A_e}$ или близкой широкой меди
$F_f$ получают из полевой модели или по измеренному $A_L$.

\criterion В источнике параметр $G$ обозначен как характерный размер
магнитопровода. Для выбранной геометрии используют только размер, однозначно
заданный исходной моделью или подтверждённый контрольным расчётом. Подстановка
произвольного габарита в формулу~\eqref{eq:fringing-factor} недопустима; при
неоднозначности рассчитывают границы либо сразу переходят к полевой модели.

\begin{figure}[htbp]
\centering
\begin{minipage}[t]{0.48\textwidth}
\centering
\includegraphics[width=\linewidth]{one_gap_magnetic_circuit.png}\\
\small а) один зазор
\end{minipage}\hfill
\begin{minipage}[t]{0.48\textwidth}
\centering
\includegraphics[width=\linewidth]{two_gap_magnetic_circuit.png}\\
\small б) два зазора
\end{minipage}
\caption{Связь геометрии магнитопровода с эквивалентной магнитной цепью
для одного и двух зазоров (по Liao и Chen~\cite[рисунки~3, 4]{liao})}
\label{fig:gap-magnetic-circuits}
\end{figure}

Эквивалентные схемы на рисунке~\ref{fig:gap-magnetic-circuits} поясняют,
почему длины зазоров нельзя только сложить и подставить в формулу
однородной магнитной цепи. Поток рассеяния каждого зазора замыкается по
собственному пути, а взаимное влияние путей зависит от расстояния между
зазорами. Аналитический расчёт используют для первого приближения, после
чего геометрию проверяют в полевой модели.

\begin{figure}[htbp]
\centering
\begin{minipage}[t]{0.48\textwidth}
\centering
\includegraphics[width=\linewidth]{fringing_field_one_gap.png}\\
\small а) один сосредоточенный зазор
\end{minipage}\hfill
\begin{minipage}[t]{0.48\textwidth}
\centering
\includegraphics[width=\linewidth]{fringing_field_two_gaps.png}\\
\small б) зазор разделён на несколько участков
\end{minipage}
\caption{Распределение магнитного поля в окрестности зазора при разных
вариантах его размещения (по Liao и Chen~\cite[рисунки~9, 10]{liao})}
\label{fig:gap-fringing-comparison}
\end{figure}

Карты на рисунке~\ref{fig:gap-fringing-comparison} следует читать по
площади области с повышенной напряжённостью поля, а не только по
максимальному значению. Чем больше эта область пересекает медь обмотки или
окружающие металлические детали, тем выше добавочные вихревые потери.

Минимальное число витков по насыщению определяют по формуле
\begin{equation}
 N_B=
 \left\lceil
 \frac{\left|\Psi\right|_{\max}}
 {A_{\min}B_{\text{доп,пред}}}
 \right\rceil,
 \label{eq:n-min}
\end{equation}
где $N_B$ --- минимальное число витков по предварительному ограничению
магнитной индукции;\par\noindent
$B_{\text{доп,пред}}$ --- предварительно допустимая магнитная индукция при
максимальной температуре по формуле~\eqref{eq:b-allow-preliminary},
Тл;\par\noindent
$\lceil\,\rceil$ --- округление вверх до целого числа.

Различают коэффициент заполнения окна медью
\begin{equation}
 k_{\text{Cu}}=
 \frac{N S_{\text{Cu}}}{A_{w,\text{пол}}},
 \label{eq:window-copper-fill}
\end{equation}
и коэффициент геометрической занятости окна
\begin{equation}
 k_{\text{зан}}=
 \frac{N S_{\text{нар}}}{A_{w,\text{пол}}},
 \label{eq:window-occupancy}
\end{equation}
где $k_{\text{Cu}}$ --- отношение площади меди всех витков к полезной
площади окна;\par\noindent
$k_{\text{зан}}$ --- отношение наружной площади изолированных
проводников к полезной площади окна;\par\noindent
$A_{w,\text{пол}}$ --- полезная площадь окна после вычета каркаса,
межобмоточной изоляции и технологически недоступных зон, м$^2$.

Максимальное число витков по геометрии определяют по формуле
\begin{equation}
 N_w=
 \left\lfloor
 \frac{k_{\text{зан,доп}}A_{w,\text{пол}}}
 {S_{\text{нар}}}
 \right\rfloor,
 \label{eq:n-window-max}
\end{equation}
где $N_w$ --- максимальное число витков по окну;\par\noindent
$k_{\text{зан,доп}}$ --- допустимая геометрическая занятость окна;\par\noindent
$\lfloor\,\rfloor$ --- округление вниз до целого числа.

\recommend Для первого чертежа проволочной обмотки принимают
$k_{\text{зан,доп}}=0{,}65\ldots0{,}75$. Значение до $0{,}80$ допускают
только при проверенной машинной укладке. Это проектная рекомендация,
а не свойство материала. Коэффициент $k_{\text{Cu}}$ сопоставляют с
таблицей~\ref{tab:kw-j}; смешивать $k_{\text{Cu}}$ и $k_{\text{зан}}$
не допускается.

Если $N_B>N_w$, кандидатный типоразмер немедленно исключают. Если
$N_B\leq N_w$, последовательно проверяют все практически допустимые
целые $N$ от $N_B$ до $N_w$ либо до меньшей границы, заданной длиной
провода и допустимой ёмкостью.

Для каждого $N$ требуемый коэффициент индуктивности с зазором равен
\begin{equation}
 A_L(N)=\frac{L_{\text{ном}}}{N^2},
 \label{eq:al-required}
\end{equation}
где $A_L(N)$ --- требуемый коэффициент индуктивности, Гн/виток$^2$;\par\noindent
$L_{\text{ном}}$ --- номинальная индуктивность, Гн.

Вклад зазора относительно магнитного сопротивления феррита оценивают по
формуле
\begin{equation}
 k_{\mathcal R}=
 \frac{\mathcal R_g}{\mathcal R_c}
 =\frac{A_{L0}}{A_L}-1,
 \label{eq:gap-reluctance-ratio}
\end{equation}
где $k_{\mathcal R}$ --- отношение магнитного сопротивления зазора
$\mathcal R_g$ к магнитному сопротивлению феррита $\mathcal R_c$;\par\noindent
$A_{L0}$ --- коэффициент индуктивности замкнутого комплекта без зазора,
Гн/виток$^2$.

\recommend Для дросселя, индуктивность которого должна слабо зависеть
от разброса проницаемости, на первом проходе желательно
$k_{\mathcal R}\geq4$. При меньшем отношении обязательно рассчитывают
чувствительность $L$ к допуску $A_{L0}$ и температуре. Число $4$
является проектным критерием устойчивости, а не универсальной границей.

При фиксированном зазоре влияние допуска замкнутого сердечника
предварительно оценивают по формуле
\begin{equation}
 A_L(A_{L0,\text{ф}})=
 \left[
 \frac{1}{A_L(N)}
 -\frac{1}{A_{L0,\text{ном}}}
 +\frac{1}{A_{L0,\text{ф}}}
 \right]^{-1},
 \label{eq:al0-tolerance-sensitivity}
\end{equation}
где $A_L(A_{L0,\text{ф}})$ --- коэффициент индуктивности при
фактическом $A_{L0}$, Гн/виток$^2$;\par\noindent
$A_{L0,\text{ном}}$ --- номинальный коэффициент замкнутого комплекта,
Гн/виток$^2$;\par\noindent
$A_{L0,\text{ф}}$ --- минимальное или максимальное паспортное значение
замкнутого комплекта, Гн/виток$^2$.

После вычисления $L=N^2A_L(A_{L0,\text{ф}})$ обе границы должны
укладываться в допуск индуктивности. Формула является оценкой
чувствительности линейной магнитной цепи; допуск физического зазора,
выпучивание и нелинейность проверяют отдельно.

Без учёта выпучивания поля начальную оценку суммарного физического
зазора получают по формуле
\begin{equation}
 g_0=\mu_0 A_e
 \left(
 \frac{1}{A_L}-\frac{1}{A_{L0}}
 \right),
 \label{eq:gap-ideal}
\end{equation}
где $g_0$ --- начальная оценка суммарного зазора, м.

С учётом коэффициента выпучивания $F_f$ зазор находят
итерационно:
\begin{equation}
 g_{r+1}=\mu_0 A_e F_f(g_r)
 \left(
 \frac{1}{A_L}-\frac{1}{A_{L0}}
 \right),
 \label{eq:gap-fringing}
\end{equation}
где $g_r$, $g_{r+1}$ --- суммарный зазор на текущей и следующей
итерациях, м;\par\noindent
$r$ --- номер итерации;\par\noindent
$F_f$ --- коэффициент выпучивания из
формулы~\eqref{eq:fringing-factor}.

Итерацию прекращают при выполнении
\begin{equation}
 \frac{\left|g_{r+1}-g_r\right|}{g_{r+1}}
 \leq\varepsilon_g,
 \label{eq:gap-iteration-stop}
\end{equation}
где $\varepsilon_g$ --- заданная относительная погрешность расчёта
зазора. Для аналитического выбора принимают
$\varepsilon_g=10^{-3}$; после этого зазор всё равно уточняют по FEM и
измерению $A_L$.

Если зазор разделён на $n_g$ одинаковых участков, начальная длина одного
участка равна
\begin{equation}
 g_i=\frac{g}{n_g},
 \label{eq:gap-split}
\end{equation}
где $g_i$ --- длина одного участка зазора, м;\par\noindent
$n_g$ --- число участков зазора.

Предварительные потери магнитопровода определяют непосредственно из
паспортной плотности потерь:
\begin{equation}
 P_{\text{Fe,пред}}=
 p_v\!\left(f_{\text{с}},\widehat B_{\text{AC}},
 B_{\text{DC}},T\right)V_e,
 \label{eq:core-loss-preliminary}
\end{equation}
где $P_{\text{Fe,пред}}$ --- предварительные потери магнитопровода, Вт;\par\noindent
$p_v$ --- удельные объёмные потери материала в расчётной рабочей точке,
Вт/м$^3$;\par\noindent
$V_e$ --- эффективный объём магнитопровода, м$^3$.

Предварительное сопротивление обмотки постоянному току равно
\begin{equation}
 R_{\text{DC,пред}}=
 \rho_{\text{Cu}}(T)
 \frac{N l_{\text{вит,ср}}}{S_{\text{Cu}}},
 \label{eq:rdc-preliminary}
\end{equation}
где $R_{\text{DC,пред}}$ --- предварительное сопротивление обмотки, Ом;\par\noindent
$l_{\text{вит,ср}}$ --- средняя длина витка по чертежу кандидата, м.

Для первого сравнения вариантов используют
\begin{equation}
 P_{\text{Cu,пред}}=
 I_{\text{ср}}^{\,2}R_{\text{DC,пред}}
 +\sum_{h=1}^{H}
 I_{h,\text{действ}}^{\,2}
 R_{\text{AC,пред}}(f_h),
 \label{eq:copper-loss-preliminary}
\end{equation}
где $P_{\text{Cu,пред}}$ --- предварительные потери обмотки, Вт;\par\noindent
$H$ --- число учитываемых гармоник;\par\noindent
$I_{h,\text{действ}}$ --- действующее значение $h$-й гармоники тока,
А;\par\noindent
$R_{\text{AC,пред}}(f_h)$ --- предварительное сопротивление переменному
току на частоте $f_h$ для фактической геометрии обмотки, Ом;\par\noindent
$f_h$ --- частота $h$-й гармоники, Гц.

Если паспорт не содержит $p_v$ в требуемой рабочей точке или отсутствует
модель $R_{\text{AC,пред}}$, вариант не объявляют рассчитанным: используют
консервативную измеренную границу либо сохраняют его как
«требующий верификации». Подробный гармонический расчёт выполняют по
разделу~\ref{sec:losses} после выбора короткого списка. Подстановка
$I_{\text{действ}}^2R_{\text{DC,пред}}$ даёт только нижнюю оценку потерь,
если скин-эффект, эффект близости или поле зазора увеличивают сопротивление.

Для каждого сочетания «материал~-- типоразмер~-- $N$~-- $g$»
заполняют таблицу~\ref{tab:turn-gap-iteration}.

\begin{table}[H]
\caption{Обязательные результаты итерации числа витков и зазора}
\label{tab:turn-gap-iteration}
\centering
\scriptsize
\begin{tabularx}{\textwidth}{>{\raggedright\arraybackslash}p{35mm}
>{\raggedright\arraybackslash}p{53mm} X}
\toprule
Рассчитываемая величина & Формула или источник & Решение \\
\midrule
$A_L(N)$ & формула~\eqref{eq:al-required} &
$A_L<A_{L0,\min}$; зазор физически реализуем \\
$k_{\mathcal R}$ & формула~\eqref{eq:gap-reluctance-ratio} &
индуктивность не чрезмерно чувствительна к проницаемости \\
$g$ и число участков & формулы~\eqref{eq:gap-fringing},
\eqref{eq:gap-split} &
допуск зазора обеспечивает допуск $L$; поле рассеяния допустимо \\
$B_{\max}$, $B_{\text{DC}}$, $\widehat B_{\text{AC}}$ &
формулы~\eqref{eq:bmax-candidate}--\eqref{eq:bdc} &
выполнены насыщение и паспортный диапазон потерь \\
$P_{\text{Fe,пред}}$ & формула~\eqref{eq:core-loss-preliminary} &
потери входят в тепловой бюджет \\
$k_{\text{Cu}}$, $k_{\text{зан}}$ &
формулы~\eqref{eq:window-copper-fill},
\eqref{eq:window-occupancy} &
обмотка размещается по чертежу \\
$R_{\text{DC,пред}}$, $P_{\text{Cu,пред}}$ &
формулы~\eqref{eq:rdc-preliminary},
\eqref{eq:copper-loss-preliminary} &
рост потерь от добавления витков оправдан снижением потерь сердечника \\
Длина витка, ёмкость, поле зазора &
чертёж, FEM и раздел~\ref{sec:parasitics} &
нет неприемлемого ухудшения электромагнитных помех \\
\bottomrule
\end{tabularx}
\end{table}

Направление изменения конструкции принимают по
таблице~\ref{tab:iteration-direction}.

\begin{table}[H]
\caption{Правила изменения числа витков и типоразмера}
\label{tab:iteration-direction}
\centering
\small
\begin{tabularx}{\textwidth}{>{\raggedright\arraybackslash}p{48mm} X}
\toprule
Невыполненное условие & Действие \\
\midrule
$B_{\max}$ или $P_{\text{Fe}}$ слишком велики &
увеличить $N$; при нехватке окна перейти к большему $A_e$ и $A_w$ либо к
материалу с меньшими потерями \\
$P_{\text{Cu}}$, длина провода или ёмкость слишком велики &
уменьшить $N$, если сохраняется запас по $B$; иначе увеличить типоразмер
или активное сечение проводника \\
Зазор слишком мал для воспроизводимого изготовления &
увеличить $N$, поскольку требуемый $A_L$ уменьшится и зазор возрастёт \\
Зазор слишком велик и создаёт сильное поле рассеяния &
уменьшить $N$, увеличить $A_e$ либо разделить зазор на несколько участков \\
$N_B>N_w$ &
увеличить окно или перейти к проводнику с меньшей наружной площадью; не
уменьшать число витков ниже ограничения по индукции \\
Нет данных материала для рассчитанных $f$, $\widehat B_{\text{AC}}$, $T$ &
выбрать другую марку материала либо выполнить собственное измерение \\
\bottomrule
\end{tabularx}
\end{table}

\source Формулы выпучивания и диапазоны коэффициента $q$ приведены Liao
и Chen~\cite{liao}. Влияние положения и числа зазоров на потери обмотки
рассмотрено в \cite{papamanolis,ewald_litz,meng}.

\subsection{Формирование короткого списка до проектирования обмотки}

На этом этапе вариант не фиксируют окончательно и не требуют результата
FEM. В коротком списке оставляют два--четыре аналитически допустимых
сочетания, для которых выполняется
\begin{equation}
 B_{\max,\text{анал}}
 \leq B_{\text{доп,пред}}\!\left(T_{\max}\right),
 \label{eq:b-acceptance}
\end{equation}
\begin{whereblock}
\whereitem{$B_{\max,\text{анал}}$}{наибольшая аналитическая оценка
магнитной индукции, средней по минимальному сечению, для наихудшего
сочетания тока, допуска индуктивности и площади сечения, Тл;}
\whereitem{$B_{\text{доп,пред}}(T_{\max})$}{предварительно допустимая
индукция при максимальной температуре по
формуле~\eqref{eq:b-allow-preliminary}, Тл.}
\end{whereblock}

Дополнительно проверяют существование технологического зазора и возможность
разместить требуемую активную площадь меди. Если условия не выполнены ни для
одного $N$, переходят к следующей паре «типоразмер~-- материал» и повторяют
перебор.

В карточку каждого кандидата передают:
\begin{enumerate}
 \item изготовителя, форму, типоразмер, марку материала и код заказа
 обеих половин или комплекта;
 \item кандидатное число витков;
 \item суммарный зазор, число участков, размер каждого участка и допуск;
 \item требуемую площадь меди, полезные размеры окна, каркас и
изоляционные ограничения;
 \item предварительные $B_{\max}$, $P_{\text{Fe}}$ и бюджет
$P_{\text{Cu}}$.
\end{enumerate}

Полное обозначение провода и послойная укладка появляются только после
раздела~\ref{sec:winding}; затем геометрию передают в FEM по
разделу~\ref{sec:fem}. Окончательную спецификацию выпускают после проверки
всех ограничений раздела~\ref{sec:optimization}.

\section{ГИБРИДНЫЕ МАГНИТОПРОВОДЫ С МНОЖЕСТВЕННЫМИ ЗАЗОРАМИ}
\label{sec:hybrid-core}

\criterion Раздел является дополнительной ветвью расчёта. Если обычный
ферритовый или порошковый магнитопровод удовлетворяет ограничениям по
индуктивности, потерям, температуре и габариту, настоящий раздел пропускают
и переходят непосредственно к разделу~\ref{sec:winding}.

\subsection{Назначение и варианты конструкции}

\phys Гибридный магнитопровод объединяет ферритовое основное тело с одной
или несколькими вставками центрального стержня из материала с меньшей
относительной магнитной проницаемостью. Вставка Sendust создаёт объёмное
магнитное сопротивление, поэтому часть энергии запасается не только в
дискретных воздушных зазорах, но и в порошковом материале. Ферритовые ярма и
наружные стержни сохраняют малые высокочастотные потери, а характеристика
$L(I)$ центральной ветви становится более пологой.

Подход применяют, когда высота, каркас и наружный контур магнитопровода уже
заданы, а полностью ферритовый вариант не обеспечивает требуемую
дифференциальную индуктивность при максимальном токе. Использование Sendust не
назначают по умолчанию: гибрид сравнивают с увеличением типоразмера,
полностью ферритовым многозазорным вариантом и готовым порошковым
магнитопроводом.

В работе Liao и Chen исследованы пять вариантов, приведённых в
таблице~\ref{tab:hybrid-topologies}. Обозначение FE соответствует ферритовой
вставке, SE --- вставке Sendust~\cite[рисунки~2--5]{liao}.

\begin{table}[H]
\caption{Варианты центрального стержня и воздушных зазоров}
\label{tab:hybrid-topologies}
\centering
\small
\begin{tabularx}{\textwidth}{>{\raggedright\arraybackslash}p{20mm}
>{\raggedright\arraybackslash}p{46mm}
>{\centering\arraybackslash}p{23mm} X}
\toprule
Тип & Центральная ветвь & Число зазоров & Назначение варианта \\
\midrule
I & Две штатные ферритовые половины & 1 &
Классическая топология и база сравнения \\
II-FE & Одна отдельная ферритовая вставка & 2 &
Проверка влияния разделения зазора без смены материала \\
II-SE & Одна отдельная вставка Sendust & 2 &
Увеличение тока при минимальном числе дополнительных деталей \\
III-SE & Две отдельные вставки Sendust & 3 &
Дополнительное распределение магнитного сопротивления и краевого поля \\
III-FE & Две отдельные ферритовые вставки & 3 &
Контрольный полностью ферритовый многозазорный вариант \\
\bottomrule
\end{tabularx}
\end{table}

\subsection{Расчёт эквивалентной магнитной цепи}

\stepgoal{Для заданных материалов, размеров вставок, числа витков и требуемой
индуктивности определить суммарное магнитное сопротивление и длины отдельных
воздушных зазоров.}

Предварительное суммарное магнитное сопротивление составной последовательной
ветви определяют по формуле
\begin{equation}
 \mathcal R_{\Sigma}=
 \mathcal R_{\text{Fe}}+
 \sum_{k=1}^{n_s}\mathcal R_{s,k}+
 \sum_{i=1}^{n_g}\mathcal R_{g,i},
 \label{eq:hybrid-reluctance-total}
\end{equation}
где
\begin{align}
 \mathcal R_{\text{Fe}}&=
 \frac{l_{\text{Fe}}}
 {\mu_0\mu_{r,\text{Fe}}A_{\text{Fe}}},
 \label{eq:hybrid-reluctance-ferrite}\\
 \mathcal R_{s,k}&=
 \frac{l_{s,k}}
 {\mu_0\mu_{r,s,k}A_{s,k}},
 \label{eq:hybrid-reluctance-segment}\\
 \mathcal R_{g,i}&=
 \frac{g_i}
 {\mu_0F_{f,i}A_{g,i}}.
 \label{eq:hybrid-reluctance-gap}
\end{align}
\begin{whereblock}
\whereitem{$\mathcal R_{\Sigma}$}{суммарное магнитное сопротивление основной
магнитной цепи, А/Вб;}
\whereitem{$\mathcal R_{\text{Fe}}$}{магнитное сопротивление оставшейся
ферритовой части, А/Вб;}
\whereitem{$\mathcal R_{s,k}$}{магнитное сопротивление $k$-й вставки
центрального стержня, А/Вб;}
\whereitem{$\mathcal R_{g,i}$}{магнитное сопротивление $i$-го воздушного
зазора, А/Вб;}
\whereitem{$n_s$}{число твёрдых вставок центрального стержня;}
\whereitem{$n_g$}{число дискретных воздушных зазоров;}
\whereitem{$l_{\text{Fe}}$}{длина пути в феррите без длин вставок и
воздушных зазоров, м;}
\whereitem{$l_{s,k}$}{длина $k$-й вставки, м;}
\whereitem{$g_i$}{физическая длина $i$-го зазора, м;}
\whereitem{$A_{\text{Fe}}$, $A_{s,k}$, $A_{g,i}$}{площади поперечного
сечения феррита, вставки и зазора соответственно, м$^2$;}
\whereitem{$\mu_{r,\text{Fe}}$, $\mu_{r,s,k}$}{относительные магнитные
проницаемости феррита и материала вставки в рассматриваемой рабочей точке;}
\whereitem{$F_{f,i}$}{коэффициент выпучивания поля $i$-го зазора.}
\end{whereblock}

В линейном первом приближении секущая индуктивность равна
\begin{equation}
 L_{\text{сек}}\approx
 \frac{N^2}{\mathcal R_{\Sigma}},
 \label{eq:hybrid-inductance}
\end{equation}
где $N$ --- число витков обмотки.

Для заданных $N$ и $L_{\text{ном}}$ требуемые зазоры должны удовлетворять
условию
\begin{equation}
 \sum_{i=1}^{n_g}\frac{g_i}{F_{f,i}A_{g,i}}=
 \mu_0\left[
 \frac{N^2}{L_{\text{ном}}}
 -\frac{l_{\text{Fe}}}{\mu_0\mu_{r,\text{Fe}}A_{\text{Fe}}}
 -\sum_{k=1}^{n_s}
 \frac{l_{s,k}}{\mu_0\mu_{r,s,k}A_{s,k}}
 \right].
 \label{eq:hybrid-gap-balance}
\end{equation}
\begin{whereblock}
\whereitem{$L_{\text{ном}}$}{требуемая номинальная индуктивность, Гн.}
\end{whereblock}

Если правая часть формулы~\eqref{eq:hybrid-gap-balance} неположительна,
вставки уже создают магнитное сопротивление, превышающее требуемое. Такой
вариант нельзя исправить уменьшением воздушных зазоров: изменяют число
витков, длину или материал вставки либо площадь её сечения.

Для каждого зазора начальный коэффициент выпучивания поля рассчитывают
отдельно:
\begin{equation}
 F_{f,i}=1+q_i\frac{g_i}{\sqrt{A_{g,i}}}
 \ln\left(\frac{2G_i}{g_i}\right),
 \label{eq:hybrid-fringing-factor}
\end{equation}
\begin{whereblock}
\whereitem{$q_i$}{коэффициент формы поперечного сечения у $i$-го зазора;}
\whereitem{$G_i$}{характерный размер зоны намотки или магнитопровода у
$i$-го зазора, м.}
\end{whereblock}

Формулы~\eqref{eq:hybrid-reluctance-total}--
\eqref{eq:hybrid-fringing-factor} обобщают магнитные цепи одного, двух и
трёх зазоров из работы Liao и Chen~\cite[формулы~(3)--(21)]{liao}.
Коэффициент $F_{f,i}$ зависит от искомого $g_i$, поэтому систему решают
итерационно. При пересечении областей краевого поля соседних зазоров
скалярные коэффициенты перестают быть независимыми; окончательные значения
$g_i$, $L_{\text{сек}}$ и $L_{\text{диф}}$ получают из нелинейной FEM-модели
и измерений.

\subsection{Последовательность проектирования}

\method Гибридный вариант рассчитывают как дополнительную ветвь общей
методики:
\begin{enumerate}
 \item по электрическому режиму определяют $L_{\text{ном}}$,
 $L_{\text{диф,min}}$, $I_{\max}$, $I_{\text{действ}}$ и максимальную
 запасаемую энергию;
 \item выбирают феррит основного тела и получают из измерений или данных
 изготовителя зависимости $B(H,T)$ и $p_v(f,B_{\text{AC}},B_{\text{DC}},T)$
 для феррита и материала вставки;
 \item формируют варианты I, II-FE, II-SE, III-SE и III-FE в одинаковом
 наружном габарите и с одинаковым доступным окном;
 \item для каждого целого $N$ задают длины вставок и решают
 формулы~\eqref{eq:hybrid-reluctance-total} и
 \eqref{eq:hybrid-gap-balance} относительно $g_i$;
 \item пересчитывают $F_{f,i}$, индукцию в каждом материале, заполнение
 окна и предварительные потери; итерацию повторяют до сходимости;
 \item проектируют одну и ту же либо эквивалентную по активной площади
 обмотку для всех сравниваемых вариантов;
 \item в нелинейной магнитостатической модели получают $\Psi(I,T)$,
 $L_{\text{сек}}(I,T)$, $L_{\text{диф}}(I,T)$, локальный $B_{\max}$ и
 распределение поля каждого зазора;
 \item в частотной и тепловой моделях определяют потери феррита, вставок,
 обмотки и металлических деталей, затем выполняют совместную температурную
 итерацию;
 \item изготавливают варианты с одинаковыми условиями намотки и измеряют
 $L(I,T)$, потери, температуры и коэффициент полезного действия устройства.
\end{enumerate}

\criterion Гибрид принимают только при одновременном выполнении требований
к минимальной дифференциальной индуктивности, потерям, температуре,
габаритам, допускам всех $g_i$ и механической фиксации вставок. Рост тока
насыщения сам по себе не является достаточным критерием выбора.

\subsection{Сравнение тока насыщения и краевого поля}

В эксперименте Liao и Chen использованы PQ3530, номинальная индуктивность
$440\,\text{мкГн}\pm7\,\%$, частота коммутации $65\,\text{кГц}$ и корректор
коэффициента мощности (power factor correction, PFC) мощностью
$500\,\text{Вт}$. Током насыщения авторы называют ток, при котором
индуктивность выходит за установленный допуск, а не ток достижения
справочной индукции насыщения материала. Результаты относительно
классического типа I приведены в
таблице~\ref{tab:hybrid-saturation-results}.

\begin{table}[H]
\caption{Ток насыщения, мощность и краевое поле относительно типа I}
\label{tab:hybrid-saturation-results}
\centering
\scriptsize
\begin{tabularx}{\textwidth}{>{\raggedright\arraybackslash}p{17mm}
>{\raggedleft\arraybackslash}p{18mm}
>{\raggedleft\arraybackslash}p{21mm}
>{\raggedleft\arraybackslash}p{21mm}
>{\raggedleft\arraybackslash}p{22mm}
>{\raggedleft\arraybackslash}p{21mm} X}
\toprule
Тип & $I_{\text{нас}}$, А & $\Delta I_{\text{нас}}$, \% &
$P_{\text{нас}}$, Вт & $\Delta P_{\text{нас}}$, \% &
$B_{\text{край,max}}$, Тл & Изменение поля, \% \\
\midrule
I      & $5{,}99$ & $0{,}00$  & $478{,}37$ & $0{,}00$  & $0{,}428$ & $0{,}00$ \\
II-FE  & $5{,}85$ & $-2{,}34$ & $468{,}86$ & $-1{,}99$ & $0{,}387$ & $-9{,}58$ \\
II-SE  & $6{,}51$ & $+8{,}68$ & $519{,}76$ & $+8{,}65$ & $0{,}359$ & $-16{,}12$ \\
III-SE & $6{,}40$ & $+6{,}84$ & $510{,}00$ & $+6{,}61$ & $0{,}383$ & $-10{,}51$ \\
III-FE & $5{,}47$ & $-8{,}68$ & $439{,}25$ & $-8{,}18$ & $0{,}438$ & $+2{,}34$ \\
\bottomrule
\end{tabularx}
\end{table}

\source{Токи и мощности --- Liao и Chen~\cite[таблица~9]{liao}; максимальные
значения краевого поля --- результаты FEM из рисунков~9--13 той же работы.
Относительные изменения рассчитаны относительно типа I.}

Наибольший ток получен у типа II-SE: прирост составляет $8{,}68\,\%$.
Вариант III-SE даёт $6{,}84\,\%$, но дополнительный зазор не улучшает ток
по сравнению с II-SE. Полностью ферритовые II-FE и III-FE не повышают ток
насыщения. Значение $B_{\text{край,max}}$ у III-FE выше классического на
$2{,}34\,\%$, поэтому увеличение числа зазоров не гарантирует монотонного
уменьшения локального поля.

\subsection{Потери, коэффициент полезного действия и нагрев}

Потери составного магнитопровода рассчитывают раздельно по материалам:
\begin{equation}
 P_{\text{маг}}=
 \int_{V_{\text{Fe}}}
 p_{v,\text{Fe}}(f,B_{\text{AC}},B_{\text{DC}},T)\dd V+
 \sum_{k=1}^{n_s}\int_{V_{s,k}}
 p_{v,s,k}(f,B_{\text{AC}},B_{\text{DC}},T)\dd V,
 \label{eq:hybrid-core-loss}
\end{equation}
\begin{whereblock}
\whereitem{$P_{\text{маг}}$}{суммарные потери феррита и вставок, Вт;}
\whereitem{$V_{\text{Fe}}$}{объём ферритовой части, м$^3$;}
\whereitem{$V_{s,k}$}{объём $k$-й вставки, м$^3$;}
\whereitem{$p_{v,\text{Fe}}$, $p_{v,s,k}$}{локальные объёмные плотности
потерь соответствующих материалов, Вт/м$^3$.}
\end{whereblock}

К потерям по формуле~\eqref{eq:hybrid-core-loss} добавляют гармонические
потери обмотки по формуле~\eqref{eq:copper-loss-harmonic} и потери в
окружающих деталях. Использовать один коэффициент Steinmetz для всего
гибридного объёма не допускается.

Статья не содержит сопоставимой таблицы собственных потерь всех пяти
дросселей. В ней приведён коэффициент полезного действия всего PFC при
$P_o=500\,\text{Вт}$. Поэтому таблица~\ref{tab:hybrid-converter-losses}
показывает не потери одного дросселя, а суммарные потери преобразователя,
восстановленные по формуле
\begin{equation}
 P_{\text{пот,пр}}=P_o\left(\frac{1}{\eta}-1\right),
 \label{eq:converter-loss-from-efficiency}
\end{equation}
\begin{whereblock}
\whereitem{$P_{\text{пот,пр}}$}{суммарные потери преобразователя, Вт;}
\whereitem{$P_o$}{выходная мощность, Вт;}
\whereitem{$\eta$}{измеренный коэффициент полезного действия.}
\end{whereblock}

\begin{table}[H]
\caption{Суммарные потери PFC при 500 Вт: значение, Вт / изменение к типу I, \%}
\label{tab:hybrid-converter-losses}
\centering
\small
\begin{tabularx}{\textwidth}{>{\raggedright\arraybackslash}p{20mm}
*{4}{>{\centering\arraybackslash}X}}
\toprule
Тип & $90\,\text{В}/47\,\text{Гц}$ & $115\,\text{В}/60\,\text{Гц}$ &
$230\,\text{В}/50\,\text{Гц}$ & $264\,\text{В}/63\,\text{Гц}$ \\
\midrule
I      & $48{,}366 / 0{,}00$   & $33{,}846 / 0{,}00$   & $17{,}224 / 0{,}00$   & $15{,}730 / 0{,}00$ \\
II-FE  & $47{,}465 / -1{,}86$  & $32{,}368 / -4{,}37$  & $16{,}636 / -3{,}41$  & $15{,}039 / -4{,}39$ \\
II-SE  & $47{,}645 / -1{,}49$  & $33{,}276 / -1{,}68$  & $17{,}438 / +1{,}24$  & $16{,}156 / +2{,}71$ \\
III-SE & $48{,}968 / +1{,}24$  & $33{,}220 / -1{,}85$  & $16{,}529 / -4{,}03$  & $15{,}305 / -2{,}70$ \\
III-FE & $48{,}065 / -0{,}62$  & $33{,}447 / -1{,}18$  & $17{,}170 / -0{,}31$  & $15{,}783 / +0{,}34$ \\
\bottomrule
\end{tabularx}
\end{table}

\source{Коэффициенты полезного действия --- Liao и
Chen~\cite[таблица~11]{liao}; потери и относительные изменения рассчитаны по
формуле~\eqref{eq:converter-loss-from-efficiency}. Отрицательное изменение
означает уменьшение потерь относительно типа I.}

Тип II-FE имеет наименьшие суммарные потери в трёх из четырёх режимов, но
его ток насыщения на $2{,}34\,\%$ ниже классического и он не обеспечивает
требуемые $500\,\text{Вт}$ в испытании статьи. Тип II-SE имеет наибольший ток,
однако при $230$ и $264\,\text{В}$ суммарные потери преобразователя выше
классического на $1{,}24$ и $2{,}71\,\%$ соответственно. Тип III-SE даёт
прирост тока $6{,}84\,\%$ и меньшие суммарные потери в трёх режимах, но при
$90\,\text{В}$ потери увеличиваются на $1{,}24\,\%$.

Результаты теплового испытания после $20\,\text{мин}$ при
$115\,\text{В}$, $60\,\text{Гц}$, $500\,\text{Вт}$ и температуре окружающей
среды $26\,^{\circ}\text{С}$ приведены в
таблице~\ref{tab:hybrid-temperature-results}.

\begin{table}[H]
\caption{Температуры относительно классического типа I}
\label{tab:hybrid-temperature-results}
\centering
\small
\begin{tabularx}{\textwidth}{>{\raggedright\arraybackslash}p{24mm}
*{4}{>{\centering\arraybackslash}X}}
\toprule
Тип & $T_{\text{серд}}$, $^{\circ}\text{С}$ &
$\Delta T_{\text{серд}}$, К & $T_{\text{обм}}$, $^{\circ}\text{С}$ &
$\Delta T_{\text{обм}}$, К \\
\midrule
I      & $64{,}8$ & $0{,}0$  & $67{,}5$ & $0{,}0$ \\
II-FE  & $60{,}6$ & $-4{,}2$ & $67{,}6$ & $+0{,}1$ \\
II-SE  & $60{,}2$ & $-4{,}6$ & $64{,}1$ & $-3{,}4$ \\
III-SE & $60{,}3$ & $-4{,}5$ & $62{,}4$ & $-5{,}1$ \\
III-FE & $58{,}1$ & $-6{,}7$ & $63{,}0$ & $-4{,}5$ \\
\bottomrule
\end{tabularx}
\end{table}

\source{Liao и Chen~\cite[таблица~13]{liao}. Разность температур рассчитана
относительно типа I. Температура после $20\,\text{мин}$ не доказывает
достижение установившегося теплового режима.}

\begin{table}[H]
\caption{Краткий выбор по результатам статьи}
\label{tab:hybrid-short-choice}
\centering
\small
\begin{tabularx}{\textwidth}{>{\raggedright\arraybackslash}p{43mm}
>{\raggedright\arraybackslash}p{28mm} X}
\toprule
Критерий & Тип & Численный результат относительно типа I \\
\midrule
Максимальный ток насыщения & II-SE & $+8{,}68\,\%$ \\
Максимальная испытанная мощность до критерия насыщения & II-SE &
$+8{,}65\,\%$ \\
Минимум краевого поля & II-SE & $-16{,}12\,\%$ \\
Минимум суммарных потерь PFC в большинстве режимов & II-FE &
$-1{,}86\ldots-4{,}39\,\%$, но ток насыщения $-2{,}34\,\%$ \\
Компромисс «ток и потери» при средней и высокой входной сети & III-SE &
ток $+6{,}84\,\%$; потери от $-1{,}85$ до $-4{,}03\,\%$ \\
\bottomrule
\end{tabularx}
\end{table}

\recommend Для режима с максимальным током при низком входном напряжении
первым гибридным кандидатом является II-SE. Если приоритетом являются потери
на средней и высокой входной сети при сохранении запаса по току, дополнительно
проверяют III-SE. Численные преимущества из
таблиц~\ref{tab:hybrid-saturation-results}--
\ref{tab:hybrid-short-choice} относятся только к испытанному PQ3530 при
$65\,\text{кГц}$ и не являются масштабным коэффициентом для другого
магнитопровода или частоты.

\section{ПРОЕКТИРОВАНИЕ ОБМОТКИ}
\label{sec:winding}

\subsection{Требуемое сечение и размещение в окне}

\stepgoal{До создания FEM-модели получить полное обозначение проводника,
число и диаметр жил, наружный размер, число слоёв, длину каждого витка и
чертёж укладки. Без этих данных поле и потери обмотки не определены.}

\begin{figure}[htbp]
\centering
\includegraphics[width=0.90\textwidth]{solid_and_litz_prototypes.png}
\caption{Экспериментальные дроссели с литцендратом и одножильным круглым
проводом на одинаковых магнитопроводах (по Ewald и
Biela~\cite[рисунок~1]{ewald_litz})}
\label{fig:solid-litz-prototypes}
\end{figure}

Рисунок~\ref{fig:solid-litz-prototypes} показывает, что тип проводника
изменяет не только высокочастотное сопротивление, но и фактическую геометрию
обмотки: наружный размер, положение витков относительно зазора, длину витка
и условия охлаждения. Поэтому сравнивают законченные варианты укладки, а не
одинаковую площадь голой меди.

\begin{figure}[htbp]
\centering
\includegraphics[
  page=3,
  trim=55bp 565bp 55bp 45bp,
  clip,
  width=\textwidth
]{source_ewald_litz.pdf}
\caption{Геометрические параметры магнитопровода и обмотки, входящие в
двумерную модель потерь проводника (по Ewald и
Biela~\cite[рисунок~2]{ewald_litz})}
\label{fig:winding-core-geometry}
\end{figure}

Геометрию на рисунке~\ref{fig:winding-core-geometry} используют как
контрольный перечень исходных размеров: ширины и высоты окна, сечения
стержней, длины зазора, размеров проводника и расстояния меди до кромки
зазора. Если хотя бы один из этих размеров ещё не выбран, расчёт потерь
обмотки в полевой модели является предварительным.

\begin{table}[H]
\caption{Свойства электротехнической меди, используемые в расчёте}
\label{tab:copper-properties}
\centering
\small
\begin{tabularx}{\textwidth}{>{\raggedright\arraybackslash}X
>{\centering\arraybackslash}p{31mm}
>{\raggedleft\arraybackslash}p{48mm}}
\toprule
Величина & Обозначение & Значение при $20\,^{\circ}\text{С}$ \\
\midrule
Удельное электрическое сопротивление &
$\rho_{\text{Cu},20}$ &
$1{,}724\cdot10^{-8}\,\text{Ом}\cdot\text{м}$ \\
Температурный коэффициент сопротивления &
$\alpha_{\text{Cu}}$ &
$3{,}93\cdot10^{-3}\,\text{К}^{-1}$ \\
Относительная магнитная проницаемость меди &
$\mu_{r,\text{Cu}}$ &
$0{,}999\,991$ \\
Абсолютная магнитная проницаемость меди &
$\mu_{\text{Cu}}=\mu_0\mu_{r,\text{Cu}}$ &
$1{,}256\,625\,75\cdot10^{-6}\,\text{Гн/м}$ \\
\bottomrule
\end{tabularx}
\end{table}

\source Удельное сопротивление и температурный коэффициент даны для
стандартной отожжённой меди по IEC~60028~\cite{iec60028}. Относительная
магнитная проницаемость приведена по таблице~2 руководства
Paulter~\cite[таблица~2]{nist_copper}. Для фактической марки провода
значения уточняют по паспорту: отжиг, сплав и покрытие изменяют свойства.

Удельное электрическое сопротивление меди при температуре $\vartheta$
определяют по формуле
\begin{equation}
 \rho_{\text{Cu}}(\vartheta)=
 \rho_{\text{Cu},20}
 \left[1+\alpha_{\text{Cu}}(\vartheta-20\,^{\circ}\text{С})\right],
 \label{eq:rho-temperature}
\end{equation}
\begin{whereblock}
\whereitem{$\rho_{\text{Cu}}(\vartheta)$}{удельное сопротивление меди при
температуре $\vartheta$, Ом$\cdot$м;}
\whereitem{$\rho_{\text{Cu},20}$}{удельное сопротивление меди при
$20\,^{\circ}\text{С}$, Ом$\cdot$м;}
\whereitem{$\alpha_{\text{Cu}}$}{температурный коэффициент сопротивления
меди, К$^{-1}$;}
\whereitem{$\vartheta$}{температура меди, $^{\circ}\text{С}$.}
\end{whereblock}

Предварительное сечение уже определено по
формуле~\eqref{eq:copper-area}, а возможность размещения --- по
формулам~\eqref{eq:window-copper-fill} и
\eqref{eq:window-occupancy}. На текущем шаге предварительный проводник
заменяют конкретным каталожным исполнением и разрабатывают послойный
чертёж укладки.

\recommend Значения $J_{\text{доп}}$ и $k_{\text{Cu}}$ из
таблицы~\ref{tab:kw-j} используют только
для первого варианта. Окончательное сечение выбирают по вычисленным
потерям и температуре. Наружную геометрию проверяют отдельно по
$k_{\text{зан}}$ и по чертежу; низкий $k_{\text{Cu}}$ не гарантирует,
что толстый литцендрат физически размещается в окне.

\subsection{Глубина проникновения тока}

Глубину проникновения тока в медь определяют по формуле
\begin{equation}
 \delta_{\text{Cu}}(f,\vartheta)=
 \sqrt{\frac{\rho_{\text{Cu}}(\vartheta)}
 {\pi f\mu_{\text{Cu}}}},
 \label{eq:skin-depth}
\end{equation}
\begin{whereblock}
\whereitem{$\delta_{\text{Cu}}$}{глубина проникновения тока, м;}
\whereitem{$f$}{частота рассматриваемой гармоники, Гц;}
\whereitem{$\mu_{\text{Cu}}$}{абсолютная магнитная проницаемость меди по
таблице~\ref{tab:copper-properties}, Гн/м.}
\end{whereblock}

Для предварительного выбора жилы при несинусоидальном токе допускается
использовать эффективную частоту Sullivan и Zhang~\cite[приложение~A]{sullivan}
\begin{equation}
 f_{\text{эф}}=
 \frac{
 \sqrt{\dfrac{1}{T_{\text{с}}}
 \displaystyle\int_0^{T_{\text{с}}}
 \left(\dfrac{di}{dt}\right)^2\dd t}}
 {2\pi I_{\text{действ}}},
 \label{eq:effective-frequency}
\end{equation}
\begin{whereblock}
\whereitem{$f_{\text{эф}}$}{эффективная частота для первого выбора
литцендрата, Гц;}
\whereitem{$i(t)$}{полный мгновенный ток, включая постоянную составляющую,
А;}
\whereitem{$T_{\text{с}}$}{период формы тока, с.}
\end{whereblock}

\recommend Для синусоидального тока предварительно принимают
$d_s\leq\delta_{\text{Cu}}$; граница $d_s\leq2\delta_{\text{Cu}}$ допустима
только как грубый фильтр. Для импульсного тока глубину проникновения для
первого кандидата вычисляют при $f_{\text{эф}}$, но окончательные потери
суммируют по гармоникам. Ни частота коммутации, ни условие
$d_s<2\delta_{\text{Cu}}$ сами по себе не учитывают эффект близости и поле
зазора.

\subsection{Литцендрат и круглый одножильный провод}

\method Провод выбирают в следующем порядке: сначала задают несколько
стандартных диаметров жилы, затем для каждого диаметра рассчитывают число жил
по сечению и по высокочастотным потерям, после чего проверяют наружный размер
кабеля и укладку. Выбор только по таблице «диаметр~-- частота» не является
расчётом литцендрата.

\begin{figure}[htbp]
\centering
\includegraphics[
  page=1,
  trim=335bp 43bp 65bp 540bp,
  clip,
  width=0.58\textwidth
]{source_sullivan_litz.pdf}
\caption{Последовательность выбора диаметра и числа жил литцендрата
(по Sullivan и Zhang~\cite[рисунок~1]{sullivan})}
\label{fig:litz-selection-flowchart}
\end{figure}

Схему на рисунке~\ref{fig:litz-selection-flowchart} применяют как цикл:
назначают доступный диаметр жилы, вычисляют требуемое число жил, проверяют
потери и окно, затем повторяют расчёт для соседних стандартных диаметров.
Результатом является не одно число, а сравнение нескольких реально
изготавливаемых кабелей.

Минимальное число жил по активному сечению определяют как
\begin{equation}
 n_{\text{сеч}}=
 \left\lceil
 \frac{4S_{\text{Cu,треб}}}{\pi d_s^2}
 \right\rceil,
 \label{eq:litz-count-area}
\end{equation}
\begin{whereblock}
\whereitem{$n_{\text{сеч}}$}{минимальное число жил по условию активного
сечения;}
\whereitem{$S_{\text{Cu,треб}}$}{требуемая суммарная площадь меди,
м$^2$;}
\whereitem{$d_s$}{диаметр меди одной жилы без изоляции, м.}
\end{whereblock}

Для каждого диаметра жилы экономичное число жил предварительно определяют по
формуле Sullivan и Zhang~\cite[формула~(2)]{sullivan}
\begin{equation}
 n_e=k_s\frac{\delta_s^2b_w}{N_s},
 \label{eq:litz-economic-count}
\end{equation}
\begin{whereblock}
\whereitem{$n_e$}{экономичное число элементарных жил;}
\whereitem{$k_s$}{коэффициент выбранного диаметра жилы, мм$^{-3}$;}
\whereitem{$\delta_s$}{глубина проникновения тока, мм;}
\whereitem{$b_w$}{ширина секции обмотки поперёк поверхности, через которую
секция обменивается магнитным полем с соседней секцией, мм;}
\whereitem{$N_s$}{число витков от поверхности нулевого поля до границы
рассматриваемой секции.}
\end{whereblock}

\begin{figure}[htbp]
\centering
\includegraphics[
  page=2,
  trim=50bp 65bp 306bp 490bp,
  clip,
  width=0.72\textwidth
]{source_sullivan_litz.pdf}
\caption{Определение ширины секции $b_w$ и числа витков $N_s$ для расчёта
эффекта близости (по Sullivan и Zhang~\cite[рисунок~2]{sullivan})}
\label{fig:litz-bw-ns}
\end{figure}

Для простой неразделённой обмотки принимают $N_s=N$. Для чередующихся секций
$N_s$ считают от плоскости нулевого поля согласно фактической укладке. Для
дросселя с дискретным зазором, около которого поле существенно неоднородно,
$b_w$ заменяют эффективной шириной из источника либо получают потери из
двумерной или трёхмерной модели; геометрическую ширину окна без такого
уточнения подставлять не следует.

\begin{figure}[htbp]
\centering
\includegraphics[
  page=3,
  trim=115bp 79bp 378bp 598bp,
  clip,
  width=0.48\textwidth
]{source_sullivan_litz.pdf}
\caption{Расположение обмотки относительно дискретного зазора
(по Sullivan и Zhang~\cite[рисунок~3]{sullivan})}
\label{fig:winding-gap-distance}
\end{figure}

\recommend Расстояние от ближайшей меди до кромки зазора увеличивают до
тех пор, пока снижение добавочных потерь подтверждается расчётом. Чрезмерное
удаление обмотки увеличивает среднюю длину витка, сопротивление, поле
рассеяния и паразитную ёмкость, поэтому расстояние выбирают совместно с
числом слоёв и положением выводов.

\begin{table}[H]
\caption{Коэффициенты экономичного выбора литцендрата}
\label{tab:litz-k}
\centering
\small
\begin{tabular}{>{\raggedleft\arraybackslash}p{29mm}
>{\raggedleft\arraybackslash}p{24mm}
>{\raggedleft\arraybackslash}p{30mm}
>{\raggedleft\arraybackslash}p{31mm}}
\toprule
$d_s$, мм & Номер AWG & Экономичный $F_R$ & $k_s$, мм$^{-3}$ \\
\midrule
$0{,}202$ & 32 & $1{,}06$ & 130 \\
$0{,}127$ & 36 & $1{,}13$ & 771 \\
$0{,}080$ & 40 & $1{,}25$ & 4400 \\
$0{,}050$ & 44 & $1{,}47$ & 24000 \\
$0{,}040$ & 46 & $1{,}60$ & 54000 \\
$0{,}032$ & 48 & $1{,}68$ & 115000 \\
\bottomrule
\end{tabular}
\end{table}

Коэффициент увеличения сопротивления при выбранном числе жил оценивают по
формуле~\cite[формула~(3)]{sullivan}
\begin{equation}
 F_R=\frac{R_{\text{AC}}}{R_{\text{DC}}}
 =1+\frac{(\pi n_sN_s)^2d_s^6}
 {192\delta_s^4b_w^2},
 \label{eq:litz-fr}
\end{equation}
\begin{whereblock}
\whereitem{$F_R$}{коэффициент увеличения сопротивления;}
\whereitem{$R_{\text{AC}}$}{сопротивление переменному току на расчётной
частоте, Ом;}
\whereitem{$R_{\text{DC}}$}{сопротивление постоянному току, Ом;}
\whereitem{$n_s$}{фактическое число жил.}
\end{whereblock}

Во всех длинах формул~\eqref{eq:litz-economic-count} и \eqref{eq:litz-fr}
используют миллиметры. Обязательное условие по току имеет вид
$n_s\geq n_{\text{сеч}}$. Значение $n_e$ является центром экономичной области,
а не нижним ограничением. Если вариант около $n_e$ размещается в окне,
проверяют также $0{,}75n_e$ и $1{,}25n_e$. Если $n_e$ не размещается,
определяют наибольшее допустимое по окну $n_s$ и сравнивают его с вариантами
другого $d_s$. Уменьшение относительно $n_e$ более чем на $25\,\%$ не
запрещено, но означает, что выбранный диаметр жилы, вероятно, не даёт
наилучшего соотношения потерь и стоимости.

Для предварительной проверки наружного размера эквивалентный диаметр пучка
определяют по формуле
\begin{equation}
 d_{\text{пуч}}=d_s\sqrt{\frac{n_s}{k_{\text{уп}}}},
 \label{eq:litz-bundle-diameter}
\end{equation}
\begin{whereblock}
\whereitem{$d_{\text{пуч}}$}{эквивалентный наружный диаметр пучка до общей
оплётки, м;}
\whereitem{$k_{\text{уп}}$}{отношение площади голой меди к площади пучка с
межжильной изоляцией.}
\end{whereblock}

\recommend До получения паспорта кабеля принимают $k_{\text{уп}}=0{,}55$.
Значение $0{,}60\ldots0{,}70$ применяют только по данным изготовителя.
Окончательный наружный диаметр $d_{\text{нар}}$, шаги скрутки, тип общей
изоляции и минимальный радиус изгиба берут из паспорта, а не из
формулы~\eqref{eq:litz-bundle-diameter}.

\begin{figure}[htbp]
\centering
\includegraphics[
  page=3,
  trim=45bp 35bp 305bp 585bp,
  clip,
  width=0.65\textwidth
]{source_ewald_litz.pdf}
\caption{Поперечное сечение литцендрата и одножильного круглого провода,
использованное при сравнении потерь (по Ewald и
Biela~\cite[рисунок~3]{ewald_litz})}
\label{fig:litz-solid-cross-section}
\end{figure}

Рисунок~\ref{fig:litz-solid-cross-section} поясняет различие между
диаметром отдельной медной жилы $d_s$, размером всего пучка и наружным
размером кабеля с изоляцией. В расчёт окна подставляют наружный размер, а в
расчёт активного сопротивления --- суммарную площадь меди.

Максимальное число одиночных жил в первой операции скрутки оценивают
по рекомендации Sullivan и Zhang
\begin{equation}
 n_{1,\max}=4\left(\frac{\delta_s}{d_s}\right)^2,
 \label{eq:litz-first-bunch}
\end{equation}
\begin{whereblock}
\whereitem{$n_{1,\max}$}{максимальное число жил в элементарном пучке.}
\end{whereblock}

Если $n_s>n_{1,\max}$, задают многоступенчатую конструкцию, например
$4/5/5$, и проверяют, что произведение чисел ступеней равно $n_s$.
Обычный пучок параллельных проводов без транспозиции не считают
литцендратом: ток между жилами может распределяться неравномерно.

Кандидатное число жил выбирают не меньше
$n_{\text{сеч}}$ и сопоставляют с экономическим числом из
формулы~\eqref{eq:litz-economic-count}. По рекомендации
Sullivan и Zhang~\cite{sullivan} исследуют также варианты примерно на
$25\,\%$ выше и ниже экономического числа, если они удовлетворяют
условиям по сечению и окну.

Для одножильного круглого провода требуемый диаметр меди равен
\begin{equation}
 d_{\text{одн,треб}}=
 \sqrt{\frac{4S_{\text{Cu,треб}}}{\pi}},
 \label{eq:solid-wire-diameter}
\end{equation}
\begin{whereblock}
\whereitem{$d_{\text{одн,треб}}$}{требуемый диаметр голой медной жилы, м.}
\end{whereblock}

Из стандартного ряда выбирают ближайший больший диаметр и проверяют его
$R_{\text{AC}}/R_{\text{DC}}$. Если
$d_{\text{одн,треб}}>\delta_{\text{Cu}}$, одножильный провод не отклоняют
автоматически, но обязательно сравнивают с литцендратом, фольгой или
параллельными транспонированными жилами по полным потерям.

Число витков в одном слое круглого кабеля предварительно определяют по
формуле
\begin{equation}
 N_{\text{сл}}=
 \left\lfloor
 \frac{b_{\text{нам}}}
 {d_{\text{нар}}+s_{\text{вит}}}
 \right\rfloor,
 \label{eq:turns-per-layer}
\end{equation}
\begin{whereblock}
\whereitem{$N_{\text{сл}}$}{число витков, размещаемых в одном слое;}
\whereitem{$b_{\text{нам}}$}{доступная ширина намотки после вычета полей и
изоляции, м;}
\whereitem{$d_{\text{нар}}$}{наибольший паспортный наружный диаметр кабеля,
м;}
\whereitem{$s_{\text{вит}}$}{технологический промежуток между соседними
витками, м.}
\end{whereblock}

Число слоёв равно
\begin{equation}
 m_{\text{сл}}=
 \left\lceil\frac{N}{N_{\text{сл}}}\right\rceil,
 \label{eq:winding-layer-count}
\end{equation}
\begin{whereblock}
\whereitem{$m_{\text{сл}}$}{число слоёв обмотки.}
\end{whereblock}

\recommend При ручной намотке без согласованной технологии принимают
$s_{\text{вит}}\geq0{,}1d_{\text{нар}}$ и резервируют не менее
$5\,\%$ ширины намотки с каждой стороны. Эти значения являются
технологическим запасом первого чертежа; окончательные размеры задаёт
изготовитель. Межслойную изоляцию и радиус изгиба добавляют к радиальной
толщине и проверяют на чертеже.

Сопротивление обмотки постоянному току при температуре $\vartheta$
определяют по формуле
\begin{equation}
 R_{\text{DC}}(\vartheta)=
 \rho_{\text{Cu}}(\vartheta)
 \frac{l_{\text{Cu}}}{S_{\text{Cu}}},
 \label{eq:rdc}
\end{equation}
\begin{whereblock}
\whereitem{$R_{\text{DC}}$}{сопротивление постоянному току, Ом;}
\whereitem{$l_{\text{Cu}}$}{полная длина меди, м;}
\whereitem{$S_{\text{Cu}}$}{суммарная активная площадь меди, м$^2$.}
\end{whereblock}

\recommend Длину $l_{\text{Cu}}$ определяют послойно по чертежу.
Оценку $N l_{\text{ср}}$ допускается использовать только до разработки
укладки, где $l_{\text{ср}}$ --- средняя длина витка, м. Наружный диаметр,
кратность шага скрутки, класс термостойкости и коэффициент упаковки
литцендрата берут из паспорта изготовителя кабеля.

\begin{table}[H]
\caption{Обязательный результат выбора проволочной обмотки}
\label{tab:winding-selection-result}
\centering
\small
\begin{tabularx}{\textwidth}{>{\raggedright\arraybackslash}p{53mm} X}
\toprule
Параметр & Что должно быть записано в расчёте \\
\midrule
Жила & материал, $d_s$, покрытие, класс нагревостойкости \\
Литцендрат & $n_s$, структура скрутки, шаги скрутки, общая изоляция,
полное обозначение изготовителя \\
Активная и наружная геометрия & $S_{\text{Cu}}$, $d_{\text{нар}}$ с
верхним допуском, минимальный радиус изгиба \\
Укладка & $N_{\text{сл}}$, $m_{\text{сл}}$, межслойная изоляция, расстояние
до каждого зазора, чертёж поперечного сечения окна \\
Проверки & $k_{\text{Cu}}$, $k_{\text{зан}}$, $l_{\text{Cu}}$,
$R_{\text{DC}}(\vartheta)$, $F_R$ или $R_{\text{AC}}(f_h)$,
$P_{\text{Cu}}$, максимальная температура \\
\bottomrule
\end{tabularx}
\end{table}

\subsection{Фольговая и прямоугольная обмотка}

Для $n_{\parallel}$ параллельных фольг требуемую ширину одной фольги
определяют по формуле
\begin{equation}
 b_{\text{ф,треб}}=
 \frac{S_{\text{Cu,треб}}}
 {n_{\parallel}t_{\text{ф}}},
 \label{eq:foil-width}
\end{equation}
\begin{whereblock}
\whereitem{$b_{\text{ф,треб}}$}{требуемая ширина одной фольги, м;}
\whereitem{$n_{\parallel}$}{число электрически параллельных фольг;}
\whereitem{$t_{\text{ф}}$}{толщина меди фольги, м.}
\end{whereblock}

\recommend Первым кандидатом принимают
$t_{\text{ф}}\leq\delta_{\text{Cu}}$ на основной частоте потерь. Если
полученная ширина не размещается в окне, не увеличивают толщину без проверки:
рассматривают несколько параллельных фольг, большее окно или другую
геометрию. У зазора широкая фольга экранирует переменное поле вихревыми
токами; это меняет и $R_{\text{AC}}$, и индуктивность. Поэтому окончательные
потери вычисляют частотно-зависимой двумерной или трёхмерной моделью с
фактическими расстояниями до зазора~\cite{ewaldfoil}.

Для каждого слоя фиксируют направление тока и потенциал. Параллельные фольги
должны иметь одинаковое потокосцепление или транспозицию; простое
параллельное соединение слоёв с разным положением в окне может дать
неравномерное распределение тока.

\subsection{Планарная обмотка}

Сопротивление медных дорожек постоянному току вычисляют по формуле
\begin{equation}
 R_{\text{DC,пл}}=
 \rho_{\text{Cu}}(\vartheta)
 \sum_{k=1}^{p}
 \frac{l_k}{n_{\parallel,k}b_k t_k}
 {}+R_{\text{пер}},
 \label{eq:planar-rdc}
\end{equation}
где $R_{\text{DC,пл}}$ --- сопротивление планарной обмотки, Ом;\par\noindent
$p$ --- число участков дорожек;\par\noindent
$l_k$ --- длина $k$-го участка, м;\par\noindent
$n_{\parallel,k}$ --- число параллельных дорожек на $k$-м участке;\par\noindent
$b_k$ --- ширина дорожки, м;\par\noindent
$t_k$ --- конечная толщина меди после изготовления, м;\par\noindent
$R_{\text{пер}}$ --- суммарное сопротивление межслойных переходов, Ом.

Сопротивление цилиндрических металлизированных отверстий предварительно
определяют по формуле
\begin{equation}
 R_{\text{пер}}=
 \rho_{\text{Cu}}(\vartheta)
 \sum_{r=1}^{s}
 \frac{l_{\text{пер},r}}
 {n_{\text{пер},r}\,
 \pi\!\left(d_{\text{нар},r}^{\,2}-d_{\text{вн},r}^{\,2}\right)/4},
 \label{eq:via-resistance}
\end{equation}
где $s$ --- число групп переходов;\par\noindent
$l_{\text{пер},r}$ --- длина металлизированного цилиндра $r$-й группы,
м;\par\noindent
$n_{\text{пер},r}$ --- число параллельных переходов в группе;\par\noindent
$d_{\text{нар},r}$, $d_{\text{вн},r}$ --- наружный и внутренний диаметры
металлизации, м.

\recommend Для планарной обмотки толщину меди, минимальные ширины и
зазоры, допустимый ток перехода и остаточную толщину диэлектрика берут
из технологических норм изготовителя печатной платы. При
$t_k\gtrsim\delta_{\text{Cu}}$ и при двух и более слоях аналитический
расчёт дополняют двумерной или трёхмерной электромагнитной моделью,
учитывающей зазор и фактическое чередование токовых слоёв.

\source Частотно-зависимые модели потерь фольговых и планарных обмоток
приведены в \cite{schaefer,ewaldfoil}.

\section{РАСЧЁТ ПОТЕРЬ}
\label{sec:losses}

\subsection{Связь потерь дросселя с коэффициентом полезного действия}

Коэффициент полезного действия (КПД) относится к преобразователю или
законченному устройству, а не к идеализированному дросселю, который
периодически принимает и возвращает реактивную энергию. КПД преобразователя с
учётом потерь дросселя определяют по формуле
\begin{equation}
 \eta=
 \frac{P_{\text{вых}}}
 {P_{\text{вых}}+P_{\text{ост}}+P_{\Sigma}},
 \label{eq:converter-efficiency}
\end{equation}
\begin{whereblock}
\whereitem{$\eta$}{коэффициент полезного действия преобразователя;}
\whereitem{$P_{\text{вых}}$}{выходная активная мощность, Вт;}
\whereitem{$P_{\text{ост}}$}{сумма потерь всех элементов, кроме дросселя,
Вт;}
\whereitem{$P_{\Sigma}$}{суммарные потери дросселя, Вт.}
\end{whereblock}

Для профиля из $R$ режимов средневзвешенные потери дросселя равны
\begin{equation}
 \overline P_{\Sigma}=
 \sum_{r=1}^{R}w_rP_{\Sigma,r},
 \qquad
 \sum_{r=1}^{R}w_r=1,
 \label{eq:mission-profile-loss}
\end{equation}
\begin{whereblock}
\whereitem{$\overline P_{\Sigma}$}{средневзвешенные потери дросселя, Вт;}
\whereitem{$w_r$}{доля времени или вероятность $r$-го режима;}
\whereitem{$P_{\Sigma,r}$}{потери дросселя в $r$-м режиме, Вт.}
\end{whereblock}

\criterion Для максимального КПД сначала ограничивают потери и температуру в
наихудшем режиме, затем минимизируют $\overline P_{\Sigma}$ по фактическому
профилю. Оптимум не обязан удовлетворять
$P_{\text{Cu}}=P_{\text{Fe}}$: это равенство допустимо как начальное
приближение, но окончательное решение получают суммированием всех потерь
\cite{papamanolis,pichon}.

\subsection{Потери в обмотке}

\textbf{Цель шага} --- разделить потери от постоянной составляющей и
потери каждой гармоники с учётом скин-эффекта и эффекта близости.

Полные потери в обмотке определяют по формуле
\begin{equation}
 P_{\text{Cu}}=
 I_{\text{ср}}^{\,2}R_{\text{DC}}(\vartheta)
 +\sum_{h=1}^{H}
 I_{h,\text{действ}}^{\,2}
 R_{\text{AC}}(f_h,\vartheta),
 \label{eq:copper-loss-harmonic}
\end{equation}
где $P_{\text{Cu}}$ --- потери в обмотке, Вт;\par\noindent
$H$ --- число учитываемых гармоник;\par\noindent
$I_{h,\text{действ}}$ --- действующее значение $h$-й гармоники тока, А;\par\noindent
$R_{\text{AC}}(f_h,\vartheta)$ --- сопротивление переменному току на
частоте $f_h$ при температуре $\vartheta$, Ом;\par\noindent
$f_h$ --- частота $h$-й гармоники, Гц.

\recommend Спектр обрезают только после проверки сходимости суммы:
добавление следующей группы гармоник должно изменять $P_{\text{Cu}}$
менее чем на принятый допуск. Для проводника около зазора
$R_{\text{AC}}$ определяют с учётом поля выпучивания; коэффициент
$R_{\text{AC}}/R_{\text{DC}}$ для прямого однородного поля использовать
недостаточно.

\method Значения $I_{h,\text{действ}}$ получают из ряда Фурье рассчитанной
или измеренной формы $i(t)$. Значение $R_{\text{AC}}(f_h,\vartheta)$ получают:
для предварительного литцендрата --- по формуле~\eqref{eq:litz-fr}; для
короткого списка --- по аналитической модели поля зазора
\cite{ewald_litz,meng}; для окончательной геометрии --- частотной FEM-моделью
или четырёхпроводным измерением. Сопротивление выводов и разделки
литцендрата измеряют отдельно и относят к $P_{\text{конт}}$.

\subsection{Потери в магнитопроводе}

Для несинусоидальной периодической формы индукции удельные потери допускается
рассчитывать по улучшенному обобщённому уравнению Штейнмеца
(improved Generalized Steinmetz Equation, iGSE):
\begin{equation}
 p_v=
 \frac{1}{T_{\text{с}}}
 \int_{0}^{T_{\text{с}}}
 k_i(T,B_{\text{DC}})
 \left|\frac{dB}{dt}\right|^{\alpha}
 \left(\Delta B_{\text{пп}}\right)^{\beta-\alpha}
 dt,
 \label{eq:igse}
\end{equation}
\begin{whereblock}
\whereitem{$p_v$}{средние удельные объёмные потери, Вт/м$^3$;}
\whereitem{$k_i$}{коэффициент улучшенного обобщённого уравнения Штейнмеца
(improved Generalized Steinmetz Equation, iGSE);}
\whereitem{$B(t)$}{мгновенная магнитная индукция, восстановленная из
напряжения, Тл;}
\whereitem{$\alpha$, $\beta$}{коэффициенты аппроксимации потерь материала;}
\whereitem{$\Delta B_{\text{пп}}$}{размах переменной магнитной индукции,
Тл.}
\end{whereblock}

\begin{figure}[htbp]
\centering
\includegraphics[
  page=6,
  trim=346bp 418bp 61bp 50bp,
  clip,
  width=0.62\textwidth
]{source_muehlethaler_core_loss.pdf}
\caption{Треугольная форма магнитной индукции и зависимость потерь
магнитопровода от коэффициента заполнения импульса
(по Mühlethaler с соавторами~\cite[рисунки~7, 8]{muehlethaler})}
\label{fig:core-loss-duty-cycle}
\end{figure}

График на рисунке~\ref{fig:core-loss-duty-cycle} показывает, что одинаковых
$f$ и $\Delta B_{\text{пп}}$ недостаточно для определения потерь:
изменение длительности участков $dB/dt$ меняет среднюю мощность потерь.
Поэтому в формулу~\eqref{eq:igse} подставляют восстановленную рабочую форму
$B(t)$, а не заменяющую её синусоиду.

\criterion Формулу~\eqref{eq:igse} применяют, только если релаксационные и
малые петли гистерезиса несущественны либо их влияние уже учтено
экспериментальной аппроксимацией. При наличии пауз с $dB/dt=0$ после
изменения потока обычная iGSE может занижать потери. В таком режиме используют
измеренную карту потерь или модифицированную i$^2$GSE с коэффициентами,
полученными для выбранного материала; без этих данных результат не считают
приёмочным.

Если паспортные потери аппроксимированы обычным уравнением Штейнмеца
$p_v=kf^\alpha\widehat B^\beta$, коэффициент $k_i$ определяют по
формуле~\cite[формула~(3)]{muehlethaler}
\begin{equation}
 k_i=
 \frac{k}
 {(2\pi)^{\alpha-1}
 2^{\beta-\alpha}
 \displaystyle\int_0^{2\pi}|\cos\theta|^\alpha\dd\theta},
 \label{eq:igse-ki}
\end{equation}
\begin{whereblock}
\whereitem{$k$}{коэффициент обычного уравнения Штейнмеца в согласованных
единицах;}
\whereitem{$\theta$}{безразмерная переменная интегрирования.}
\end{whereblock}

\method Параметры $k$, $\alpha$ и $\beta$ не берут из другой марки феррита.
Их получают из таблицы изготовителя или локальной аппроксимацией нескольких
паспортных точек в диапазоне расчётных $f$, $\widehat B$ и $T$. При
постоянном подмагничивании используют отдельную измеренную карту
$p_v(f,\Delta B,B_{\text{DC}},T)$ или подтверждённый поправочный множитель.
Модель без данных по $B_{\text{DC}}$ не объявляют точной: влияние
подмагничивания на потери зависит от материала и рабочей точки
\cite{sanusi}.

\begin{figure}[htbp]
\centering
\includegraphics[width=0.78\textwidth]{dc_bias_thermal_comparison.png}
\caption{Тепловое сравнение магнитопровода без постоянного подмагничивания
и при его наличии (по Sanusi с соавторами~\cite[рисунок~1]{sanusi})}
\label{fig:dc-bias-thermal}
\end{figure}

\begin{figure}[htbp]
\centering
\includegraphics[
  page=8,
  trim=47bp 346bp 65bp 58bp,
  clip,
  width=0.92\textwidth
]{source_sanusi_dc_bias.pdf}
\caption{Измеренное влияние постоянного подмагничивания на удельные потери
для разных амплитуд, частот и материалов
(по Sanusi с соавторами~\cite[рисунки~12--15]{sanusi})}
\label{fig:dc-bias-loss-curves}
\end{figure}

Рисунки~\ref{fig:dc-bias-thermal} и
\ref{fig:dc-bias-loss-curves} подтверждают, что поправка на постоянное
подмагничивание не является одной универсальной константой. Численные
значения с графиков относятся только к испытанным материалам и режимам;
для выбранной марки используют её собственную карту потерь либо измерение
на образце.

При неоднородных поле и температуре потери интегрируют по объёму:
\begin{equation}
 P_{\text{Fe}}=
 \int_{V_e}p_v(\mathbf r)\,dV,
 \label{eq:core-loss-volume}
\end{equation}
где $P_{\text{Fe}}$ --- потери в магнитопроводе, Вт;\par\noindent
$\mathbf r$ --- радиус-вектор точки магнитопровода, м;\par\noindent
$dV$ --- элемент объёма, м$^3$.

\source Применимость iGSE, влияние формы напряжения и релаксационных
участков рассмотрены в~\cite{muehlethaler}. Влияние постоянного
подмагничивания рассмотрено в~\cite{sanusi}. Паспортные кривые материала
имеют приоритет в пределах указанных изготовителем условий.

\recommend Если отсутствуют коэффициенты для требуемых температуры,
$B_{\text{DC}}$ и формы возбуждения, значение $P_{\text{Fe}}$ не
экстраполируют как достоверное. Выполняют измерение на образце или
вводят подтверждённую консервативную границу. Для планарного
магнитопровода проверяют локальные потери у зазора, а не только
$p_v V_e$ по средней индукции.

\subsection{Дополнительные потери}

Суммарные потери дросселя определяют по формуле
\begin{equation}
 P_{\Sigma}=P_{\text{Cu}}+P_{\text{Fe}}+
 P_{\text{экр}}+P_{\text{конт}}+P_{\text{проч}},
 \label{eq:total-loss}
\end{equation}
где $P_{\Sigma}$ --- суммарные потери, Вт;\par\noindent
$P_{\text{экр}}$ --- потери в электростатических или магнитных экранах,
Вт;\par\noindent
$P_{\text{конт}}$ --- потери в выводах, соединениях и контактах, Вт;\par\noindent
$P_{\text{проч}}$ --- прочие учтённые потери, Вт.

\recommend Металлические крепёжные детали, теплоотводы, экраны и
полигоны печатной платы, расположенные в поле зазора, включают в
электромагнитную модель. Их вихревые потери могут быть сопоставимы с
потерями обмотки.

\section{ПАРАЗИТНЫЕ ПАРАМЕТРЫ И ЭЛЕКТРОМАГНИТНАЯ СОВМЕСТИМОСТЬ}
\label{sec:parasitics}

\subsection{Собственная ёмкость и собственный резонанс}

\phys Ёмкости между витками, слоями, обмоткой и сердечником создают
параллельный путь для высокочастотного тока. С ростом частоты модуль
индуктивного сопротивления сначала увеличивается, достигает максимума около
первого собственного резонанса, а затем дроссель может стать ёмкостным.
Поэтому высокая низкочастотная индуктивность не гарантирует подавления помех.

\begin{figure}[htbp]
\centering
\includegraphics[width=0.72\textwidth]{winding_capacitance_topologies.png}
\caption{Пример перехода от геометрии многослойной обмотки к сети
распределённых ёмкостей (по обзору Jia с соавторами~\cite{jia})}
\label{fig:winding-capacitance-network}
\end{figure}

На рисунке~\ref{fig:winding-capacitance-network} приведена обмотка
трансформатора, однако показанный принцип применим и к дросселю: каждая
пара соседних витков, слоёв и проводящих деталей образует частичную
ёмкость. Численные значения из этой схемы переносить на дроссель нельзя;
по ней определяют только состав узлов собственной расчётной модели.

До изготовления образца эквивалентную ёмкость для заданного распределения
потенциалов определяют из энергии электрического поля
\begin{equation}
 C_{\text{соб,расч}}=
 \frac{2W_e}{U_{\text{исп}}^2},
 \label{eq:self-capacitance-energy}
\end{equation}
\begin{whereblock}
\whereitem{$C_{\text{соб,расч}}$}{расчётная эквивалентная собственная
ёмкость, Ф;}
\whereitem{$W_e$}{энергия электрического поля электростатической модели,
Дж;}
\whereitem{$U_{\text{исп}}$}{испытательное напряжение между выводами,
В.}
\end{whereblock}

\method В электростатической модели задают фактические витки или секции,
толщину и относительную диэлектрическую проницаемость изоляции, сердечник,
каркас, экран, теплоотвод и корпус. Потенциалы секций в первом приближении
задают линейно от входного к выходному выводу. Для обмотки рядом с
переключаемым узлом дополнительно решают не менее трёх задач: напряжение между
выводами; «горячий» вывод относительно сердечника; «холодный» вывод
относительно сердечника. По ним строят трёхвыводную сеть частичных ёмкостей
«вывод~-- вывод~-- сердечник» по методике Zhao с
соавторами~\cite{zhao}. Одна величина $C_{\text{соб}}$ не воспроизводит
несколько резонансов многослойной обмотки.

\subsubsection{Физический смысл собственной ёмкости}

Собственная ёмкость дросселя не является отдельным конденсатором,
расположенным в одной точке. Она образована совокупностью частичных ёмкостей:
между соседними витками, между слоями, между началом и концом обмотки, между
обмоткой и магнитопроводом, между обмоткой и экраном, а также между обмоткой и
окружающими проводящими деталями. Напряжение на этих ёмкостях различно,
поэтому простое сложение их геометрических значений не даёт корректной
двухвыводной ёмкости.

Для инженерного анализа различают следующие величины:
\begin{enumerate}
\item $C_{\text{соб}}$ --- эквивалентная двухвыводная собственная ёмкость,
которая вместе с индуктивностью воспроизводит первый собственный резонанс;
\item $C_{\text{обм--маг}}$ --- ёмкость между обмоткой и магнитопроводом при
оговорённом соединении обоих выводов обмотки;
\item $C_{\text{обм--корп}}$ --- ёмкость между обмоткой и корпусом либо
защитным заземлением;
\item $C_{ij}$ --- частичная ёмкость между конкретными электрическими узлами
$i$ и $j$ распределённой модели.
\end{enumerate}

Эти величины отвечают на разные вопросы. Двухвыводная $C_{\text{соб}}$
определяет прежде всего первый резонанс и высокочастотный обход индуктивной
ветви между выводами. Ёмкости к магнитопроводу и корпусу определяют передачу
синфазной помехи. Измерение только $C_{\text{соб}}$ не позволяет вычислить
синфазный ток, если неизвестно, какая часть электрического поля замыкается на
корпус и какой потенциал имеет магнитопровод.

\begin{table}[htbp]
\begin{singlespace}
\caption{Влияние паразитной ёмкости дросселя}
\label{tab:capacitance-effects}
\centering
\small
\begin{tabularx}{\textwidth}{>{\raggedright\arraybackslash}p{34mm}
>{\raggedright\arraybackslash}p{50mm}
>{\raggedright\arraybackslash}X}
\toprule
Объект влияния & Физический механизм & Практическое последствие \\
\midrule
Первый собственный резонанс &
Компенсация индуктивной и ёмкостной реактивностей &
Снижение $C_{\text{соб}}$ повышает $f_{\text{СР}}$ при неизменной $L$ \\
Импеданс выше резонанса &
Высокочастотный ток проходит между выводами через распределённую ёмкостную
сеть &
Дроссель перестаёт ограничивать помеху как индуктивность и может стать путём
её обхода \\
Дифферен\-циальная помеха &
Ёмкость между входной и выходной частями обмотки передаёт составляющую
напряжения поперёк последовательного дросселя &
Ослабление фильтра после резонанса оказывается меньше расчёта по идеальной
индуктивности \\
Синфазная помеха &
Ёмкости обмотка--магнитопровод--корпус проводят ток, пропорциональный
$du/dt$ &
Растут кондуктивные и излучаемые помехи; результат зависит от подключения
магнитопровода, экрана и корпуса \\
Звон и перенапряжение &
Ёмкость образует колебательный контур с индуктивностью дросселя и
индуктивностями соединений &
Возникают узкие пики спектра, дополнительное напряжение на изоляции и
чувствительность к длине проводников \\
Потери и нагрев &
Через частичные ёмкости протекают токи смещения; энергия рассеивается в
изоляции, проводниках, экранах и демпфирующих элементах &
Растут диэлектрические, поверхностные и вихревые потери; особенно при большом
$du/dt$ \\
Распределение напряжения &
На высоких частотах напряжение распределяется по виткам нелинейно &
Первые или последние витки могут испытывать повышенную напряжённость
электрического поля \\
Повторяемость изделия &
Ёмкость зависит от шага намотки, прижима слоёв, толщины изоляции, пропитки,
положения выводов и корпуса &
Резонанс и ослабление имеют технологический разброс, который нельзя оценить
по одному образцу \\
\bottomrule
\end{tabularx}
\end{singlespace}
\end{table}

Энергия электрического поля, запасённая в эквивалентной ёмкости при
мгновенном напряжении $u$, равна
\begin{equation}
 W_C=\frac{C_{\text{экв}}u^2}{2},
 \label{eq:capacitance-energy}
\end{equation}
\begin{whereblock}
\whereitem{$W_C$}{энергия электрического поля, Дж;}
\whereitem{$C_{\text{экв}}$}{эквивалентная ёмкость рассматриваемого пути,
Ф;}
\whereitem{$u$}{мгновенное напряжение на этой ёмкости, В.}
\end{whereblock}

Для синусоидального напряжения диэлектрические потери первого приближения
оценивают по формуле
\begin{equation}
 P_{\text{диэл}}=
 2\pi f C_{\text{экв}}U_{\text{д}}^2\tan\delta_{\text{диэл}},
 \label{eq:dielectric-capacitance-loss}
\end{equation}
\begin{whereblock}
\whereitem{$P_{\text{диэл}}$}{потери в диэлектрике, Вт;}
\whereitem{$f$}{частота гармоники напряжения, Гц;}
\whereitem{$U_{\text{д}}$}{действующее значение напряжения гармоники, В;}
\whereitem{$\tan\delta_{\text{диэл}}$}{тангенс угла диэлектрических потерь.}
\end{whereblock}

Для несинусоидального напряжения потери суммируют по существенным гармоникам
либо вычисляют во временной полевой модели. Формула
\eqref{eq:dielectric-capacitance-loss} не учитывает потери в проводящих
деталях от токов смещения и поэтому не заменяет электромагнитно-тепловую
проверку.

\begin{figure}[htbp]
\centering
\includegraphics[
  angle=180,
  origin=c,
  width=0.72\textwidth
]{electrostatic_field_distribution.png}
\caption{Распределение электрического поля между обмоткой, сердечником и
окружающей областью (по Zhao с соавторами~\cite[рисунок~2]{zhao})}
\label{fig:electrostatic-field}
\end{figure}

По карте на рисунке~\ref{fig:electrostatic-field} находят участки с
наибольшей энергией электрического поля. Сначала изменяют положение
«горячих» витков и площадь их перекрытия с сердечником или экраном, затем
повторно вычисляют матрицу частичных ёмкостей.

При известной измеренной частоте первого собственного резонанса
эквивалентную собственную ёмкость оценивают по формуле
\begin{equation}
 C_{\text{соб}}=
 \frac{1}{(2\pi f_{\text{СР}})^2L_{\text{изм}}},
 \label{eq:self-capacitance}
\end{equation}
где $C_{\text{соб}}$ --- эквивалентная собственная ёмкость, Ф;\par\noindent
$f_{\text{СР}}$ --- измеренная частота первого собственного резонанса,
Гц;\par\noindent
$L_{\text{изм}}$ --- индуктивность в малосигнальном режиме на частоте,
не менее чем в десять раз ниже $f_{\text{СР}}$, Гн.

При одноветвевой эквивалентной схеме проводимость дросселя записывают как
\begin{equation}
 Y(f)=
 \frac{1}{R_{\text{AC}}(f)+j2\pi fL_{\text{диф}}(f)}
 +j2\pi f C_{\text{соб}},
 \label{eq:admittance}
\end{equation}
где $Y(f)$ --- комплексная проводимость, См;\par\noindent
$j$ --- мнимая единица;\par\noindent
$L_{\text{диф}}(f)$ --- малосигнальная дифференциальная индуктивность на
частоте $f$, Гн.

Комплексное сопротивление равно
\begin{equation}
 Z(f)=\frac{1}{Y(f)},
 \label{eq:impedance}
\end{equation}
где $Z(f)$ --- комплексное сопротивление дросселя, Ом.

\subsection{Расчёт или измерение собственной ёмкости}
\label{subsec:calculate-or-measure-capacitance}

\criterion Собственную ёмкость следует предварительно рассчитывать и
обязательно измерять на готовом образце. Эти операции не являются
взаимозаменяемыми. Расчёт нужен до изготовления для выбора жизнеспособной
геометрии и оценки запаса по собственному резонансу. Измерение нужно после
изготовления для идентификации фактической модели, учёта технологического
разброса и окончательной настройки фильтра.

Если отказаться от расчёта и измерить ёмкость только в конце, может
обнаружиться, что первый резонанс лежит внутри требуемой полосы подавления.
Исправление внешними элементами в этом случае ограничено: демпфирующая цепь
может уменьшить высоту резонансного пика, но не возвращает дросселю
индуктивный характер выше собственного резонанса. Для существенного повышения
$f_{\text{СР}}$ обычно требуется изменение самой обмотки.

В однорезонансном приближении частоту первого собственного резонанса
оценивают по формуле
\begin{equation}
 f_{\text{СР}}\approx
 \frac{1}{2\pi\sqrt{L_{\text{эф}}C_{\text{соб}}}},
 \label{eq:self-resonance-frequency}
\end{equation}
\begin{whereblock}
\whereitem{$f_{\text{СР}}$}{частота первого собственного резонанса, Гц;}
\whereitem{$L_{\text{эф}}$}{эффективная малосигнальная индуктивность в
области первого резонанса, Гн;}
\whereitem{$C_{\text{соб}}$}{эквивалентная двухвыводная собственная ёмкость,
Ф.}
\end{whereblock}

Формула~\eqref{eq:self-resonance-frequency} описывает только первый режим
простой параллельной RLC-схемы. Для многослойной, фольговой или планарной
обмотки $L_{\text{эф}}$ и распределение напряжения зависят от частоты, а
несколько частичных ёмкостей создают несколько резонансов. Поэтому значение,
полученное из $L$ и $f_{\text{СР}}$, является параметром эквивалентной модели,
а не геометрической суммой всех ёмкостей.

Если задана верхняя частота $f_{\text{треб,max}}$, до которой дроссель должен
оставаться индуктивным, допустимую собственную ёмкость первого приближения
определяют по формуле
\begin{equation}
 C_{\text{соб,max}}=
 \frac{1}{\left(2\pi k_f f_{\text{треб,max}}\right)^2L_{\max}},
 \label{eq:max-self-capacitance}
\end{equation}
\begin{whereblock}
\whereitem{$C_{\text{соб,max}}$}{предельная собственная ёмкость по условию
запаса до первого резонанса, Ф;}
\whereitem{$k_f$}{коэффициент частотного запаса, принимаемый от 3 до 5;}
\whereitem{$f_{\text{треб,max}}$}{верхняя частота, до которой требуется
индуктивный характер импеданса, Гц;}
\whereitem{$L_{\max}$}{наибольшая эффективная индуктивность среди
рассматриваемых режимов, Гн.}
\end{whereblock}

Для проверки ослабления одновременно используют минимальную индуктивность,
поскольку она даёт наименьшее индуктивное сопротивление. Иными словами,
$L_{\max}$ ограничивает минимальную частоту резонанса, а $L_{\min}$ ---
минимальное ослабление ниже резонанса. Один расчёт при номинальном значении
$L$ не охватывает оба худших случая.

Чувствительность резонансной частоты к малым изменениям параметров оценивают
по формуле
\begin{equation}
 \frac{\Delta f_{\text{СР}}}{f_{\text{СР}}}\approx
 -\frac{1}{2}
 \left(
 \frac{\Delta L}{L}+\frac{\Delta C_{\text{соб}}}{C_{\text{соб}}}
 \right),
 \label{eq:self-resonance-sensitivity}
\end{equation}
\begin{whereblock}
\whereitem{$\Delta f_{\text{СР}}$}{изменение частоты первого собственного
резонанса, Гц;}
\whereitem{$\Delta L$}{изменение индуктивности, Гн;}
\whereitem{$\Delta C_{\text{соб}}$}{изменение собственной ёмкости, Ф.}
\end{whereblock}

Например, увеличение ёмкости на $20\,\%$ при неизменной индуктивности снижает
частоту резонанса приблизительно на $10\,\%$. Это соотношение применимо только
к малым отклонениям и однорезонансной модели.

\begin{table}[htbp]
\begin{singlespace}
\caption{Распределение задач между расчётом и измерением}
\label{tab:calculation-versus-measurement}
\centering
\small
\begin{tabularx}{\textwidth}{>{\raggedright\arraybackslash}p{31mm}
>{\raggedright\arraybackslash}p{49mm}
>{\raggedright\arraybackslash}X}
\toprule
Этап & Что определяют & Решение по результату \\
\midrule
Эскизный аналитический расчёт &
Порядок величины $C_{\text{соб}}$, ожидаемый $f_{\text{СР}}$, влияние числа
слоёв, перекрытия и расстояний &
Отбрасывают варианты, у которых расчётный резонанс не имеет требуемого
запаса; точное совпадение с образцом не требуется \\
Электро\-статическая модель &
Энергию поля и матрицу частичных ёмкостей между выводами, слоями,
магнитопроводом, экраном и корпусом &
Выбирают порядок слоёв, положение «горячего» вывода, секционирование и
подключение экрана \\
Измерение отдельного дросселя &
$Z(f)$, фазу, резонансы, эквивалентные $R_s(f)$, $L_s(f)$ и
$C_{\text{соб}}$ при заданном состоянии магнитопровода &
Идентифицируют модель реального компонента и проверяют запас по
$f_{\text{СР}}$ \\
Измерение пассивного фильтра &
Коэффициент передачи или вносимое ослабление при заданных импедансах
источника и нагрузки &
Выбирают внешние конденсаторы, демпфирование и взаимное расположение
компонентов \\
Измерение работающего изделия &
Спектр дифференциальной и синфазной помехи в фактических режимах &
Подтверждают выполнение нормы с запасом; это испытание нельзя заменить
измерением одного дросселя \\
Серийная проверка &
Разброс резонансов и импеданса для образцов, изготовленных в разных точках
поля технологических допусков &
Назначают контролируемые параметры намотки и приёмочные пределы \\
\bottomrule
\end{tabularx}
\end{singlespace}
\end{table}

\subsubsection{Что допускается подстраивать после измерения}

По измеренной характеристике отдельного дросселя допускается:
\begin{enumerate}
\item подобрать внешнюю демпфирующую RC- или RCD-цепь для снижения добротности
конкретного резонанса;
\item скорректировать номиналы конденсаторов и индуктивностей соседней
ступени фильтра с учётом фактических импедансов источника и нагрузки;
\item изменить взаимное расположение компонентов, длину соединений,
ориентацию начала и конца обмотки, точку подключения экрана и путь возврата
синфазного тока;
\item разделить одну ступень фильтра на две ступени, если требуемая полоса
слишком широка для одного резонансного элемента.
\end{enumerate}

Внешняя подстройка не устраняет внутреннюю ёмкостную передачу между частями
обмотки. При подключении дополнительной ёмкости параллельно дросселю частота
резонанса первого приближения становится
\begin{equation}
 f'_{\text{СР}}\approx
 \frac{1}{2\pi\sqrt{L_{\text{эф}}
 \left(C_{\text{соб}}+C_{\text{доп}}\right)}},
 \label{eq:self-resonance-added-capacitance}
\end{equation}
\begin{whereblock}
\whereitem{$f'_{\text{СР}}$}{частота резонанса после подключения
дополнительной параллельной ёмкости, Гц;}
\whereitem{$C_{\text{доп}}$}{дополнительная внешняя ёмкость, Ф.}
\end{whereblock}

Следовательно, параллельный конденсатор снижает, а не повышает частоту
собственного резонанса. Он может быть частью намеренно настроенного фильтра,
но не является способом исправить слишком большую $C_{\text{соб}}$.
Последовательная компенсация создаёт узкополосный нуль или новый резонанс,
чувствительный к нагрузке и допускам, и не заменяет уменьшение внутренней
ёмкости для широкополосного подавления.

Если собственный резонанс расположен слишком низко, изменяют конструкцию
обмотки: уменьшают число витков или слоёв при сохранении требуемой
индуктивности другим магнитопроводом; уменьшают площадь перекрытия витков с
большой разностью потенциалов; применяют секционирование; увеличивают
межслойное расстояние; выбирают последовательность намотки с меньшей энергией
электрического поля; удаляют «горячие» витки от магнитопровода, корпуса и
экрана. Каждое изменение повторно проверяют по потерям меди, коэффициенту
заполнения окна, изоляции и тепловому режиму.

\subsubsection{Рекомендуемая последовательность проектирования}

\method Расчёт и измерение выполняют в следующем порядке:
\begin{enumerate}
\item по спектру преобразователя задают $f_{\text{треб,max}}$ и требуемое
ослабление; частоту коммутации нельзя автоматически принимать равной
$f_{\text{треб,max}}$, поскольку существенны её гармоники и резонансные пики;
\item по формуле~\eqref{eq:max-self-capacitance} задают целевое
$C_{\text{соб,max}}$ и сравнивают варианты обмотки аналитически либо по
электростатической модели;
\item изготавливают макет с воспроизводимой фиксацией витков, выводов,
изоляции и магнитопровода;
\item измеряют комплексный импеданс отдельного дросселя, затем коэффициент
передачи пассивного фильтра и после этого спектр работающего изделия;
\item по измерениям корректируют эквивалентную модель и выполняют только те
изменения, которые воздействуют на выявленный путь помехи; после изменения
намотки или экрана полный цикл измерений повторяют.
\end{enumerate}

\recommend Для единичного низкочастотного дросселя, у которого измеренный
$f_{\text{СР}}$ более чем в пять раз выше всей требуемой полосы, подробная
трёхмерная электростатическая модель обычно не нужна. Для многослойной,
фольговой или планарной обмотки, для напряжения с большим $du/dt$, а также при
работе вблизи первого резонанса расчёт частичных ёмкостей следует выполнять
до изготовления. Во всех случаях итоговое решение принимают по измеренному
$Z(f)$ и результату испытания полной схемы.

\subsubsection{Расчётный пример оценки допустимой ёмкости}

\textbf{1. Исходные данные.}

Для иллюстрации приняты:
\begin{enumerate}
\item индуктивность $L=60\,\text{мкГн}$;
\item частота коммутации $f_{\text{с}}=150\,\text{кГц}$;
\item верхняя частота требуемого индуктивного поведения
$f_{\text{треб,max}}=500\,\text{кГц}$;
\item коэффициент частотного запаса $k_f=5$;
\item сравниваемые значения $C_{\text{соб}}$ от $20$ до
$200\,\text{пФ}$.
\end{enumerate}

\textbf{2. Принятые допущения.}

Дроссель представлен параллельной LC-схемой без потерь; индуктивность не
зависит от частоты, тока и температуры; рассматривается только первый
собственный резонанс. Пример служит для определения порядка величин, а не для
приёмки конструкции.

\textbf{3. Расчётная схема и методика.}

Сначала определяют индуктивное сопротивление на частоте коммутации, затем
частоту резонанса для каждого значения ёмкости и предельную ёмкость по
заданному частотному запасу.

\textbf{4. Пошаговый расчёт.}

\textbf{Шаг 1. Определение индуктивного сопротивления на частоте
коммутации.}

Индуктивное сопротивление равно
\begin{equation}
 X_L=2\pi f_{\text{с}}L,
 \label{eq:example-inductive-reactance}
\end{equation}
\begin{whereblock}
\whereitem{$X_L$}{индуктивное сопротивление, Ом;}
\whereitem{$f_{\text{с}}$}{частота коммутации, Гц;}
\whereitem{$L$}{индуктивность дросселя, Гн.}
\end{whereblock}

\subst
\begin{equation}
 X_L=2\pi\cdot150\cdot10^3\cdot60\cdot10^{-6}
 =56{,}5\,\text{Ом}.
\end{equation}

\result На основной частоте идеальная индуктивная ветвь имеет сопротивление
$56{,}5\,\text{Ом}$. Это число не характеризует подавление гармоник выше
первого резонанса.

\textbf{Шаг 2. Определение частоты первого собственного резонанса при
$C_{\text{соб}}=50\,\text{пФ}$.}

По формуле~\eqref{eq:self-resonance-frequency}:
\begin{equation}
 f_{\text{СР}}=
 \frac{1}{2\pi\sqrt{60\cdot10^{-6}\cdot50\cdot10^{-12}}}
 =2{,}91\,\text{МГц}.
\end{equation}

\result Резонанс расположен в $19{,}4$ раза выше частоты коммутации, но
только в $5{,}82$ раза выше принятой верхней частоты требуемого индуктивного
поведения.

\textbf{Шаг 3. Определение предельной собственной ёмкости.}

По формуле~\eqref{eq:max-self-capacitance}:
\begin{equation}
 C_{\text{соб,max}}=
 \frac{1}{
 \left(2\pi\cdot5\cdot500\cdot10^3\right)^2
 \cdot60\cdot10^{-6}}
 =67{,}5\,\text{пФ}.
\end{equation}

\result По условию пятикратного запаса расчётная собственная ёмкость не
должна превышать $67{,}5\,\text{пФ}$.

\textbf{5. Результаты.}

\begin{table}[htbp]
\begin{singlespace}
\caption{Влияние собственной ёмкости на первый резонанс при
$L=60\,\text{мкГн}$}
\label{tab:self-capacitance-example}
\centering
\small
\begin{tabular}{>{\raggedleft\arraybackslash}p{32mm}
>{\raggedleft\arraybackslash}p{38mm}
>{\raggedleft\arraybackslash}p{40mm}}
\toprule
$C_{\text{соб}}$, пФ & $f_{\text{СР}}$, МГц &
$f_{\text{СР}}/f_{\text{треб,max}}$ \\
\midrule
20  & 4,59 & 9,18 \\
50  & 2,91 & 5,82 \\
100 & 2,05 & 4,10 \\
200 & 1,45 & 2,90 \\
\bottomrule
\end{tabular}
\end{singlespace}
\end{table}

\textbf{6. Вывод.}

При заданном пятикратном запасе варианты $20$ и $50\,\text{пФ}$ проходят
предварительный критерий, а варианты $100$ и $200\,\text{пФ}$ требуют
переработки обмотки либо изменения архитектуры фильтра. Даже прошедший
вариант $50\,\text{пФ}$ должен быть проверен измерением комплексного
импеданса, поскольку реальный дроссель имеет потери и несколько
распределённых резонансов.

\subsection{Построение и чтение частотных характеристик}

Термин «амплитудно-частотная характеристика дросселя» употребляют
неоднозначно. Для отдельного двухвыводного компонента строго определяют не
коэффициент передачи, а частотную зависимость комплексного импеданса
$Z(f)=U(f)/I(f)$. Коэффициент передачи появляется только после включения
дросселя в цепь с источником и нагрузкой. Поэтому в протоколе указывают
измеряемый объект и режим, а не только слово «АЧХ».

\begin{table}[htbp]
\begin{singlespace}
\caption{Различие частотных измерений дросселя, фильтра и преобразователя}
\label{tab:frequency-measurement-types}
\centering
\small
\begin{tabularx}{\textwidth}{>{\raggedright\arraybackslash}p{30mm}
>{\raggedright\arraybackslash}p{42mm}
>{\raggedright\arraybackslash}p{44mm}
>{\raggedright\arraybackslash}X}
\toprule
Объект & Измеряемая величина & Что показывает & Что не показывает \\
\midrule
Отдельный обесточенный дроссель &
$|Z(f)|$, $\arg Z(f)$, $R_s(f)$, $L_s(f)$, $Q(f)$ &
Индуктивную и ёмкостную области, собственные резонансы, потери и пригодность
эквивалентной модели &
Фактическое ослабление фильтра и соответствие изделия нормам
электромагнитной совместимости \\
Обесточенный пассивный фильтр &
$H(f)=U_{\text{вых}}/U_{\text{вх}}$, параметр $S_{21}$ или вносимое
ослабление &
Передачу малого сигнала через всю ступень при заданных импедансах источника и
нагрузки &
Спектр помехи, создаваемой работающими ключами, и изменение режима
преобразователя \\
Работающий преобразователь с LISN &
Спектр напряжения помехи, дБмкВ, при заданном детекторе и полосе &
Конечный результат совместного действия источника помехи, путей связи,
фильтра, компоновки и нагрузки &
Передаточную характеристику отдельного дросселя \\
Контур регулирования работающего преобразователя &
Петлевой коэффициент передачи $T(j2\pi f)$, запас по фазе и амплитуде &
Устойчивость системы управления &
Собственную ёмкость дросселя напрямую; это отдельное измерение \\
\bottomrule
\end{tabularx}
\end{singlespace}
\end{table}

\phys Измерение отдельного дросселя отвечает на вопрос, какой импеданс
компонент создаёт для малого возмущения каждой частоты. Измерение пассивного
фильтра отвечает на вопрос, какая доля внешнего тестового сигнала проходит от
его входа к выходу. Измерение работающего преобразователя отвечает на вопрос,
какой спектр помехи реально поступает в нормированный измерительный тракт.
Эти три результата связаны, но один не может быть непосредственно заменён
другим.

Для пассивного четырёхполюсника коэффициент передачи по напряжению равен
\begin{equation}
 H(f)=\frac{U_{\text{вых}}(f)}{U_{\text{вх}}(f)},
 \label{eq:filter-transfer-function}
\end{equation}
\begin{whereblock}
\whereitem{$H(f)$}{комплексный коэффициент передачи фильтра;}
\whereitem{$U_{\text{вых}}(f)$}{комплексное выходное напряжение, В;}
\whereitem{$U_{\text{вх}}(f)$}{комплексное входное напряжение, В.}
\end{whereblock}

Ослабление по отношению к тому же стенду без исследуемого фильтра определяют
по формуле
\begin{equation}
 A_{\text{вн}}(f)=
 20\lg\left|
 \frac{U_{\text{вых,без}}(f)}{U_{\text{вых,ф}}(f)}
 \right|,
 \label{eq:measured-insertion-loss}
\end{equation}
\begin{whereblock}
\whereitem{$A_{\text{вн}}(f)$}{измеренное вносимое ослабление, дБ;}
\whereitem{$U_{\text{вых,без}}(f)$}{выходное напряжение стенда без фильтра,
В;}
\whereitem{$U_{\text{вых,ф}}(f)$}{выходное напряжение того же стенда с
фильтром, В.}
\end{whereblock}

Значение $S_{21}$, измеренное в системе $50\,\text{Ом}$, характеризует
именно систему $50\,\text{Ом}$. В силовом преобразователе импедансы источника
помехи, нагрузки и LISN частотно зависимы и обычно не равны
$50\,\text{Ом}$ одновременно. Поэтому лабораторное измерение $S_{21}$ удобно
для сравнения вариантов, но окончательное ослабление проверяют в
представительной схеме.

\method Частотную характеристику отдельного дросселя строят как минимум по
модулю $|Z(f)|$ и фазе $\varphi_Z(f)=\arg Z(f)$. Порядок получения
характеристики:
\begin{enumerate}
\item задают диапазон от частоты не выше $f_{\text{с}}/10$ до частоты не
ниже верхней границы нормируемых помех или $5f_{\text{СР}}$;
\item измеряют дроссель анализатором импеданса либо векторным анализатором
цепей с малым сигналом; выполняют компенсацию разомкнутого и
короткозамкнутого приспособления непосредственно в плоскости выводов;
\item повторяют измерение при минимальной, номинальной и максимальной
постоянной составляющей тока, если прибор и устройство подмагничивания это
допускают;
\item строят $|Z|$ в логарифмическом масштабе, $\varphi_Z$, эквивалентные
$R_s(f)$, $L_s(f)$ и добротность $Q(f)$;
\item отмечают рабочую частоту, существенные гармоники, полосу нормируемых
помех, первый резонанс и все последующие резонансы;
\item накладывают требуемый импеданс или требуемое ослабление полной
схемы, после чего принимают решение по критериям ниже.
\end{enumerate}

\begin{figure}[htbp]
\centering
\includegraphics[width=0.82\textwidth]{impedance_measurement_setup.png}
\caption{Пример измерения частотной зависимости импеданса магнитного
компонента анализатором импеданса (по Zhao с соавторами~\cite{zhao})}
\label{fig:impedance-measurement-setup}
\end{figure}

При сборке стенда по рисунку~\ref{fig:impedance-measurement-setup}
калибровку выполняют на концах измерительного приспособления. Длинные
провода между калибровочной плоскостью и дросселем добавляют собственные
индуктивность и ёмкость и могут создать ложный резонанс.

\subsubsection{Условия измерения отдельного дросселя}

\method В протоколе фиксируют амплитуду тестового сигнала, постоянный ток,
температуру, положение магнитопровода и выводов, состояние экрана и корпуса,
тип приспособления и выполненную компенсацию. Без этих данных две кривые
$Z(f)$ могут различаться из-за стенда, а не из-за конструкции дросселя.

\begin{table}[htbp]
\begin{singlespace}
\caption{Обязательные варианты измерения паразитных параметров}
\label{tab:required-parasitic-measurements}
\centering
\small
\begin{tabularx}{\textwidth}{>{\raggedright\arraybackslash}p{35mm}
>{\raggedright\arraybackslash}p{48mm}
>{\raggedright\arraybackslash}X}
\toprule
Вариант измерения & Соединение и режим & Назначение \\
\midrule
Двухвыводный импеданс без подмагничивания &
Измерение между началом и концом обмотки; магнитопровод и экран находятся в
состоянии, соответствующем изделию &
Базовые $Z(f)$, $f_{\text{СР}}$, $R_s(f)$, $L_s(f)$ и $Q(f)$ \\
Двухвыводный импеданс с постоянным током &
Ток задают через развязывающую цепь или штатное устройство подмагничивания;
малый переменный сигнал накладывают на рабочую точку &
Дифференциальная индуктивность и положение резонансов при минимальном,
номинальном и максимальном токе \\
Измерение при температуре &
Повторение ключевых сканирований при минимальной и максимальной рабочей
температуре после установления теплового равновесия &
Учёт изменения сопротивления обмотки, проницаемости и потерь
магнитопровода \\
Обмотка относительно магнитопровода &
Оба вывода обмотки соединяют вместе; ёмкость измеряют до магнитопровода при
оговорённом состоянии экрана и корпуса &
Оценка полного синфазного пути обмотка--магнитопровод; результат не содержит
распределения потенциала вдоль обмотки \\
Начало и конец относительно магнитопровода &
Трёхвыводное измерение или идентификация матрицы частичных ёмкостей при
нескольких независимых соединениях &
Оценка асимметрии связи и выбор ориентации «горячего» вывода \\
Дроссель в фактической сборке &
Те же положения относительно печатной платы, теплоотвода, корпуса и соседних
компонентов, что в изделии &
Выявление внешних ёмкостей и взаимной магнитной связи, отсутствующих у
свободно лежащего компонента \\
\bottomrule
\end{tabularx}
\end{singlespace}
\end{table}

При измерении с постоянным током проверяют допустимые ток и энергию для
анализатора, развязывающего устройства и приспособления. Постоянную
составляющую не подают непосредственно на высокочастотный порт прибора, если
это прямо не допускается его документацией. Амплитуду малого сигнала выбирают
так, чтобы измерялась локальная дифференциальная индуктивность без заметного
перемещения рабочей точки.

При двухвыводном измерении нельзя произвольно оставлять магнитопровод
плавающим, если в изделии он соединён с экраном или ёмкостно связан с
корпусом. Обратная ошибка также недопустима: принудительное заземление
магнитопровода на стенде создаёт синфазный путь, которого может не быть в
устройстве. Требуемое состояние задают в программе испытаний.

Для векторного анализатора цепей выполняют калибровку в измерительной
плоскости и, при необходимости, исключение параметров приспособления
(de-embedding). Для анализатора импеданса выполняют компенсацию холостого
хода и короткого замыкания. Если измеряемый импеданс приближается к пределам
динамического диапазона метода, результат подтверждают другим включением
приспособления или другим прибором.

В индуктивной области параметры последовательной эквивалентной схемы
определяют по формулам
\begin{align}
 R_s(f)&=\operatorname{Re}Z(f),
 \label{eq:series-resistance-from-z}\\
 L_s(f)&=\frac{\operatorname{Im}Z(f)}{2\pi f},
 \label{eq:series-inductance-from-z}\\
 Q(f)&=\frac{|\operatorname{Im}Z(f)|}{\operatorname{Re}Z(f)}.
 \label{eq:q-from-z}
\end{align}
\begin{whereblock}
\whereitem{$R_s(f)$}{эквивалентное последовательное сопротивление, Ом;}
\whereitem{$L_s(f)$}{эквивалентная последовательная индуктивность при
$\operatorname{Im}Z>0$, Гн;}
\whereitem{$Q(f)$}{добротность дросселя;}
\whereitem{$\operatorname{Re}Z$, $\operatorname{Im}Z$}{действительная и
мнимая части комплексного сопротивления, Ом.}
\end{whereblock}

Первый собственный резонанс определяют по переходу
$\operatorname{Im}Z$ через нуль от положительного значения к отрицательному
и по локальному максимуму $|Z|$. Один только максимум модуля без фазы может
быть ошибочно принят за резонанс.

\begin{table}[htbp]
\begin{singlespace}
\caption{Информация, извлекаемая из частотной зависимости импеданса}
\label{tab:information-from-impedance}
\centering
\small
\begin{tabularx}{\textwidth}{>{\raggedright\arraybackslash}p{35mm}
>{\raggedright\arraybackslash}p{49mm}
>{\raggedright\arraybackslash}X}
\toprule
Наблюдаемый признак & Определяемая величина & Интерпретация \\
\midrule
$|Z|$ при низкой частоте и фаза около $0^\circ$ &
Сопротивление обмотки и приспособления &
После компенсации приспособления приближается к $R_{\text{DC}}$; анализатор
импеданса не всегда обеспечивает точность четырёхпроводного микроомметра \\
Рост $|Z|$ приблизительно на $20\,\text{дБ}$ на декаду и положительная фаза &
Индуктивная область и $L_s(f)$ &
Показывает полосу, где дроссель допустимо представлять последовательной
RL-схемой \\
Отклонение $L_s(f)$ и рост $R_s(f)$ до резонанса &
Частотную зависимость потерь и магнитных параметров &
Может быть вызвано поверхностным эффектом, эффектом близости, потерями
магнитопровода и началом влияния распределённой ёмкости \\
Максимум $|Z|$ совместно с переходом фазы через $0^\circ$ от плюса к минусу &
Первый собственный резонанс $f_{\text{СР}}$ &
Выше этой точки одноветвевая модель обычно имеет ёмкостный характер \\
Высота и ширина первого резонанса &
Эквивалентное параллельное сопротивление и добротность резонанса &
Высокий узкий пик означает малое демпфирование и повышенную вероятность
звона при возбуждении соответствующей гармоникой \\
Спад $|Z|$ приблизительно на $20\,\text{дБ}$ на декаду и отрицательная фаза &
Ёмкостную область и эффективную последовательную ёмкость &
Высокочастотный ток обходит индуктивную ветвь; одна измеренная ёмкость может
быть достаточна только вдали от последующих резонансов \\
Несколько максимумов и минимумов &
Распределённые режимы обмотки и связи с внешними проводниками &
Требуется многозвенная модель; формула с одной $L$ и одной
$C_{\text{соб}}$ непригодна во всей полосе \\
Сдвиг кривой при постоянном токе &
Изменение дифференциальной индуктивности и потерь из-за подмагничивания &
Проверяет применимость дросселя в реальной рабочей точке \\
Сдвиг кривой при температуре &
Изменение сопротивления меди, проницаемости и потерь магнитопровода &
Показывает худший режим, который может отсутствовать при комнатной
температуре \\
Изменение кривой при касании корпуса, экрана или магнитопровода &
Ёмкостную связь с внешними проводящими деталями &
Указывает, что двухвыводное измерение с плавающим магнитопроводом не
соответствует фактической сборке \\
\bottomrule
\end{tabularx}
\end{singlespace}
\end{table}

В ёмкостной области эффективную последовательную ёмкость для локального
участка характеристики можно оценить по формуле
\begin{equation}
 C_s(f)=
 -\frac{1}{2\pi f\operatorname{Im}Z(f)},
 \qquad \operatorname{Im}Z(f)<0,
 \label{eq:series-capacitance-from-z}
\end{equation}
где $C_s(f)$ --- эффективная последовательная ёмкость, Ф.

Полученное $C_s(f)$ нельзя автоматически считать равным
$C_{\text{соб}}$, найденной по формуле~\eqref{eq:self-capacitance}. Первая
величина описывает локальное ёмкостное поведение выше резонанса, вторая ---
положение первого резонанса в упрощённой параллельной модели. Их расхождение
является нормальным признаком распределённого компонента.

Для одиночного слабо искажённого параллельного резонанса добротность по
полосе оценивают как
\begin{equation}
 Q_{\text{рез}}\approx
 \frac{f_{\text{СР}}}{f_2-f_1},
 \label{eq:resonance-quality-bandwidth}
\end{equation}
\begin{whereblock}
\whereitem{$Q_{\text{рез}}$}{добротность первого параллельного резонанса;}
\whereitem{$f_1$, $f_2$}{нижняя и верхняя частоты, на которых
$|Z|=|Z|_{\max}/\sqrt{2}$, Гц.}
\end{whereblock}

Эту оценку не используют при перекрывающихся резонансах, асимметричном пике
или недостаточном динамическом диапазоне прибора. В таких случаях параметры
идентифицируют подгонкой комплексной многозвенной модели одновременно по
действительной и мнимой частям $Z(f)$.

Измерение одной цифры «ёмкость» LCR-метром на произвольной частоте не
заменяет построение $Z(f)$. Прибор вычисляет последовательный или параллельный
эквивалент из измеренных комплексных данных. В индуктивной области результат
в режиме измерения ёмкости может быть отрицательным, нестабильным или зависеть
от выбранной эквивалентной схемы. Для определения первого резонанса проводят
частотное сканирование, а $C_{\text{соб}}$ затем находят по резонансу и
индуктивности либо идентифицируют из комплексной модели.

Для последовательного дросселя между эквивалентными импедансами источника и
нагрузки ожидаемое ослабление относительно схемы без дросселя оценивают по
формуле
\begin{equation}
 A_Z(f)=20\lg
 \left|
 1+\frac{Z(f)}{Z_s(f)+Z_l(f)}
 \right|,
 \label{eq:series-insertion-loss}
\end{equation}
\begin{whereblock}
\whereitem{$A_Z(f)$}{расчётное ослабление напряжения, дБ;}
\whereitem{$Z_s(f)$}{эквивалентный импеданс источника помехи, Ом;}
\whereitem{$Z_l(f)$}{эквивалентный импеданс нагрузки или LISN, Ом.}
\end{whereblock}

\subsubsection{Последовательность измерений от компонента к изделию}

\method Частотную проверку выполняют тремя последовательными уровнями.

\textbf{Уровень 1 --- отдельный дроссель.} Измеряют комплексный импеданс
$Z(f)$ в обесточенном состоянии, при требуемом постоянном подмагничивании и
температуре. На этом уровне определяют, что именно ограничивает компонент:
недостаточная индуктивность, потери, первый резонанс, последующие резонансы
или связь с магнитопроводом и корпусом.

\textbf{Уровень 2 --- пассивный фильтр.} Преобразователь обесточивают, а в
исследуемый тракт вводят малый тестовый сигнал. Измеряют $H(f)$, $S_{21}$ или
вносимое ослабление при заданных импедансах источника и нагрузки. Этот уровень
учитывает внешние конденсаторы, индуктивности дорожек, взаимную связь
компонентов и демпфирование. Если измерение выполнено только в тракте
$50\,\text{Ом}$, результат не переносят в силовую схему без проверки её
реальных граничных импедансов.

\textbf{Уровень 3 --- работающее изделие.} При номинальном питании и нагрузке
измеряют спектр напряжения помехи на LISN, а при необходимости --- токовой
пробой и пробами ближнего поля. Это не АЧХ дросселя и не передаточная АЧХ
фильтра: частота по оси абсцисс является частотой спектральной составляющей,
созданной самим преобразователем. Результат зависит одновременно от спектра
коммутации, режима нагрузки, импедансов сети, компоновки, корпуса и фильтра.

Если требуется именно передаточная характеристика работающей схемы, задают
конкретные входной и выходной порты, способ малосигнального возбуждения и
рабочую точку. Такое измерение выполняют отдельным методом инжекции и не
смешивают со спектром электромагнитных помех или с измерением петлевого
коэффициента передачи регулятора.

\criterion Решение считается подтверждённым только при согласованности всех
трёх уровней: отдельный дроссель имеет требуемый импеданс; пассивный фильтр
обеспечивает требуемое вносимое ослабление; работающее изделие выполняет норму
электромагнитной совместимости с заданным запасом. Успех на одном уровне не
гарантирует успех на следующем.

\begin{figure}[htbp]
\centering
\includegraphics[width=0.78\textwidth]{self_parasitics_attenuation.png}
\caption{Расхождение измеренного ослабления фильтра с идеализированными
моделями из-за собственных паразитных параметров компонентов
(по Jia с соавторами~\cite[рисунок~1]{jia})}
\label{fig:self-parasitics-attenuation}
\end{figure}

Рисунок~\ref{fig:self-parasitics-attenuation} показывает, что после
собственных резонансов реальное ослабление может быть на десятки децибел
хуже расчёта по идеальным $L$ и $C$. Точные частоты и уровни относятся к
исследованному фильтру; для проектируемого устройства вывод делают только
по его измеренным $Z(f)$ и передаточной характеристике.

Требуемое ослабление получают из исходного спектра устройства без
проектируемого фильтрующего элемента
\begin{equation}
 A_{\text{треб}}(f)=
 U_{\text{баз}}(f)-U_{\text{норм}}(f)+M_{\text{ЭМС}},
 \label{eq:required-attenuation}
\end{equation}
\begin{whereblock}
\whereitem{$A_{\text{треб}}(f)$}{требуемое ослабление, дБ;}
\whereitem{$U_{\text{баз}}(f)$}{измеренный или рассчитанный уровень помехи
до установки проектируемого фильтрующего элемента, дБмкВ;}
\whereitem{$U_{\text{норм}}(f)$}{предельный уровень по применимому
стандарту, дБмкВ;}
\whereitem{$M_{\text{ЭМС}}$}{проектный запас по электромагнитной
совместимости, дБ.}
\end{whereblock}

\recommend При отсутствии корпоративного требования принимают
$M_{\text{ЭМС}}=6\,\text{дБ}$. Формула~\eqref{eq:series-insertion-loss}
пригодна для первой линейной оценки; окончательное ослабление рассчитывают в
схеме с паразитными ёмкостями и индуктивностями и подтверждают измерением с
тем же детектором, полосой и LISN, которые требует стандарт.

\criterion Частота $f_{\text{треб,max}}$ означает верхнюю частоту, до которой
дроссель должен оставаться индуктивным, а не автоматически верхнюю границу
измерительного стандарта. Для первого макета принимают
$f_{\text{СР}}\geq3f_{\text{треб,max}}$; для критичных помех ---
$f_{\text{СР}}\geq5f_{\text{треб,max}}$. Во всей требуемой полосе должны
одновременно выполняться $\operatorname{Im}Z>0$ и
$A_Z(f)\geq A_{\text{треб}}(f)$. Если это невозможно, изменяют секционирование
и порядок слоёв, уменьшают перекрытие витков, выбирают меньшее число витков
или вводят дополнительную фильтрующую ступень. Управление обмоточной
ёмкостью подтверждено моделями и измерениями в
\cite{zhao,heldwein}.

\begin{table}[htbp]
\begin{singlespace}
\caption{Выбор корректирующего действия по результатам измерений}
\label{tab:capacitance-corrective-actions}
\centering
\small
\begin{tabularx}{\textwidth}{>{\raggedright\arraybackslash}p{48mm}
>{\raggedright\arraybackslash}p{50mm}
>{\raggedright\arraybackslash}X}
\toprule
Результат измерения & Наиболее вероятная причина & Первое корректирующее
действие \\
\midrule
$f_{\text{СР}}<3f_{\text{треб,max}}$ у отдельного дросселя &
Чрезмерная собственная ёмкость или слишком большая индуктивность для данной
полосы &
Переработать геометрию обмотки; проверить возможность меньшего числа витков и
другого магнитопровода \\
Резонанс расположен достаточно высоко, но пик очень узкий и высокий &
Недостаточное демпфирование &
Рассчитать внешнюю демпфирующую цепь и проверить её потери; не изменять
намотку без подтверждения причины \\
$Z(f)$ дросселя соответствует расчёту, но пассивный фильтр имеет провал
ослабления &
Паразитные параметры печатной платы, конденсатора, соединений или взаимная
связь компонентов &
Изменить компоновку и идентифицировать полную схему фильтра по $S$-параметрам
или комплексным импедансам \\
Пассивный фильтр имеет требуемое ослабление, но спектр на LISN превышает
норму &
Непредставительные импедансы стенда, синфазный обход фильтра, изменение
источника помехи или связь через корпус &
Разделить дифференциальную и синфазную составляющие и локализовать путь
возврата помехи \\
Результат резко меняется при перевороте дросселя или подключении корпуса &
Асимметричная ёмкость «горячих» витков к окружающим деталям &
Изменить ориентацию начала обмотки, расстояние до корпуса либо точку
подключения экрана \\
Параметры заметно расходятся между образцами &
Нестабильный шаг намотки, положение выводов, прижим слоёв, толщина изоляции
или пропитка &
Закрепить технологические параметры и ввести контроль $f_{\text{СР}}$ либо
$|Z|$ на выбранных частотах \\
\bottomrule
\end{tabularx}
\end{singlespace}
\end{table}

\subsection{Синфазный ток и электрическое поле}

\phys Синфазный ток возникает, когда участки обмотки с большим
$du/dt$ ёмкостно связаны с сердечником, корпусом, теплоотводом или защитным
заземлением. Ток проходит через эти паразитные ёмкости и возвращается к
преобразователю вне расчётного силового контура. Последствиями являются рост
кондуктивных помех на LISN, излучение кабелей, ложные срабатывания
чувствительных цепей и дополнительный ток через изоляцию. Пиковый ток одного
фронта не сравнивают непосредственно с нормой в дБмкВ: для сравнения нужен
спектр и измерительный тракт.

Оценку импульса синфазного тока через ёмкостную связь с переключаемым
узлом выполняют по формуле
\begin{equation}
 i_{\text{СФ}}=
 C_{\text{СФ}}\frac{du_{\text{пер}}}{dt},
 \label{eq:common-mode-current}
\end{equation}
где $i_{\text{СФ}}$ --- синфазный ток, А;\par\noindent
$C_{\text{СФ}}$ --- эквивалентная ёмкость связи обмотки с
переключаемым узлом, корпусом или теплоотводом, Ф;\par\noindent
$u_{\text{пер}}$ --- напряжение переключаемого узла, В.

Для гармоники частоты $f$ действующее значение синфазного тока равно
\begin{equation}
 I_{\text{СФ}}(f)=
 2\pi f C_{\text{СФ}}|U_{\text{пер}}(f)|,
 \label{eq:common-mode-spectrum}
\end{equation}
\begin{whereblock}
\whereitem{$I_{\text{СФ}}(f)$}{действующее значение гармоники синфазного
тока, А;}
\whereitem{$U_{\text{пер}}(f)$}{действующее значение гармоники напряжения
переключаемого узла, В.}
\end{whereblock}

Ожидаемое напряжение на измерительном тракте определяют по формуле
\begin{equation}
 U_{\text{LISN,СФ}}(f)=
 \left|Z_{\text{тр}}(f)I_{\text{СФ}}(f)\right|,
 \label{eq:lisn-common-mode-voltage}
\end{equation}
\begin{whereblock}
\whereitem{$U_{\text{LISN,СФ}}(f)$}{напряжение синфазной помехи на выходе
измерительного тракта, В;}
\whereitem{$Z_{\text{тр}}(f)$}{передаточный импеданс фактических LISN,
кабелей, корпуса и приёмника, Ом.}
\end{whereblock}

Уровень напряжения в децибелах относительно $1\,\text{мкВ}$ равен
\begin{equation}
 U_{\text{СФ,дБмкВ}}=
 20\lg\left(
 \frac{U_{\text{LISN,СФ}}}{1\,\text{мкВ}}
 \right).
 \label{eq:dbuv}
\end{equation}

\method Ёмкость $C_{\text{СФ}}$ получают из трёхвыводной
электростатической модели или измеряют между соединёнными выводами обмотки и
сердечником либо корпусом при оговорённом состоянии остальных проводников.
Значение $Z_{\text{тр}}$ получают из схемы применяемой LISN и стенда. Число
$25\,\text{Ом}$ допустимо только для конкретного идеализированного соединения
двух ветвей по $50\,\text{Ом}$ и не является универсальной константой.
Расчёт повторяют для всех существенных гармоник фронта и сопоставляют с
$U_{\text{норм}}(f)-M_{\text{ЭМС}}$.

\source{Механизмы синфазных и дифференциальных путей распространения помех,
а также влияние паразитных ёмкостей рассмотрены в
\cite{zhao,jia,heldwein}.}

\begin{figure}[htbp]
\centering
\includegraphics[width=0.66\textwidth]{capacitance_cancellation_circuit.png}
\caption{Пример компенсационной обмотки для уменьшения ёмкостной передачи
высокочастотной помехи (по обзору Jia с соавторами~\cite{jia})}
\label{fig:capacitance-cancellation}
\end{figure}

Компенсационная обмотка на рисунке~\ref{fig:capacitance-cancellation}
создаёт ток, противоположный току через паразитную ёмкость. Такой приём
применяют только после определения амплитуды и фазы паразитной передачи:
ошибка числа витков или полярности повышает помеху, а добавленная обмотка
может создать новый резонанс.

\recommend Для уменьшения электрического поля и синфазных помех:
\begin{enumerate}
 \item размещают вывод и витки с наибольшим $du/dt$ дальше от корпуса,
 теплоотвода, заземлённых полигонов и низкопотенциальных витков;
 \item минимизируют площадь перекрытия участков обмотки с большим
 переменным напряжением;
 \item порядок слоёв выбирают по результатам расчёта матрицы частичных
 ёмкостей, а не только по минимальной длине проводника;
 \item электростатический экран подключают в одной определённой точке и
 не замыкают в короткозамкнутый виток;
 \item после добавления экрана повторно определяют $P_{\text{экр}}$,
 $C_{\text{СФ}}$ и $f_{\text{СР}}$.
\end{enumerate}

\criterion Конструкция по синфазным помехам приемлема, если рассчитанный и
затем измеренный спектр на LISN ниже нормы не менее чем на
$M_{\text{ЭМС}}$ во всём диапазоне. Если превышение создают отдельные
резонансные пики, сначала уменьшают соответствующую ёмкостную связь и
площадь контура; увеличение индуктивности дросселя без анализа пути
возврата синфазного тока может не дать улучшения.

\subsection{Поле рассеяния зазора}

\phys Поле рассеяния создают выпучивание потока у зазора и площадь контура
обмотки. Переменное поле, проникая в медные полигоны, теплоотводы, крепёж и
экраны, наводит вихревые токи. Обычный проводящий металл не «усиливает» поле:
он перераспределяет его и нагревается из-за вихревых потерь. Ферромагнитная
деталь может дополнительно концентрировать поток и поэтому моделируется с
фактической магнитной характеристикой. Нагрев окружающих деталей уменьшает
КПД и может повысить температуру обмотки даже при малых собственных потерях
сердечника.

\begin{figure}[htbp]
\centering
\includegraphics[width=0.76\textwidth]{slotted_shield_eddy_currents.png}
\caption{Разрыв замкнутого пути вихревого тока прорезью в проводящем
экране (по обзору Jia с соавторами~\cite{jia})}
\label{fig:slotted-shield-eddy-current}
\end{figure}

Прорезь на рисунке~\ref{fig:slotted-shield-eddy-current} разрывает
короткозамкнутый виток в плоскости экрана и уменьшает соответствующие
вихревые потери. При этом токи по другим путям, нагрев кромок и ёмкостная
связь экрана с обмоткой сохраняются, поэтому экран включают в
электромагнитную и тепловую модели целиком.

Потери в окружающей проводящей детали определяют по объёму
\begin{equation}
 P_{\text{мет}}=
 \int_{V_{\text{мет}}}
 \frac{|\mathbf J_{\text{вихр}}|^2}
 {\sigma_{\text{мет}}}\dd V,
 \label{eq:metal-eddy-loss}
\end{equation}
\begin{whereblock}
\whereitem{$P_{\text{мет}}$}{вихревые потери в рассматриваемой детали, Вт;}
\whereitem{$V_{\text{мет}}$}{объём детали, м$^3$;}
\whereitem{$\mathbf J_{\text{вихр}}$}{плотность наведённого вихревого тока,
А/м$^2$;}
\whereitem{$\sigma_{\text{мет}}$}{электрическая проводимость материала
детали при расчётной температуре, См/м.}
\end{whereblock}

Для оценки воздействия поля на малый контур чувствительной цепи используют
закон электромагнитной индукции
\begin{equation}
 U_{\text{инд,действ}}\approx
 2\pi f B_{\perp,\text{действ}}A_{\text{конт}},
 \label{eq:induced-loop-voltage}
\end{equation}
\begin{whereblock}
\whereitem{$U_{\text{инд,действ}}$}{действующее значение наведённого
напряжения, В;}
\whereitem{$B_{\perp,\text{действ}}$}{действующая составляющая магнитной
индукции, нормальная к площади контура, Тл;}
\whereitem{$A_{\text{конт}}$}{площадь контура чувствительной цепи, м$^2$.}
\end{whereblock}

\method До моделирования задают контрольные объекты: ближайший медный
полигон, теплоотвод, крепёж, корпус и контуры измерительных либо управляющих
цепей. Для каждого указывают материал, толщину, расстояние и допустимый
нагрев или наведённое напряжение. В частотной модели рассчитывают каждую
существенную гармонику, суммируют потери и строят карты
$|\mathbf H|$, $|\mathbf B|$ и $|\mathbf J_{\text{вихр}}|$. Результат без
координат контрольной поверхности не воспроизводим.

Для уменьшения поля рассеяния:
\begin{enumerate}
 \item применяют несколько меньших зазоров вместо одного большого,
 если это допускает конструкция;
 \item удаляют проводники, переходы и металлические детали от областей
 максимального выпучивания поля;
 \item располагают прямой и обратный высокочастотные токи рядом, чтобы
 уменьшить площадь токового контура;
 \item проверяют поле $\mathbf H$ и плотность вихревых потерь в
 окружающих проводниках на трёхмерной модели;
 \item оценивают излучаемое поле на контрольной поверхности и
 чувствительность результата к положению проводов.
\end{enumerate}

\recommend Магнитный экран применяют после оптимизации геометрии
зазора и токового контура. Экран может уменьшить внешнее поле, но
увеличить массу, температуру и потери. Его эффективность подтверждают
расчётом и измерением в составе преобразователя.

\criterion Для каждой металлической детали должно выполняться
$P_{\text{мет}}\leq P_{\text{мет,доп}}$ и ограничение по температуре; для
каждого чувствительного контура ---
$U_{\text{инд,действ}}\leq U_{\text{доп}}$. При отсутствии числового
требования сравнивают не менее трёх вариантов положения, выбирают вариант с
минимальными потерями и полем и фиксируют контрольную поверхность для
последующего измерения. Влияние поля зазора на распределение тока и потери
обмотки подтверждено в \cite{schaefer,ewald_litz,meng}.

\section{ВЕРИФИКАЦИЯ МЕТОДОМ КОНЕЧНЫХ ЭЛЕМЕНТОВ}
\label{sec:fem}

\subsection{Виды постановок и источники свойств}

\phys FEM не заменяет исходные данные и не является одной универсальной
операцией. Магнитостатическая задача определяет потокосцепление и насыщение,
частотная электромагнитная --- распределение тока и вихревые потери,
электростатическая --- частичные ёмкости, тепловая --- температуры. Геометрия
обмотки к началу этих расчётов уже должна быть выбрана по
разделу~\ref{sec:winding}.

Настоящий раздел объединяет в одну последовательность полевые проверки,
необходимость которых была установлена при расчёте потерь по
разделу~\ref{sec:losses} и паразитных параметров по
разделу~\ref{sec:parasitics}. Сначала фиксируют геометрию кандидата, затем
последовательно решают магнитную, частотную, электростатическую и тепловую
задачи; изменение геометрии возвращает расчёт к началу этой
последовательности.

\begin{table}[H]
\caption{Исходные данные полевых моделей и места их получения}
\label{tab:fem-input-sources}
\centering
\scriptsize
\begin{tabularx}{\textwidth}{>{\raggedright\arraybackslash}p{38mm}
>{\raggedright\arraybackslash}p{53mm} X}
\toprule
Данные & Что вводят в модель & Откуда получают \\
\midrule
Магнитопровод & фактические размеры, скругления, стыки, прокладки,
ориентацию половин & размерный чертёж и измерение образца \\
Нелинейность феррита & семейство $B(H,T)$ либо $\mu(B,T)$ во всём диапазоне
тока & паспорт материала; при отсутствии --- измеренная кривая конкретного
комплекта \\
Потери феррита & карта $p_v(f,\Delta B,B_{\text{DC}},T)$ или локальные
$k$, $\alpha$, $\beta$ & паспортные кривые и собственные измерения; за
пределами карты экстраполяцию не считают подтверждением \\
Проводники & $\rho(\vartheta)$, $\mu_r$, покрытие, точные сечения и
соединения & таблица~\ref{tab:copper-properties}, паспорт провода и чертёж
обмотки \\
Литцендрат & явные жилы либо проверенная однородная модель, шаги скрутки,
коэффициент упаковки & паспорт кабеля и калибровка модели по
$R_{\text{AC}}(f)$ \\
Диэлектрики & толщины, $\varepsilon_r(f,T)$, электрическая проводимость &
паспорт эмали, каркаса, прокладок и печатной платы \\
Окружающие детали & геометрия, $\sigma(T)$, $\mu_r(B,T)$ металлических
деталей & сборочный чертёж и паспорт материала \\
Тепловые свойства & теплопроводность, теплоёмкость, плотность, контактные
сопротивления & паспорта материалов; контакты и коэффициенты теплоотдачи
калибруют испытанием \\
\bottomrule
\end{tabularx}
\end{table}

\subsection{Последовательность решения}
\label{subsec:fem-sequence}

\begin{enumerate}
\item В нелинейной магнитостатической задаче рассчитывают набор токов от
нуля до аварийного максимума при минимальной и максимальной температуре.
Сохраняют $\Psi(I,T)$, $B_{\max}$, энергию и поле на контрольной поверхности.
\item Дифференциальную индуктивность получают малосигнальным возмущением или
центральной конечной разностью. Расчёт повторяют при нескольких приращениях
тока.
\item Для каждой существенной гармоники решают частотную задачу с
фактической проводимостью и геометрией. Из интеграла джоулевых потерь получают
$R_{\text{AC}}(f_h)$ и потери в окружающем металле.
\item В электростатической задаче получают матрицу частичных ёмкостей,
$C_{\text{соб}}$ и $C_{\text{СФ}}$ для нескольких комбинаций потенциалов.
\item Распределения потерь передают в тепловую задачу. Обновлённые
температуры возвращают в электромагнитные свойства до сходимости.
\end{enumerate}

\subsection{Минимальный состав геометрической модели}

Модель FEM должна включать:
\begin{enumerate}
 \item нелинейные характеристики $B(H,T)$ магнитного материала;
 \item фактическую геометрию всех зазоров и немагнитных прокладок;
 \item обмотку с достаточной детализацией для расчёта локальных потерь;
 \item окружающие проводящие детали, экраны, теплоотвод и полигоны
 печатной платы в области существенного поля;
 \item граничные условия на удалении, исключающем влияние границы
 расчётной области на контролируемые величины.
\end{enumerate}

\recommend Сначала решают двумерную модель симметричного поперечного
сечения для перебора, затем трёхмерную модель окончательного варианта с
выводами и окружающими деталями. Внешнюю границу удаляют или применяют
бесконечные элементы; размер области увеличивают до изменения $L$,
$B_{\max}$ и поля на контрольной поверхности менее чем на
$1\,\%$.

Секущую индуктивность определяют по формуле
\begin{equation}
 L_{\text{сек}}(I,T)=\frac{\Psi(I,T)}{I},
 \label{eq:secant-inductance}
\end{equation}
\begin{whereblock}
\whereitem{$L_{\text{сек}}$}{секущая индуктивность, Гн;}
\whereitem{$\Psi(I,T)$}{потокосцепление обмотки, вычисленное решателем как
сумма потоков всех витков, Вб$\cdot$виток;}
\whereitem{$I$}{постоянный ток рабочей точки, А;}
\whereitem{$T$}{температура магнитопровода, К.}
\end{whereblock}

Дифференциальную индуктивность определяют по формуле
\begin{equation}
 L_{\text{диф}}(I,T)=\frac{d\Psi(I,T)}{dI}.
 \label{eq:differential-inductance}
\end{equation}

При отсутствии малосигнального решателя производную определяют центральной
разностью
\begin{equation}
 L_{\text{диф}}(I,T)\approx
 \frac{\Psi(I+\Delta I,T)-\Psi(I-\Delta I,T)}
 {2\Delta I},
 \label{eq:differential-inductance-fd}
\end{equation}
\begin{whereblock}
\whereitem{$\Delta I$}{малое приращение тока, А.}
\end{whereblock}

\recommend Сначала принимают $\Delta I=0{,}01I_{\max}$, затем повторяют с
$0{,}005I_{\max}$ и $0{,}02I_{\max}$. Если $L_{\text{диф}}$ изменяется более
чем на $1\,\%$, уменьшают шаг тока и сгущают сетку. В точке $I=0$ используют
одностороннюю разность или малосигнальный решатель.

\recommend Для ограничения пульсации тока используют
$L_{\text{диф}}$, а для потокосцепления и энергии --- полную
характеристику $\Psi(I,T)$. Подмена этих величин одним паспортным
$A_L$ недопустима при заметном насыщении.

Сопротивление обмотки для гармоники получают из рассчитанных джоулевых
потерь
\begin{equation}
 R_{\text{AC}}(f_h,\vartheta)=
 \frac{P_{\text{Cu},h}}
 {I_{h,\text{действ}}^2},
 \label{eq:rac-from-fem}
\end{equation}
\begin{whereblock}
\whereitem{$P_{\text{Cu},h}$}{джоулевы потери в обмотке от одной гармоники,
Вт;}
\whereitem{$I_{h,\text{действ}}$}{заданное действующее значение тока
гармоники, А.}
\end{whereblock}

\recommend Если частотный решатель использует амплитуду, а не действующее
значение фазора, перед вычислением формулы~\eqref{eq:rac-from-fem} приводят
ток и потери к одной системе. В отчёте явно записывают соглашение решателя;
иначе возможна ошибка в два раза по мощности.

\subsection{Настройка зазора по модели}

После первого решения длину зазора допускается корректировать по
отношению требуемой и рассчитанной секущих индуктивностей:
\begin{equation}
 g_{r+1}=g_r
 \frac{L_{\text{сек},r}}{L_{\text{треб}}},
 \label{eq:gap-update}
\end{equation}
где $g_r$, $g_{r+1}$ --- суммарный зазор на текущей и следующей
итерациях, м;\par\noindent
$L_{\text{сек},r}$ --- рассчитанная секущая индуктивность на текущей
итерации, Гн;\par\noindent
$L_{\text{треб}}$ --- требуемая индуктивность в контрольной точке, Гн;\par\noindent
$r$ --- номер итерации.

\recommend Формула~\eqref{eq:gap-update} является локальным
итерационным правилом, а не физическим законом. При сильной
нелинейности изменяют $g$ ограниченными шагами и на каждой итерации
повторяют нелинейное решение.

\subsection{Сходимость сетки и энергетическая проверка}

Относительное изменение контролируемой величины при измельчении сетки
определяют как
\begin{equation}
 \varepsilon_Q=
 \frac{\left|Q_m-Q_{m-1}\right|}
 {\max\!\left(\left|Q_m\right|,Q_{\min}\right)},
 \label{eq:mesh-convergence}
\end{equation}
где $\varepsilon_Q$ --- относительное изменение результата;\par\noindent
$Q_m$, $Q_{m-1}$ --- значения контролируемой величины на двух
последовательных сетках;\par\noindent
$Q_{\min}$ --- малое опорное значение, предотвращающее деление на
число, близкое нулю.

Энергию поля в линейной контрольной точке сопоставляют с
\begin{equation}
 W_{\text{FEM}}\approx\frac{L_{\text{сек}}I^2}{2},
 \label{eq:fem-energy-check}
\end{equation}
где $W_{\text{FEM}}$ --- энергия, интегрированная по расчётной области,
Дж.

\recommend Используют не менее трёх последовательно сгущаемых сеток.
Контролируют отдельно $L$, $B_{\max}$, $P_{\text{Cu}}$,
$P_{\text{Fe}}$ и потери в окружающих деталях. Локальный максимум в
остром геометрическом угле не принимают без исследования зависимости
от радиуса скругления и сетки.

\begin{table}[H]
\caption{Обязательные результаты FEM-верификации}
\label{tab:fem-results}
\centering
\small
\begin{tabularx}{\textwidth}{>{\raggedright\arraybackslash}p{50mm} X}
\toprule
Постановка & Результат, передаваемый в следующий шаг \\
\midrule
Нелинейная магнитостатическая &
таблица $\Psi(I,T)$, $L_{\text{сек}}$, $L_{\text{диф}}$, локальный
$B_{\max}$ и карта внешнего поля \\
Частотная электромагнитная &
$R_{\text{AC}}(f_h,T)$, $P_{\text{Cu},h}$, потери каждого металлического
объекта, чувствительность к положению обмотки \\
Электростатическая &
матрица частичных ёмкостей, $C_{\text{соб}}$, $C_{\text{СФ}}$, ожидаемые
резонансы \\
Тепловая &
$\vartheta_{\text{обм,max}}$, $\vartheta_{\text{серд,max}}$, температуры
окружающих деталей и сходимость совместной итерации \\
Численная проверка &
три сетки, размер внешней области, $\varepsilon_Q$ для каждой принимаемой
величины и энергетический баланс \\
\bottomrule
\end{tabularx}
\end{table}

\section{ТЕПЛОВОЙ РАСЧЁТ И СОВМЕСТНАЯ ИТЕРАЦИЯ}
\label{sec:thermal}

Объёмную плотность тепловыделения каждого компонента задают по формуле
\begin{equation}
 q_{v,k}=\frac{P_k}{V_k},
 \label{eq:heat-density}
\end{equation}
где $q_{v,k}$ --- средняя объёмная плотность тепловыделения $k$-го
компонента, Вт/м$^3$;\par\noindent
$P_k$ --- потери в компоненте, Вт;\par\noindent
$V_k$ --- объём компонента, м$^3$.

Для сосредоточенной предварительной модели превышение температуры
оценивают по формуле
\begin{equation}
 \Delta\vartheta=P_{\Sigma}R_{\theta},
 \label{eq:thermal-resistance}
\end{equation}
где $\Delta\vartheta$ --- превышение температуры, К;\par\noindent
$R_{\theta}$ --- тепловое сопротивление от дросселя к окружающей среде,
К/Вт.

\recommend Значение $R_{\theta}$ не является константой формы
сердечника. Его получают из тепловой модели или испытания при заданных
ориентации, корпусе, обдуве, печатной плате и расстоянии до соседних
компонентов.

Допустимую температуру обмотки назначают как наименьшее из ограничений
эмали, общей изоляции литцендрата, каркаса, пропитки, припоя и требований
надёжности. Допустимую температуру магнитопровода назначают как наименьшее из
паспортной рабочей температуры материала, температуры, до которой имеются
данные $B(H,T)$ и потерь, и ограничения клея либо прокладок. Температуру Кюри
в этот перечень как рабочий предел не включают.

Электромагнитный и тепловой расчёты выполняют итерационно:
\begin{enumerate}
 \item задают начальные температуры сердечника и обмотки;
 \item вычисляют свойства материалов и потери;
 \item передают распределение потерь в тепловую модель;
 \item обновляют температуры и повторяют расчёт;
 \item прекращают итерации, когда изменение максимальной температуры и
 суммарных потерь не превышает установленного допуска.
\end{enumerate}

\criterion Для учебного расчёта принимают критерий сходимости
$1\,\text{К}$ по максимальной температуре и $1\,\%$ по суммарным потерям на
двух последовательных итерациях. Если расчётная температура превышает
допустимую, сначала уменьшают источник потерь, затем тепловое сопротивление;
увеличение класса нагревостойкости не восстанавливает потерянный КПД.

\section{ОПТИМИЗАЦИЯ И КРИТЕРИИ ПРИЁМКИ}
\label{sec:optimization}

\subsection{Переменные проекта и целевые функции}

\phys Вектор здесь --- условная упорядоченная запись проектных параметров.
Она может содержать числовые, целочисленные и категориальные параметры,
например марку материала. Вектор проектных переменных описывает конструкцию,
а вектор целевых функций --- числовые результаты, которые требуется
уменьшить или увеличить. Запись вектора сама по себе не выполняет
оптимизацию.

Конструкцию представляют вектором
\begin{equation}
 \mathbf x=
 [m_c,\;s_c,\;N,\;n_g,\;g_i,\;d_s,\;n_s,\;
 \mathcal L_w,\;r_g,\;\mathcal C],
 \label{eq:design-vector}
\end{equation}
\begin{whereblock}
\whereitem{$\mathbf x$}{вектор проектных переменных;}
\whereitem{$m_c$}{марка магнитного материала;}
\whereitem{$s_c$}{типоразмер магнитопровода;}
\whereitem{$N$}{число витков;}
\whereitem{$n_g$, $g_i$}{число зазоров и длина одного зазора, м;}
\whereitem{$d_s$, $n_s$}{диаметр и число жил литцендрата, м и шт.;}
\whereitem{$\mathcal L_w$}{схема послойной укладки обмотки;}
\whereitem{$r_g$}{минимальное расстояние от меди до кромки зазора, м;}
\whereitem{$\mathcal C$}{вариант охлаждения и экранирования.}
\end{whereblock}

Все рассматриваемые результаты приводят к направлению «меньше --- лучше»:
\begin{equation}
 \mathbf F=
 \left[
 \overline P_{\Sigma},\;
 P_{\Sigma,\max},\;
 V_{\text{изд}},\;
 -M_{\text{ДМ}},\;
 -M_{\text{СФ}},\;
 H_{\text{конт,max}}
 \right],
 \label{eq:objective-vector}
\end{equation}
\begin{whereblock}
\whereitem{$\mathbf F$}{вектор целевых функций;}
\whereitem{$\overline P_{\Sigma}$}{потери по профилю эксплуатации,
Вт;}
\whereitem{$P_{\Sigma,\max}$}{потери в наихудшем режиме, Вт;}
\whereitem{$V_{\text{изд}}$}{габаритный объём дросселя, м$^3$;}
\whereitem{$M_{\text{ДМ}}$, $M_{\text{СФ}}$}{минимальные запасы по
дифференциальным и синфазным помехам во всей нормируемой полосе, дБ;}
\whereitem{$H_{\text{конт,max}}$}{наибольшая напряжённость магнитного поля
на заданной контрольной поверхности, А/м.}
\end{whereblock}

Запас по синфазной помехе определяют по формуле
\begin{equation}
 M_{\text{СФ}}=
 \min_{f\in\mathcal B}
 \left[
 U_{\text{норм,СФ}}(f)-
 U_{\text{расч,СФ}}(f)
 \right],
 \label{eq:cm-margin}
\end{equation}
\begin{whereblock}
\whereitem{$\mathcal B$}{нормируемая полоса частот;}
\whereitem{$U_{\text{норм,СФ}}$}{предельный уровень синфазной помехи,
дБмкВ;}
\whereitem{$U_{\text{расч,СФ}}$}{расчётный или измеренный уровень
синфазной помехи, дБмкВ.}
\end{whereblock}

Запас $M_{\text{ДМ}}$ определяют аналогично после разделения помех на
дифференциальную и синфазную составляющие.

\subsection{Соответствие подходу Papamanolis и вложенный перебор}
\label{subsec:papamanolis-search}

Метод Papamanolis с соавторами~\cite{papamanolis} не является отдельной
методикой расчёта магнитопровода. Это способ организации оптимизации:
при заданных магнитопроводе и типе литцендрата перебирают частоту
коммутации $f_s$, относительный размах пульсации тока $r_I$ и целое число
витков $N$. Индуктивность $L$ следует из требуемой пульсации, а длина
воздушного зазора $g$ --- из сочетания $L$ и $N$. Для каждой точки
рассчитывают потери и температуру, после чего сравнивают только физически
реализуемые варианты.

\begin{table}[H]
\caption{Соответствие настоящей методики подходу Papamanolis}
\label{tab:papamanolis-compliance}
\centering
\scriptsize
\begin{tabularx}{\textwidth}{
>{\raggedright\arraybackslash}p{31mm}
>{\raggedright\arraybackslash}p{48mm}
>{\raggedright\arraybackslash}X}
\toprule
Элемент расчёта & Подход Papamanolis & Настоящая методика и результат сверки \\
\midrule
Пространство поиска &
$f_s$, $r_I$ или $L$, целое $N$ при заданных сердечнике и диаметре жилы &
Совместный перебор $N$ и $g$ уже предусмотрен; ниже явно добавлен внешний
цикл по $f_s$, когда частота является переменной \\
Электрический режим &
Понижающий преобразователь при постоянных напряжениях и коэффициенте
заполнения импульса $0{,}5$ &
Используются фактические $u_L(t)$ и $i_L(t)$, поэтому область применимости
шире и включает другие схемы и несимметричные формы \\
Потери в обмотке &
Постоянная составляющая, скин-эффект и эффект близости с учётом размещения
витков &
Соответствует; дополнительно введены гармоническая сумма, поле зазора,
потери выводов и окружающих деталей \\
Потери в магнитопроводе &
Улучшенное обобщённое уравнение Штейнмеца, измеренная карта потерь,
температура и постоянное подмагничивание &
Соответствует; использование локальных данных выбранного материала и
учёт $B_{\text{DC}}$ обязательны \\
Физические ограничения &
Насыщение, заполнение окна, предельный зазор, температура и укладка витков &
Соответствует; дополнительно проверяются допуски, собственный резонанс,
электромагнитная совместимость и поле рассеяния \\
Температурная итерация &
Электромагнитно-тепловой расчёт до сходимости потерь и температуры &
Соответствует разделу~\ref{sec:thermal} \\
Частичная нагрузка &
После поиска минимума при полной нагрузке сравниваются потери при
$50\,\%$ и $20\,\%$ нагрузки и стоимость &
Соответствует средневзвешенным потерям по формуле
\eqref{eq:mission-profile-loss}; профиль эксплуатации задаётся техническим
заданием \\
Выбор близких решений &
Сохраняются варианты с потерями до $1{,}2$ от глобального минимума для
последующего выбора по стоимости и частичной нагрузке &
Соответствует идее множества Парето; порог $2\,\%$ ниже применяется только
для выбора почти максимального КПД, а порог $20\,\%$ --- для широкого
исследовательского списка \\
\bottomrule
\end{tabularx}
\end{table}

\result Настоящая методика согласуется со структурой расчёта
Papamanolis, но не должна заменяться её упрощёнными формулами. Формулы
статьи используют для начальной оценки и построения сетки поиска;
окончательное решение получают по полной модели потерь, температурной
итерации и обязательным ограничениям.

\subsubsection{Параметризация задачи}

Относительный размах пульсации тока определяют по формуле
\begin{equation}
 r_I=
 \frac{\Delta I_{\text{пп}}}{I_{\text{DC}}},
 \label{eq:papamanolis-ripple}
\end{equation}
\begin{whereblock}
\whereitem{$r_I$}{отношение размаха пульсации тока к постоянной
составляющей тока;}
\whereitem{$\Delta I_{\text{пп}}$}{размах тока дросселя от минимального
до максимального значения, А;}
\whereitem{$I_{\text{DC}}$}{постоянная составляющая тока дросселя, А.}
\end{whereblock}

Для рассмотренного в статье понижающего преобразователя при
$U_{\text{вых}}/U_{\text{вх}}=0{,}5$ индуктивность связана с пульсацией
соотношением
\begin{equation}
 L=
 \frac{U_{\text{вых}}}
 {2f_s I_{\text{DC}}r_I},
 \label{eq:papamanolis-buck-l}
\end{equation}
\begin{whereblock}
\whereitem{$L$}{индуктивность дросселя, Гн;}
\whereitem{$U_{\text{вых}}$}{постоянное выходное напряжение
преобразователя, В;}
\whereitem{$f_s$}{частота коммутации преобразователя, Гц.}
\end{whereblock}

Формула~\eqref{eq:papamanolis-buck-l} не является общей формулой
понижающего преобразователя. При другом коэффициенте заполнения импульса,
изменяющихся напряжениях или иной топологии $L$ определяют по фактической
вольт-секундной площади и форме тока согласно
разделу~\ref{sec:electrical-mode}.

В сокращённой постановке статьи вектор переменных имеет вид
\begin{equation}
 \mathbf x_P=[f_s,\;r_I,\;N],
 \label{eq:papamanolis-vector}
\end{equation}
\begin{whereblock}
\whereitem{$\mathbf x_P$}{вектор переменных сокращённой задачи
Papamanolis;}
\whereitem{$N$}{целое число витков обмотки.}
\end{whereblock}

Значения $L$ и $g$ не перебирают независимо от $\mathbf x_P$: значение
$L$ рассчитывают из электрического режима, а $g$ --- из требуемого
$A_L=L/N^2$ по подразделу~\ref{subsec:turn-gap-loop}. При фиксированной
частоте из вектора~\eqref{eq:papamanolis-vector} исключают $f_s$ и
перебирают только $r_I$ или эквивалентно $L$, затем $N$.

\subsubsection{Аналитическая оценка числа витков}

При фиксированных $f_s$, $L$, сердечнике и проводнике статья сводит
потери в меди и магнитопроводе к зависимости
\begin{equation}
 P_0(N)=P_{\text{Cu}}+P_{\text{Fe}}=
 c_{\text{Cu}}N^2+c_{\text{Fe}}N^{-\beta},
 \label{eq:papamanolis-reduced-loss}
\end{equation}
\begin{whereblock}
\whereitem{$P_0$}{сумма потерь в обмотке и магнитопроводе в сокращённой
модели, Вт;}
\whereitem{$P_{\text{Cu}}$}{потери в обмотке, Вт;}
\whereitem{$P_{\text{Fe}}$}{потери в магнитопроводе, Вт;}
\whereitem{$c_{\text{Cu}}$}{положительный коэффициент потерь в обмотке
для фиксированной рабочей точки, Вт;}
\whereitem{$c_{\text{Fe}}$}{положительный коэффициент потерь в
магнитопроводе для фиксированной рабочей точки, Вт;}
\whereitem{$\beta$}{локальный показатель степени зависимости потерь
магнитопровода от амплитуды магнитной индукции.}
\end{whereblock}

Рост первой составляющей связан с увеличением длины и сопротивления
обмотки, а уменьшение второй --- со снижением переменной магнитной
индукции при увеличении $N$. В непрерывном приближении оптимальное число
витков равно
\begin{equation}
 N_{\text{opt}}=
 \left(
 \frac{\beta c_{\text{Fe}}}
 {2c_{\text{Cu}}}
 \right)^{1/(2+\beta)}.
 \label{eq:papamanolis-nopt}
\end{equation}

В точке~\eqref{eq:papamanolis-nopt} сокращённая модель даёт соотношение
\begin{equation}
 \left.
 \frac{P_{\text{Fe}}}{P_{\text{Cu}}}
 \right|_{N=N_{\text{opt}}}
 =
 \frac{2}{\beta}.
 \label{eq:papamanolis-loss-ratio}
\end{equation}

Формулы~\eqref{eq:papamanolis-nopt} и
\eqref{eq:papamanolis-loss-ratio} применяют только для начальной точки
перебора. Они предполагают непрерывное $N$, постоянный коэффициент
заполнения окна, степенную модель потерь и отсутствие активных ограничений.
В реальной конструкции целое число витков, укладка, насыщение, зазор,
постоянное подмагничивание и температура сдвигают минимум; равенство
\eqref{eq:papamanolis-loss-ratio} не является критерием приёмки.

Для оценки ширины области близких решений вводят относительные потери
\begin{equation}
 \lambda_N=
 \frac{P_0(N)}{P_0(N_{\text{opt}})}
 =
 \frac{2}{2+\beta}
 \left[
 \frac{\beta}{2}
 \left(\frac{N}{N_{\text{opt}}}\right)^2+
 \left(\frac{N}{N_{\text{opt}}}\right)^{-\beta}
 \right],
 \label{eq:papamanolis-flat-turns}
\end{equation}
\begin{whereblock}
\whereitem{$\lambda_N$}{отношение потерь при выбранном числе витков к
минимальным потерям сокращённой модели.}
\end{whereblock}

Для исследовательского допуска $\lambda_N\leq1{,}2$ и
$2<\beta<3$ статья приводит приближённые границы
\begin{align}
 N_{\text{кв,min}}&\approx
 (0{,}6343+0{,}0505\beta)N_{\text{opt}},
 \label{eq:papamanolis-nq-min}\\
 N_{\text{кв,max}}&\approx
 (1{,}4835-0{,}0610\beta)N_{\text{opt}},
 \label{eq:papamanolis-nq-max}
\end{align}
где $N_{\text{кв,min}}$, $N_{\text{кв,max}}$ --- нижняя и верхняя
границы квазиоптимального диапазона числа витков. Эти границы используют
для формирования широкого списка, но каждый целый вариант внутри него
пересчитывают полной моделью.

\subsubsection{Двухусловное начальное приближение}

Предлагаемое в статье руководство из двух уравнений в обозначениях
настоящей методики записывают как систему
\begin{equation}
 \begin{cases}
 B_{\text{DC}}(L,N)+\widehat B_{\text{AC}}(f_s,N)
 =B_{\text{доп}},\\
 N=N_{\text{opt}}(L,f_s).
 \end{cases}
 \label{eq:papamanolis-two-equation-guide}
\end{equation}

Первое условие задаёт решение вблизи ограничения по максимальной
индукции, второе уравновешивает тенденции потерь в меди и
магнитопроводе. Систему~\eqref{eq:papamanolis-two-equation-guide} решают
численно относительно $L$ и непрерывного $N$, затем проверяют соседние
целые числа витков. При сильном влиянии $B_{\text{DC}}$ на потери
магнитопровода точный минимум может находиться ниже границы насыщения,
поэтому это начальное приближение, а не замена перебора.

\paragraph{Основной режим: частота коммутации фиксирована.}

Если частота задана схемотехникой и алгоритмом управления, внешний
частотный цикл отсутствует. Расчёт выполняют в следующем порядке:
\begin{enumerate}
\item принимают $f_s=f_{\text{ном}}$ и для каждого режима профиля
восстанавливают $u_L(t)$, $i_L(t)$, гармоники и температуры;
\item задают сетку допустимых $r_I$ либо $L$ по требованиям к пульсации,
динамике регулирования и режиму непрерывного тока;
\item для каждой точки определяют диапазон целых $N$ от ограничения по
индукции до ограничения по окну;
\item для каждого $N$ рассчитывают $A_L=L/N^2$, физический зазор, его
число участков и поле выпучивания;
\item исключают варианты, нарушающие ограничения по $B_{\max}$,
$L_{\text{диф,min}}$, окну, технологической длине зазора и собственной
резонансной частоте;
\item строят фактическую укладку проводника и рассчитывают
$P_{\text{Cu}}$ по гармоникам, включая поле зазора;
\item рассчитывают $P_{\text{Fe}}$ по рабочей форме $B(t)$ с учётом
$f_s$, температуры и $B_{\text{DC}}$;
\item добавляют потери выводов, экранов и окружающих проводящих деталей,
после чего замыкают температурную итерацию;
\item для каждого $r_I$ или $L$ сохраняют вариант с минимальной целевой
функцией и несколько соседних целых $N$;
\item выбирают итоговую конструкцию по профилю эксплуатации и множеству
Парето, а не только по минимуму потерь при полной нагрузке.
\end{enumerate}

\paragraph{Перебор частоты --- применять только при переменной частоте
коммутации.}

\recommend Частоту $f_s$ включают в пространство поиска, только если
разработчик вправе выбрать её на этапе проектирования либо преобразователь
действительно работает в заданном диапазоне переменной частоты. Если
$f_s$ фиксирована техническим заданием, повторяют только внутренний
перебор $r_I$, $L$, $N$ и $g$ при $f_s=f_{\text{ном}}$.

Когда частоту выбирают до изготовления преобразователя, формируют
дискретное множество $\mathcal F_s$ и для каждого его элемента заново
выполняют весь внутренний расчёт:
\begin{equation}
 P_{\Sigma,\text{opt}}(f_s)=
 \min_{\mathbf x,\;r_I}
 P_{\Sigma}(\mathbf x,r_I,f_s),
 \qquad f_s\in\mathcal F_s,
 \label{eq:papamanolis-frequency-outer-loop}
\end{equation}
\begin{whereblock}
\whereitem{$P_{\Sigma,\text{opt}}(f_s)$}{минимальные допустимые потери
дросселя для заданной частоты, Вт;}
\whereitem{$\mathcal F_s$}{множество разрешённых кандидатных частот, Гц.}
\end{whereblock}

Минимизацию в формуле~\eqref{eq:papamanolis-frequency-outer-loop}
выполняют только среди вариантов, удовлетворяющих всем обязательным
ограничениям. Каждая точка кривой
$P_{\Sigma,\text{opt}}(f_s)$ соответствует отдельной оптимизированной
конструкции дросселя; нельзя считать, что один и тот же дроссель
автоматически имеет параметры всех точек этой кривой.

Если частота изменяется во время работы одного физического
преобразователя, $N$, $g$, магнитопровод и обмотку выбирают один раз, а
все частотные режимы оценивают совместно:
\begin{equation}
 J_f(\mathbf x)=
 \sum_{q=1}^{Q}
 w_q P_{\Sigma}
 \bigl(\mathbf x;f_{s,q},U_q,I_q,T_q\bigr),
 \qquad
 \sum_{q=1}^{Q}w_q=1,
 \label{eq:variable-frequency-profile}
\end{equation}
\begin{whereblock}
\whereitem{$J_f$}{средневзвешенные потери одного дросселя в
частотно-переменном режиме, Вт;}
\whereitem{$q$}{номер рабочего режима;}
\whereitem{$Q$}{число учитываемых рабочих режимов;}
\whereitem{$w_q$}{доля времени или вероятность $q$-го режима;}
\whereitem{$f_{s,q}$, $U_q$, $I_q$, $T_q$}{частота коммутации,
характерное напряжение, ток и температура $q$-го режима, Гц, В, А и К.}
\end{whereblock}

Для выбора частоты всего преобразователя минимизируют не только потери
дросселя, но и системную целевую функцию
\begin{equation}
 J_{\text{сист}}(f_s)=
 \sum_{q=1}^{Q}w_q
 \left[
 P_{\Sigma,q}(f_s)+P_{\text{ост},q}(f_s)
 \right],
 \label{eq:system-frequency-objective}
\end{equation}
\begin{whereblock}
\whereitem{$J_{\text{сист}}$}{средневзвешенные потери преобразователя,
учитываемые при выборе частоты, Вт;}
\whereitem{$P_{\Sigma,q}$}{потери дросселя в $q$-м режиме, Вт;}
\whereitem{$P_{\text{ост},q}$}{потери остальных элементов
преобразователя, включая коммутационные потери полупроводников, Вт.}
\end{whereblock}

Статья минимизирует прежде всего потери дросселя. Поэтому найденный в ней
рост эффективности дросселя при повышении частоты нельзя непосредственно
переносить на преобразователь: коммутационные потери полупроводников и
требования электромагнитной совместимости обычно усиливаются с ростом
$f_s$. Гармоники $hf_s$ не являются кандидатными частотами внешнего
перебора; они учитываются внутри каждой рабочей точки при расчёте
$P_{\text{Cu}}$, $P_{\text{Fe}}$ и помех.

\subsubsection{Интерпретация результата частотного перебора}

В сокращённой модели статьи частотная зависимость минимальных потерь
приближённо имеет вид
\begin{equation}
 P_{0,\text{opt}}(f_s)\propto
 f_s^{\,2(\alpha-\beta)/(2+\beta)},
 \label{eq:papamanolis-frequency-scaling}
\end{equation}
\begin{whereblock}
\whereitem{$P_{0,\text{opt}}$}{минимальная сумма потерь в обмотке и
магнитопроводе при заданной частоте, Вт;}
\whereitem{$\alpha$}{локальный показатель степени зависимости потерь
магнитопровода от частоты;}
\whereitem{$\beta$}{локальный показатель степени зависимости потерь
магнитопровода от амплитуды магнитной индукции.}
\end{whereblock}

Формула~\eqref{eq:papamanolis-frequency-scaling} задаёт три области:
\begin{enumerate}
\item при $\alpha<\beta$ минимальные потери дросселя уменьшаются с ростом
частоты, поскольку уменьшаются вольт-секундная площадь и переменная
магнитная индукция;
\item при $\alpha\approx\beta$ минимальные потери слабо зависят от
частоты, поэтому возникает широкая плоская область близких решений;
\item при $\alpha>\beta$ дальнейший рост частоты увеличивает потери
магнитопровода и перестаёт быть выгодным даже на уровне дросселя.
\end{enumerate}

Показатели $\alpha$ и $\beta$ не являются постоянными для всей карты
материала: их локальные значения зависят от $f_s$,
$\widehat B_{\text{AC}}$, $B_{\text{DC}}$ и температуры. Поэтому
частоту пересечения $\alpha\approx\beta$ определяют по данным конкретного
материала, а не переносят из статьи.

В подробной модели статьи при $f_s\gtrsim200\,\text{кГц}$ оптимальная
индуктивность изменяется слабо. Из
формулы~\eqref{eq:papamanolis-buck-l} при почти постоянном $L$ следует
\begin{equation}
 r_{I,\text{opt}}(f_s)\propto\frac{1}{f_s},
 \label{eq:papamanolis-ripple-frequency-trend}
\end{equation}
где $r_{I,\text{opt}}$ --- относительный размах пульсации тока,
соответствующий минимуму потерь при заданной частоте. Следовательно,
уменьшение оптимальной пульсации при росте $f_s$ не означает, что авторы
каждый раз увеличивают $L$: в плоской высокочастотной области они получают
примерно одну и ту же индуктивность при разных сочетаниях $f_s$ и $r_I$.

\begin{table}[H]
\caption{Контрольные результаты статьи Papamanolis}
\label{tab:papamanolis-check-values}
\centering
\small
\begin{tabularx}{\textwidth}{
>{\raggedright\arraybackslash}p{39mm}
>{\raggedleft\arraybackslash}p{19mm}
>{\raggedleft\arraybackslash}p{19mm}
>{\raggedleft\arraybackslash}p{18mm}
>{\raggedleft\arraybackslash}p{22mm}
>{\raggedright\arraybackslash}X}
\toprule
Расчётная точка & $f_s$, кГц & $r_I$, \% & $N$ & $L$, мкГн &
$P_0$, Вт \\
\midrule
Низкочастотная точка OP1 & 80 & 110 & 22 & 114 & 4,32 \\
Глобальный минимум OP2 & 375 & 18 & 18 & 148 & 2,37; в подробной
сеточной модели около 2,4 \\
\bottomrule
\end{tabularx}
\end{table}

Значения таблицы~\ref{tab:papamanolis-check-values} относятся только к
понижающему преобразователю мощностью $2\,\text{кВт}$, магнитопроводу
E55/28/21 из N87, литцендрату с диаметром жилы $100\,\text{мкм}$,
естественной конвекции и условиям статьи. Их используют для проверки
воспроизведения авторского алгоритма, но не как рекомендуемые параметры
другого дросселя. В диапазоне $40\ldots375\,\text{кГц}$ ускоренное
руководство статьи отличалось от подробной модели не более чем на
$17{,}8\,\%$; в экспериментальной проверке среднее абсолютное расхождение
расчётных и измеренных суммарных потерь составило $4{,}85\,\%$, а
максимальное --- $12{,}7\,\%$~\cite{papamanolis}.

\subsection{Построение множества Парето и выбор одного варианта}

Вариант $\mathbf x_a$ доминирует вариант $\mathbf x_b$, если он не хуже по
каждой компоненте $\mathbf F$ и строго лучше хотя бы по одной. Множество
недоминируемых вариантов называют множеством Парето.

\method Разработчик выполняет оптимизацию в следующем порядке:
\begin{enumerate}
\item формирует таблицу всех физически реализуемых сочетаний материала,
типоразмера, $N$, зазоров и проводников;
\item рассчитывает для каждого сочетания все компоненты $\mathbf F$ и
ограничения подраздела~\ref{subsec:mandatory-constraints};
\item исключает варианты, нарушающие хотя бы одно обязательное ограничение;
\item среди оставшихся удаляет доминируемые варианты и строит таблицу
множества Парето;
\item находит минимальные средневзвешенные потери
$\overline P_{\Sigma,\min}$;
\item для задачи «максимум КПД и минимум помех» оставляет варианты
$\overline P_{\Sigma}\leq1{,}02\overline P_{\Sigma,\min}$ и среди них
выбирает вариант с наибольшим
$\min(M_{\text{ДМ}},M_{\text{СФ}})$;
\item если несколько вариантов равны по запасу помех в пределах
$1\,\text{дБ}$, выбирает меньшие $P_{\Sigma,\max}$ и
 $H_{\text{конт,max}}$, затем меньшие объём и стоимость.
\end{enumerate}

\begin{figure}[htbp]
\centering
\includegraphics[
  page=6,
  trim=32bp 158bp 324bp 54bp,
  clip,
  height=0.62\textheight,
  keepaspectratio
]{source_papamanolis_optimization.pdf}
\caption{Линии равных суммарных потерь и положение оптимальной амплитуды
пульсации тока (по Papamanolis с соавторами%
~\cite[рисунки~3, 4]{papamanolis})}
\label{fig:loss-optimum-contours}
\end{figure}

На рисунке~\ref{fig:loss-optimum-contours} минимум расположен внутри
допустимой области, а не обязательно при минимальной пульсации. Уменьшение
пульсации требует большей индуктивности и числа витков, из-за чего могут
вырасти сопротивление, собственная ёмкость и потери меди.

\begin{figure}[htbp]
\centering
\includegraphics[
  page=24,
  trim=160bp 245bp 120bp 80bp,
  clip,
  width=0.82\textwidth
]{source_pichon_optimization.pdf}
\caption{Изменение минимумов потерь при переборе частоты коммутации и
пульсации тока (по Pichon с соавторами%
~\cite[рисунки~22, 23]{pichon})}
\label{fig:loss-optimum-frequency-ripple}
\end{figure}

Графики на рисунке~\ref{fig:loss-optimum-frequency-ripple} используют как
пример чтения результата перебора: для каждого сочетания частоты и
пульсации сначала находят реализуемую конструкцию, затем сравнивают её
потери. Численные координаты минимумов из статьи не являются рекомендацией
для другого преобразователя.

\recommend Не сводят все критерии к одной взвешенной сумме до
построения множества Парето. Сначала исключают недопустимые варианты,
затем выбирают компромисс между потерями, температурой, габаритом и
помехами. Порог $2\,\%$ является правилом настоящей методики для сохранения
почти максимального КПД; при ином приоритете его задают в техническом
задании. Многокритериальный характер выбора и разнообразие близких решений
показаны в \cite{papamanolis,pichon}.

\subsection{Обязательные ограничения}
\label{subsec:mandatory-constraints}

Проект принимают при одновременном выполнении условий:
\begin{align}
 L_{\text{диф,min}}&\geq L_{\text{треб,min}},
 \label{eq:constraint-l}\\
 B_{\max,\text{FEM,нх}}+u_{B,\text{числ}}
 &\leq k_BB_S(T_{\max}),
 \label{eq:acceptance-b-fem}\\
 \vartheta_{\text{обм,max}}&\leq\vartheta_{\text{обм,доп}},
 \label{eq:constraint-tw}\\
 \vartheta_{\text{серд,max}}&\leq\vartheta_{\text{серд,доп}},
 \label{eq:constraint-tc}\\
 k_{\text{Cu}}&\leq k_{\text{Cu,доп}},
 \label{eq:constraint-kw}\\
 k_{\text{зан}}&\leq k_{\text{зан,доп}},
 \label{eq:constraint-kocc}\\
 f_{\text{СР}}&\geq k_{\text{СР}}f_{\text{треб,max}}.
 \label{eq:constraint-srf}
\end{align}
Здесь $L_{\text{диф,min}}$ --- минимальная дифференциальная
индуктивность во всех режимах, Гн;
$L_{\text{треб,min}}$ --- минимально требуемая индуктивность, Гн;
$B_{\max,\text{FEM,нх}}$ --- наибольший локальный максимум магнитной
индукции по всем наихудшим сочетаниям, рассчитанный методом конечных
элементов, Тл;
$u_{B,\text{числ}}$ --- верхняя оценка численной и модельной
неопределённости по формуле~\eqref{eq:b-numerical-uncertainty}, Тл;
$k_B$ --- коэффициент использования паспортной индукции насыщения;
$B_S(T_{\max})$ --- паспортная индукция насыщения при наибольшей расчётной
температуре, Тл;
$\vartheta_{\text{обм,max}}$, $\vartheta_{\text{серд,max}}$ ---
максимальные температуры обмотки и сердечника, $^{\circ}\text{С}$;
$\vartheta_{\text{обм,доп}}$, $\vartheta_{\text{серд,доп}}$ ---
допустимые температуры, $^{\circ}\text{С}$;
$k_{\text{Cu,доп}}$ --- допустимый коэффициент заполнения окна медью;
$k_{\text{зан,доп}}$ --- допустимый коэффициент геометрической занятости
окна.

Дополнительно должны выполняться требования технического задания к
изоляции, электрической прочности, путям утечки и зазорам, механической
прочности, шуму, массе и электромагнитной совместимости. Во всей
нормируемой полосе должны выполняться
$M_{\text{ДМ}}\geq M_{\text{ЭМС}}$,
$M_{\text{СФ}}\geq M_{\text{ЭМС}}$ и ограничения на потери в окружающих
металлических деталях и наведённое напряжение по
подразделу о поле рассеяния.

\section{СОСТАВ ИТОГОВОЙ РАСЧЁТНОЙ ДОКУМЕНТАЦИИ}

Итоговый расчёт должен содержать:
\begin{enumerate}
 \item таблицу исходных данных, режимов и допусков;
 \item указание источника каждой паспортной величины и коэффициента;
 \item формы $u_L(t)$, $i(t)$ и их гармонический состав;
 \item расчёт числа витков, зазоров, индукции и энергии;
 \item чертёж укладки обмотки и расчёт заполнения окна;
 \item раздельные потери по компонентам и гармоникам;
 \item результаты FEM с исследованием сеточной сходимости;
 \item температурное поле при наихудшем режиме;
 \item частотную зависимость $Z(f)$ и частоту первого резонанса;
 \item оценку синфазного тока и внешнего магнитного поля;
 \item таблицу запасов по каждому обязательному ограничению;
 \item программу и протокол испытания опытного образца.
\end{enumerate}

Расчёт считается завершённым после сопоставления модели с измерениями
$L(I,T)$, $R_{\text{AC}}(f,T)$, $Z(f)$, температуры, суммарных потерь и
помех в составе преобразователя.

\section{ЗАКЛЮЧЕНИЕ}

Методика связывает электрический режим, магнитный расчёт, конструкцию
обмотки, потери, паразитные параметры, численную модель и испытание в одну
замкнутую процедуру. Результатом расчёта является не только число витков и
зазор, а полностью определённая конструкция с установленными источниками
исходных данных, допусками и измеримыми критериями приёмки.

Максимальный коэффициент полезного действия получают сравнением полных
потерь по рабочему профилю, а минимальные помехи --- сравнением синфазных и
дифференциальных спектров с применимой нормой. Увеличение сердечника,
уменьшение плотности тока или повышение частоты собственного резонанса по
отдельности не гарантируют оптимума. При отсутствии паспортных данных
величину обозначают как неопределённую и включают её определение в программу
моделирования или испытаний.

Структура, нумерация, запись формул, таблиц и списка источников приняты с
учётом требований ГОСТ~Р~2.105--2019 и ГОСТ~7.32--2017
\cite{gost2105,gost732}.

\begin{thebibliography}{99}

\bibitem{obrusnik}
Обрусник~В.~П.
\textit{Магнитные элементы электронных устройств: руководство к
организации самостоятельной работы студентов}.
Томск: Томский государственный университет систем управления и
радиоэлектроники, 2019. 61~с.

\bibitem{legostaev}
Легостаев~Н.~С.
\textit{Магнитные элементы электронных устройств: учебное пособие}.
Томск: Эль Контент, 2014. 186~с.

\bibitem{gost2105}
ГОСТ~Р~2.105--2019. Единая система конструкторской документации.
Общие требования к текстовым документам. Введён 2021-02-01.
С изменениями №~1, №~2 и опубликованными поправками.

\bibitem{gost732}
ГОСТ~7.32--2017. Система стандартов по информации, библиотечному и
издательскому делу. Отчёт о научно-исследовательской работе. Структура и
правила оформления. Введён 2018-07-01.

\bibitem{codata}
Tiesinga~E., Mohr~P.~J., Newell~D.~B., Taylor~B.~N.
CODATA recommended values of the fundamental physical constants: 2022.
\textit{Journal of Physical and Chemical Reference Data}. 2024. Vol.~53.
No.~3. Art.~033105.
DOI: \url{https://doi.org/10.1063/5.0216647}.

\bibitem{nist_copper}
Paulter~N.~G.
\textit{Users' Guide for Hand-Held and Walk-Through Metal Detectors}.
NIJ Guide 600--00. Washington, DC: U.S. Department of Justice,
National Institute of Justice, 2001.
\url{https://tsapps.nist.gov/publication/get_pdf.cfm?pub_id=30185}.

\bibitem{iec60028}
IEC~60028:1925.
\textit{International standard of resistance for copper}. Edition~2.0.
Geneva: International Electrotechnical Commission, 1925.
\url{https://webstore.iec.ch/en/publication/98}.

\bibitem{tdkcore}
TDK Electronics.
\textit{Ferrites and accessories. E 120/71/40. Core B66614}.
October 2024.

\bibitem{tdkn88}
TDK Electronics.
\textit{SIFERRIT materials N88. Material properties}.
2023.

\bibitem{liao}
Liao~H., Chen~J.-F.
Design process of high-frequency inductor with multiple air-gaps in the
dimensional limitation.
\textit{The Journal of Engineering}. 2022. P.~16--33.
DOI: \url{https://doi.org/10.1049/tje2.12087}.

\bibitem{papamanolis}
Papamanolis~P., Guillod~T., Krismer~F., Kolar~J.~W.
Minimum loss operation and optimal design of high-frequency inductors for
defined core and litz wire.
\textit{IEEE Open Journal of Power Electronics}. 2020. Vol.~1.
P.~469--487.
DOI: \url{https://doi.org/10.1109/OJPEL.2020.3027452}.

\bibitem{pichon}
Pichon~H., Lembeye~Y., Crebier~J.-C.
Accurate efficiency and power densities optimization of output inductor of
buck derived converters.
\textit{Applied Sciences}. 2022. Vol.~12. No.~18. Art.~9330.
DOI: \url{https://doi.org/10.3390/app12189330}.

\bibitem{ewald_litz}
Ewald~T., Biela~J.
Analytical winding loss and inductance models for gapped inductors with
litz or solid wires.
\textit{IEEE Transactions on Power Electronics}. 2022. Vol.~37. No.~12.
P.~15127--15139.
DOI: \url{https://doi.org/10.1109/TPEL.2022.3187155}.

\bibitem{meng}
Meng~Q., Biela~J.
Analytical models for HF-losses of litz wire in inductors with arbitrary
winding and gap arrangements.
\textit{IEEE Open Journal of Power Electronics}. 2025. Vol.~6.
P.~1534--1546.
DOI: \url{https://doi.org/10.1109/OJPEL.2025.3604564}.

\bibitem{muehlethaler}
Mühlethaler~J., Biela~J., Kolar~J.~W.
Improved core-loss calculation for magnetic components employed in power
electronic systems.
\textit{Proceedings of APEC}. 2011. P.~1729--1736.
DOI: \url{https://doi.org/10.1109/APEC.2011.5744829}.

\bibitem{sanusi}
Sanusi~B.~N., Zambach~M., Frandsen~C., Beleggia~M.,
J\o{}rgensen~A., Ouyang~Z.
Investigation and modeling of DC bias impact on core losses at high
frequency.
\textit{IEEE Transactions on Power Electronics}. 2023. Vol.~38. No.~6.
P.~7444--7458.
DOI: \url{https://doi.org/10.1109/TPEL.2023.3249106}.

\bibitem{zhao}
Zhao~H., Huang~Z., Dalal~D.~N., J\o{}rgensen~J.~K.,
J\o{}rgensen~A.~B., Wang~X., Munk-Nielsen~S.
Parasitic capacitance modeling of copper-foiled medium-voltage filter
inductors considering fringe electrical field.
\textit{IEEE Transactions on Power Electronics}. 2021. Vol.~36. No.~7.
P.~8181--8192.
DOI: \url{https://doi.org/10.1109/TPEL.2020.3048226}.

\bibitem{jia}
Jia~N., Xue~L., Cui~H.
Mitigating EMI noise in propagation paths: Review of parasitic and coupling
effects in power electronic packages, filters, and systems.
\textit{IEEE Open Journal of Power Electronics}. 2024. Vol.~5.
P.~352--368.
DOI: \url{https://doi.org/10.1109/OJPEL.2024.3357832}.

\bibitem{sullivan}
Sullivan~C.~R., Zhang~R.~Y.
Simplified design method for litz wire.
\textit{Proceedings of APEC}. 2014. P.~2667--2674.
DOI: \url{https://doi.org/10.1109/APEC.2014.6803681}.

\bibitem{schaefer}
Schäfer~M., Bortis~D., Kolar~J.~W.
Optimal design of highly efficient and highly compact PCB winding
inductors.
\textit{2018 IEEE 19th Workshop on Control and Modeling for Power
Electronics (COMPEL)}. 2018. P.~1--8.
DOI: \url{https://doi.org/10.1109/COMPEL.2018.8460166}.

\bibitem{ewaldfoil}
Ewald~T., Biela~J.
Frequency-dependent inductance and winding loss model for gapped foil
inductors.
\textit{IEEE Transactions on Power Electronics}. 2022. Vol.~37. No.~10.
P.~12370--12379.
DOI: \url{https://doi.org/10.1109/TPEL.2022.3169620}.

\bibitem{heldwein}
Heldwein~M.~L., Kolar~J.~W.
Winding capacitance cancellation for three-phase EMC input filters.
\textit{IEEE Transactions on Power Electronics}. 2008. Vol.~23. No.~4.
P.~2062--2074.
DOI: \url{https://doi.org/10.1109/TPEL.2008.924820}.

\end{thebibliography}

\end{document}
